\begin{comment}
    Euler-Maruyama, Milstein, Implicit Euler, Runge-Kutta, SDE
    https://web.mit.edu/10.001/Web/Course_Notes/Differential_Equations_Notes/node5.html
\end{comment}

\subsection{Stochastic Taylor Expansion}
\noindent Deterministic Taylor expansions are well known. To begin, we shall consider the solution $X_t$ of a $1$-dimensional ordinary differential equation

\begin{equation}\label{eq:DifferentialEquation}
    \frac{d}{d t} X_t = a(X_t),
\end{equation}

\noindent with initial value $X_{t_0}$, for $t \in [t_0, T]$ where $0 \leq t_0 < T$. We can write this in the equivalent integral equation form

\begin{equation}\label{eq:DeterministicIntegralEquation}
    X_t = X_{t_0} + \int_{t_0}^t a(X_s) \, ds.
\end{equation}

\noindent To justify the following constructions, we shall require the function $a$ to satisfy appropriate properties, for example, to be sufficiently smooth and to have a linear growth bound. Let $f : \mathbb{R} \to \mathbb{R}$ be a continuously differentiable function. Then, by the chain rule, we have

\begin{equation}\label{eq:ChainRule}
    \frac{d}{d t} f(X_t) = a(X_t) \frac{\partial}{\partial x} f(X_t),
\end{equation}

\noindent which, using the operator

\begin{equation}\label{eq:OperatorDefinition}
    L = a \frac{\partial}{\partial x},
\end{equation}

\noindent we can express as the integral relation

\begin{equation}\label{eq:IntegralRelation}
    f(X_t) = f(X_{t_0}) + \int_{t_0}^t L f(X_s) \, ds
\end{equation}

\noindent for all $t \in [t_0, T]$. When $f(x) \equiv x$ we have $Lf = a$, $L^2 f = La, \dots$ and \eqref{eq:IntegralRelation} reduces to

\begin{equation}\label{eq:ReducedIntegralRelation}
    X_t = X_{t_0} + \int_{t_0}^t a(X_s) \, ds,
\end{equation}

\noindent that is to equation \eqref{eq:IntegralEquation}.

If we now apply the relation~\eqref{eq:IntegralRelation} to the function $f = a$ in the integral in~\eqref{eq:ReducedIntegralRelation}, we obtain

\begin{equation}\label{eq:FirstTaylorExpansion}
    \begin{split}
        X_t &= X_{t_0} + \int_{t_0}^t \left( a(X_{t_0}) + \int_{t_0}^s La(X_z) \, dz \right) \, ds \\
        &= X_{t_0} + a(X_{t_0}) \int_{t_0}^t \, ds + \int_{t_0}^t \int_{t_0}^s La(X_z) \, dz \, ds,
    \end{split}
\end{equation}

\noindent which is the simplest nontrivial Taylor expansion for $X_t$. We can apply \eqref{eq:IntegralRelation} again to the function $f = La$ in the double integral to derive

\begin{equation}\label{eq:SecondTaylorExpansion}
    X_t = X_{t_0} + a(X_{t_0}) \int_{t_0}^t \, ds + La(X_{t_0}) \int_{t_0}^t \int_{t_0}^s \, dz \, ds + R_3
\end{equation}

\noindent with remainder

\begin{equation}\label{eq:RemainderTerm}
    R_3 = \int_{t_0}^t \int_{t_0}^s \int_{t_0}^z L^2 a(X_u) \, du \, dz \, ds,
\end{equation}

\noindent for $t \in [t_0, T]$. For a general $r+1$ times continuously differentiable function $f : \mathbb{R} \to \mathbb{R}$ this method gives the classical Taylor formula in integral form:

\begin{equation}\label{eq:GeneralTaylorFormula}
    \begin{split}
        f(X_t) &= f(X_{t_0}) + \sum_{l=1}^r \frac{(t-t_0)^l}{l!} L^l f(X_{t_0}) \\
        &\quad + \int_{t_0}^t \dots \int_{t_0}^{s_2} L^{r+1} f(X_{s_1}) \, ds_1 \dots ds_{r+1}
    \end{split}
\end{equation}

\noindent for $t \in [t_0, T]$ and $r = 1, 2, 3, \dots$.

The Taylor formula \eqref{eq:GeneralTaylorFormula} has proven to be a very useful tool in both theoretical and practical investigations, particularly in numerical analysis. It allows the approximation of a sufficiently smooth function in a neighbourhood of a given point to any desired order of accuracy. This expansion depends on the values of the function and some of its higher derivatives at the expansion point, weighted by corresponding multiple time integrals. In addition, there is a remainder term which contains the next multiple time integral, but now with a time-dependent integrand.

We shall indicate it here for the solution $X_t$ of the 1-dimensional Ito stochastic differential equation in integral form

\begin{equation} \label{eq:ItoIntegralForm}
    X_t = X_{t_0} + \int_{t_0}^t a(X_s) \, ds + \int_{t_0}^t b(X_s) \, dW_s
\end{equation}

\noindent for $t \in [t_0, T]$, where the second integral in \eqref{eq:ItoIntegralForm} is an Ito stochastic integral and the coefficients $a$ and $b$ are sufficiently smooth real-valued functions satisfying a linear growth bound. Then, for any twice continuously differentiable function $f : \mathbb{R} \to \mathbb{R}$ the Ito formula \cref{eq:ItoFormulaApplication} gives

\begin{equation} \label{eq:ItoFormulaApplication}
\begin{split}
    f(X_t) &= f(X_{t_0}) \\
    &\quad + \int_{t_0}^t \left( a(X_s)\frac{\partial}{\partial x}f(X_s) + \frac{1}{2}b^2(X_s)\frac{\partial^2}{\partial x^2}f(X_s) \right) \, ds \\
    &\quad + \int_{t_0}^t b(X_s)\frac{\partial}{\partial x}f(X_s) \, dW_s \\
    &= f(X_{t_0}) + \int_{t_0}^t L^0 f(X_s) \, ds + \int_{t_0}^t L^1 f(X_s) \, dW_s,
\end{split}
\end{equation}

\noindent for $t \in [t_0, T]$. Here we have introduced the operators

\begin{equation} \label{eq:OperatorL0}
    L^0 = a \frac{\partial}{\partial x} + \frac{1}{2}b^2 \frac{\partial^2}{\partial x^2}
\end{equation}

\noindent and

\begin{equation} \label{eq:OperatorL1}
    L^1 = b \frac{\partial}{\partial x}.
\end{equation}

\noindent Obviously, for $f(x) \equiv x$ we have $L^0 f = a$ and $L^1 f = b$, in which case \eqref{eq:ItoFormulaApplication} reduces to the original Ito equation for $X_t$, that is to

\begin{equation} \label{eq:OriginalItoEquation}
    X_t = X_{t_0} + \int_{t_0}^t a(X_s) \, ds + \int_{t_0}^t b(X_s) \, dW_s.
\end{equation}

\noindent In analogy with the above deterministic expansions, if we apply the Ito formula \eqref{eq:ItoFormulaApplication} to the functions $f = a$ and $f = b$ in \eqref{eq:OriginalItoEquation} we obtain

\begin{equation} \label{eq:ItoCoefficientExpansion}
\begin{split}
    X_t &= X_{t_0} \\
    &\quad + \int_{t_0}^t \left( a(X_{t_0}) + \int_{t_0}^s L^0 a(X_z) \, dz + \int_{t_0}^s L^1 a(X_z) \, dW_z \right) \, ds \\
    &\quad + \int_{t_0}^t \left( b(X_{t_0}) + \int_{t_0}^s L^0 b(X_z) \, dz + \int_{t_0}^s L^1 b(X_z) \, dW_z \right) \, dW_s \\
    &= X_{t_0} + a(X_{t_0}) \int_{t_0}^t ds + b(X_{t_0}) \int_{t_0}^t dW_s + R
\end{split}
\end{equation}

\noindent with remainder

\begin{equation} \label{eq:FirstRemainder}
\begin{split}
    R &= \int_{t_0}^t \int_{t_0}^s L^0 a(X_z) \, dz \, ds + \int_{t_0}^t \int_{t_0}^s L^1 a(X_z) \, dW_z \, ds \\
    &\quad + \int_{t_0}^t \int_{t_0}^s L^0 b(X_z) \, dz \, dW_s + \int_{t_0}^t \int_{t_0}^s L^1 b(X_z) \, dW_z \, dW_s.
\end{split}
\end{equation}

\noindent This is the simplest nontrivial Ito-Taylor expansion. We can continue it, for instance, by applying the Ito formula \eqref{eq:ItoFormulaApplication} to $f = L^1 b$ in \eqref{eq:ItoCoefficientExpansion}, in which case we get

\begin{equation} \label{eq:HigherOrderExpansion}
\begin{split}
    X_t &= X_{t_0} + a(X_{t_0}) \int_{t_0}^t ds + b(X_{t_0}) \int_{t_0}^t dW_s \\
    &\quad + L^1 b(X_{t_0}) \int_{t_0}^t \int_{t_0}^s dW_z \, dW_s + \bar{R}
\end{split}
\end{equation}

\noindent with remainder

\begin{equation} \label{eq:SecondRemainder}
\begin{split}
    \bar{R} &= \int_{t_0}^t \int_{t_0}^s L^0 a(X_z) \, dz \, ds + \int_{t_0}^t \int_{t_0}^s L^1 a(X_z) \, dW_z \, ds \\
    &\quad + \int_{t_0}^t \int_{t_0}^s L^0 b(X_z) \, dz \, dW_s + \int_{t_0}^t \int_{t_0}^s \int_{t_0}^z L^0 L^1 b(X_u) \, du \, dW_z \, dW_s \\
    &\quad + \int_{t_0}^t \int_{t_0}^s \int_{t_0}^z L^1 L^1 b(X_u) \, dW_u \, dW_z \, dW_s.
\end{split}
\end{equation}

\noindent We shall now consider another representation for the increments of a function of an Ito process $X$, which we shall call the \textit{Stratonovich-Taylor expansion}. For this we start with the $1$-dimensional Stratonovich stochastic differential equation in integral form

\begin{equation} \label{eq:StratonovichIntegralForm}
    X_t = X_{t_0} + \int_{t_0}^t \underline{a}(X_s) \, ds + \int_{t_0}^t b(X_s) \circ dW_s
\end{equation}

\noindent for $t \in [t_0, T]$, where the second integral in \eqref{eq:StratonovichIntegralForm} is a Stratonovich stochastic integral and the coefficients $\underline{a}$ and $b$ are sufficiently smooth real valued functions satisfying a linear growth bound. Consider that

\begin{equation*}
    \underline{a} = a - \frac{1}{2} bb'
\end{equation*}

\noindent the \cref{eq:ItoIntegralForm} and the Stratonovich equation \eqref{eq:StratonovichIntegralForm} are equivalent and have the same solutions. The solution of a Stratonovich SDE transforms according to the deterministic chain rule, so for any twice continuously differentiable function $f : \mathbb{R} \to \mathbb{R}$ we have

\begin{equation} \label{eq:StratonovichItoFormula}
\begin{split}
    f(X_t) &= f(X_{t_0}) + \int_{t_0}^t \underline{a}(X_s)\frac{\partial}{\partial x}f(X_s) \, ds+ \int_{t_0}^t b(X_s)\frac{\partial}{\partial x}f(X_s) \circ dW_s, \\
    &= f(X_{t_0}) + \int_{t_0}^t \underline{L}^0 f(X_s) \, ds + \int_{t_0}^t \underline{L}^1 f(X_s) \circ dW_s,
\end{split}
\end{equation}

\noindent for $t \in [t_0, T]$, with the operators

\begin{equation} \label{eq:StratonovichOperatorL0}
    \underline{L}^0 = \underline{a} \frac{\partial}{\partial x}
\end{equation}

\noindent and

\begin{equation} \label{eq:StratonovichOperatorL1}
    \underline{L}^1 = b \frac{\partial}{\partial x}.
\end{equation}

\noindent Obviously, for $f(x) \equiv x$ we have $\underline{L}^0 f = \underline{a}$ and $\underline{L}^1 f = b$, in which case \eqref{eq:StratonovichItoFormula} reduces to the original Stratonovich equation

\begin{equation} \label{eq:StratonovichIdentity}
    X_t = X_{t_0} + \int_{t_0}^t \underline{a}(X_s) \, ds + \int_{t_0}^t b(X_s) \circ dW_s.
\end{equation}

\noindent Analogously with the Ito case just considered, we can apply \eqref{eq:StratonovichItoFormula} to the integrand functions $f = \underline{a}$ and $f = b$ in \eqref{eq:StratonovichOperatorL0}. This gives

\begin{equation} \label{eq:StratonovichCoefficientExpansion}
\begin{split}
    X_t &= X_{t_0} \\
    &\quad + \int_{t_0}^t \biggl( \underline{a}(X_{t_0}) + \int_{t_0}^s \underline{L}^0 \underline{a}(X_z) \, dz + \int_{t_0}^s \underline{L}^1 \underline{a}(X_z) \circ dW_z \biggr) \, ds \\
    &\quad + \int_{t_0}^t \biggl( b(X_{t_0}) + \int_{t_0}^s \underline{L}^0 b(X_z) \, dz + \int_{t_0}^s \underline{L}^1 b(X_z) \circ dW_z \biggr) \circ dW_s \\
    &= X_{t_0} + \underline{a}(X_{t_0}) \int_{t_0}^t ds + b(X_{t_0}) \int_{t_0}^t \circ dW_s + R
\end{split}
\end{equation}

\noindent with remainder

\begin{equation} \label{eq:StratonovichRemainderFirst}
\begin{split}
    R &= \int_{t_0}^t \int_{t_0}^s \underline{L}^0 \underline{a}(X_z) \, dz \, ds + \int_{t_0}^t \int_{t_0}^s \underline{L}^1 \underline{a}(X_z) \circ dW_z \, ds \\
    &\quad + \int_{t_0}^t \int_{t_0}^s \underline{L}^0 b(X_z) \, dz \circ dW_s + \int_{t_0}^t \int_{t_0}^s \underline{L}^1 b(X_z) \circ dW_z \circ dW_s.
\end{split}
\end{equation}

\noindent This is the simplest nontrivial Stratonovich-Taylor expansion of $f(X_t)$. We can continue expanding, for instance, by applying \eqref{eq:StratonovichItoFormula} to the integrand $f = \underline{L}^1 b$ in \eqref{eq:StratonovichCoefficientExpansion} to obtain

\begin{equation} \label{eq:StratonovichNextExpansion}
\begin{split}
    X_t &= X_{t_0} + \underline{a}(X_{t_0}) \int_{t_0}^t ds + b(X_{t_0}) \int_{t_0}^t \circ dW_s \\
    &\quad + \underline{L}^1 b(X_{t_0}) \int_{t_0}^t \int_{t_0}^s \circ dW_z \circ dW_s + \bar{R}
\end{split}
\end{equation}

\noindent with remainder

\begin{equation} \label{eq:StratonovichRemainderSecond}
\begin{split}
    \bar{R} &= \int_{t_0}^t \int_{t_0}^s \underline{L}^0 \underline{a}(X_z) \, dz \, ds + \int_{t_0}^t \int_{t_0}^s \underline{L}^1 \underline{a}(X_z) \circ dW_z \, ds \\
    &\quad + \int_{t_0}^t \int_{t_0}^s \underline{L}^0 b(X_z) \, dz \circ dW_s \\
    &\quad + \int_{t_0}^t \int_{t_0}^s \int_{t_0}^z \underline{L}^0 \underline{L}^1 b(X_u) \, du \circ dW_z \circ dW_s \\
    &\quad + \int_{t_0}^t \int_{t_0}^s \int_{t_0}^z \underline{L}^1 \underline{L}^1 b(X_u) \circ dW_u \circ dW_z \circ dW_s.
\end{split}
\end{equation}

\noindent Similar expansions hold for multi-dimensional Ito processes satisfying nonautonomous stochastic differential equations. We shall refer to the nonautonomous $d$-dimensional Ito equation

\begin{equation} \label{eq:MultiDimIto}
    X_t = X_{t_0} + \int_{t_0}^t a(s, X_s) \, ds + \sum_{j=1}^m \int_{t_0}^t b^j(s, X_s) \, dW^j_s
\end{equation}

\noindent for $t \in [t_0, T]$ and the equivalent Stratonovich equation

\begin{equation} \label{eq:MultiDimStratonovich}
    X_t = X_{t_0} + \int_{t_0}^t \underline{a}(s, X_s) \, ds + \sum_{j=1}^m \int_{t_0}^t b^j(s, X_s) \circ dW^j_s
\end{equation}

\noindent for $t \in [t_0, T]$, where

\begin{equation} \label{eq:MultiDimCorrection}
    \underline{a}^i = a^i - \frac{1}{2} \sum_{j=1}^m \sum_{k=1}^d b^{k,j} \frac{\partial b^{i,j}}{\partial x^k}
\end{equation}

\noindent for $i = 1, 2, \dots, d$. In both of these SDEs $W = \{W_t, t \in [0, T]\}$ is a standard $m$-dimensional Wiener process adapted to an increasing family of sub-$\sigma$-algebras $\{\mathcal{A}_t, t \in [0, T]\}$.

\subsection{Strong Convergence Criterion}
In simulation, filtering, and estimation of \Ito{} processes, the sample paths of the approximation must closely track those of the underlying process. The primary criterion for this proximity is the absolute error at the final time instant $T$.
\begin{equation}
    \epsilon(\delta) = \E(|X_T - Y_N|).
    \label{eq:StrongConvergenceError}
\end{equation}
This error is bounded by the root mean square error via the \textit{Lyapunov inequality}:
\begin{equation}
    \epsilon(\delta) = \E(|X_T - Y_N|) \le \sqrt{\E(|X_T - Y_N|^2)}.
    \label{eq:LyapunovBound}
\end{equation}

\begin{defn}{Strong Order of Convergence}{StrongOrderOfConvergence}
An approximating process $Y$ converges in the \textit{strong sense with order} $\gamma \in (0, \infty]$ if there exists a finite constant $K$ and a positive constant $\delta_0$ such that
\begin{equation}
    \E(|X_T - Y_N|) \le K \delta^\gamma
    \label{eq:StrongConvergenceOrder}
\end{equation}
for any time discretization with maximum step size $\delta \in (0, \delta_0)$.
\end{defn}

\begin{xtip}
In the deterministic limit where diffusion vanishes ($b \equiv 0$), the strong convergence criterion reduces to the standard error criterion for ordinary differential equations (ODEs). Stochastic schemes generically exhibit lower convergence orders than their deterministic counterparts because the regularity of the Wiener process dictates that increments $\Delta W_n$ scale as $\delta^{1/2}$ in the root mean square sense rather than $\delta$. This intrinsic degradation results in the Euler approximation achieving a strong order of $\gamma = 0.5$ for SDEs in contrast to the order $1.0$ obtained for ODEs.
\end{xtip}

\subsection{Weak Convergence Criterion}
In many practical scenarios, such as moment estimation or option pricing, a precise pathwise approximation of the \Ito{} process is not required. Instead, the focus is often on the expectation of a function of the process at a final time $T$, such as the first two moments $\E(X_T)$ and $\E((X_T)^2)$, or more generally $\E(g(X_T))$ for some function $g$. In these cases, it suffices to approximate the probability distribution of the random variable $X_T$ rather than the sample paths. This leads to a weaker convergence criterion than the strong sense.

\begin{defn}{Weak Order of Convergence}{WeakOrderOfConvergence}
A time discrete approximation $Y$ converges in the \textit{weak sense with order} $\beta \in (0, \infty]$ if for any polynomial $g$ there exists a finite constant $K$ and a positive constant $\delta_0$ such that
\begin{equation}
    |\E(g(X_T)) - \E(g(Y_N))| \le K \delta^\beta
    \label{eq:WeakConvergenceOrder}
\end{equation}
for any time discretization with maximum step size $\delta \in (0, \delta_0)$.
\end{defn}

\begin{xtip}
In the limit of vanishing diffusion where $b \equiv 0$, the weak convergence criterion with $g(x) \equiv x$ collapses to the standard deterministic convergence criterion for ordinary differential equations. For stochastic systems, the weak order $\beta$ frequently diverges from the strong order $\gamma$, exemplified by the Euler approximation achieving a higher weak order of $\beta = 1.0$ compared to its strong order of $\gamma = 0.5$. This distinction becomes critical when the objective involves estimating moments or probability distributions rather than simulating specific sample paths.
\end{xtip}

\subsection{Euler-Maruyama}
\begin{comment}
    \begin{equation}
    \mathbb{E} [\phi(X)] \simeq \frac{1}{N} \sum_{i=1}^{N} \phi(\hat{X}_i) \notag
    derivative hedging using reinforcement learning
    improve basket trading using bayesian optimisation
    robust portfolio optimisation using SPO
\end{equation}
\end{comment}
To apply the Monte Carlo method in option pricing, we generate a sequence $(Y_1, \dots, Y_N)$ of sample realisations of a random variable, such that the empirical mean
\begin{equation}
    \E [\phi(X)] \approx \frac{1}{N} \sum_{i=1}^{N} \phi(Y_i)
    \label{eq:MonteCarloEstimator}
\end{equation}
converges to the expected value $\E [\phi(X)]$ by the Strong Law of Large Numbers.

Random samples for the solution of a stochastic differential equation defined on the probability space $(\Omega, \scF, \Prob)$ of the form
\begin{equation}
    \di X_t = b(t, X_t) \di t + a(t, X_t) \di W_t,
    \label{eq:SDE}
\end{equation}
where $(W_t)_{t\in\mathbb{R}_+}$ is a standard Brownian motion, can be generated by time discretisation on a grid $\pi_{\delta} = \{0 = t_0 < t_1 < \dots < t_M = T\}$.

More precisely, the Euler-Maruyama discretisation scheme for the stochastic differential equation, \cref{eq:SDE}, is given by the recursive approximation $Y^\delta$:
\begin{align}
    Y_{t_{k+1}}^{\delta} &= Y_{t_{k}}^{\delta} + \int_{t_{k}}^{t_{k+1}} b(t_k, Y^\delta_{t_k}) \di s + \int_{t_{k}}^{t_{k+1}} a(t_k, Y^\delta_{t_k}) \di W_s \notag \\
    &\approx Y_{t_{k}}^{\delta} + b(t_k, Y_{t_{k}}^{\delta}) (t_{k+1} - t_k) + a(t_k, Y_{t_{k}}^{\delta}) (W_{t_{k+1}} - W_{t_k}),
    \label{eq:EulerScheme}
\end{align}
where the increments satisfy $W_{t_{k+1}} - W_{t_k} \sim \Nor (0,t_{k+1} - t_{k})$ for $k = 0,1,\dots,M-1$.

Let $T > 0$ and consider the stochastic differential equation \cref{eq:SDE} on the time interval $[0, T]$. Suppose that the initial data satisfies
\begin{align}
    \E \left(|X_0|^2\right) &< \infty, \label{eq:InitialMoment} \\
    \E \left(|X_0 - Y^\delta_0|^2\right)^{1/2} &\leq K_1 \delta^{1/2}, \label{eq:InitialApproximationError}
\end{align}
and that the coefficients $a: [0, T] \times \R^d \to \R^{d \times m}$ and $b: [0, T] \times \R^d \to \R^d$ satisfy the global Lipschitz condition
\begin{equation}
    |a(t, x) - a(t, y)| + |b(t, x) - b(t, y)| \leq K_2 |x - y| \label{eq:LipschitzConditionSDENumerics}
\end{equation}
and the linear growth condition
\begin{equation}
    |a(t, x)| + |b(t, x)| \leq K_3 (1 + |x|) \label{eq:LinearGrowthCondition}
\end{equation}
for all $t \in [0, T]$ and $x, y \in \R^d$. Furthermore, assume the coefficients satisfy the Hölder continuity condition in time:
\begin{equation}
    |a(s, x) - a(t, x)| + |b(s, x) - b(t, x)| \leq K_4 (1 + |x|) |s - t|^{1/2} \label{eq:HolderContinuity}
\end{equation}
for all $s, t \in [0, T]$ and $x \in \R^d$. The constants $K_1, \dots, K_4$ must not depend on the discretisation step size $\delta$.

\begin{thm}{Strong Convergence of the Euler-Maruyama Scheme}{StrongConvergenceEuler}
    Then, for the Euler-Maruyama approximation $Y^\delta$ defined on the grid $\pi_\delta$, the following strong error estimate holds:
    \begin{equation}
        \E \left(|X_T - Y^\delta(T)|\right) \leq C \delta^{1/2}, \label{eq:StrongConvergenceBound}
    \end{equation}
    where the constant $C$ depends only on $K_1, \dots, K_4$ and $T$, but does not depend on $\delta$.
\end{thm}

We illustrate Euler’s method via the geometric Brownian motion SDE

\begin{Pcode}{eulerMaruyama.py}
    T, N, R = 1.0, 2**8, 4
    dt = T / N
    mu, sigma, X0 = 1.0, 0.5, 1.0

    np.random.seed(4)
    dB = np.sqrt(dt) * np.random.randn(N)

    t = np.linspace(0, T, N + 1)
    B = np.r_[0, np.cumsum(dB)]
    Y = X0 * np.exp((mu - 0.5 * sigma**2) * t + sigma * B)

    X = np.r_[X0, X0 * np.cumprod(1 + mu * dt + sigma * dB)]

    dB_c = dB.reshape(-1, R).sum(axis=1)
    X_c = np.r_[X0, X0 * np.cumprod(1 + mu*dt*R + sigma*dB_c)]
\end{Pcode}

\begin{figure}[H]
    \centering
    %% Creator: Matplotlib, PGF backend
%%
%% To include the figure in your LaTeX document, write
%%   \input{<filename>.pgf}
%%
%% Make sure the required packages are loaded in your preamble
%%   \usepackage{pgf}
%%
%% Also ensure that all the required font packages are loaded; for instance,
%% the lmodern package is sometimes necessary when using math font.
%%   \usepackage{lmodern}
%%
%% Figures using additional raster images can only be included by \input if
%% they are in the same directory as the main LaTeX file. For loading figures
%% from other directories you can use the `import` package
%%   \usepackage{import}
%%
%% and then include the figures with
%%   \import{<path to file>}{<filename>.pgf}
%%
%% Matplotlib used the following preamble
%%   \def\mathdefault#1{#1}
%%   \everymath=\expandafter{\the\everymath\displaystyle}
%%   \IfFileExists{scrextend.sty}{
%%     \usepackage[fontsize=10.000000pt]{scrextend}
%%   }{
%%     \renewcommand{\normalsize}{\fontsize{10.000000}{12.000000}\selectfont}
%%     \normalsize
%%   }
%%   \usepackage{amsmath}
%%   \usepackage{amssymb}
%%   \usepackage{amsfonts}
%%   \usepackage{libertine}
%%   \usepackage[libertine]{newtxmath}
%%   \makeatletter\@ifpackageloaded{underscore}{}{\usepackage[strings]{underscore}}\makeatother
%%
\begingroup%
\makeatletter%
\begin{pgfpicture}%
\pgfpathrectangle{\pgfpointorigin}{\pgfqpoint{6.288009in}{4.038523in}}%
\pgfusepath{use as bounding box, clip}%
\begin{pgfscope}%
\pgfsetbuttcap%
\pgfsetmiterjoin%
\definecolor{currentfill}{rgb}{1.000000,1.000000,1.000000}%
\pgfsetfillcolor{currentfill}%
\pgfsetlinewidth{0.000000pt}%
\definecolor{currentstroke}{rgb}{1.000000,1.000000,1.000000}%
\pgfsetstrokecolor{currentstroke}%
\pgfsetdash{}{0pt}%
\pgfpathmoveto{\pgfqpoint{0.000000in}{0.000000in}}%
\pgfpathlineto{\pgfqpoint{6.288009in}{0.000000in}}%
\pgfpathlineto{\pgfqpoint{6.288009in}{4.038523in}}%
\pgfpathlineto{\pgfqpoint{0.000000in}{4.038523in}}%
\pgfpathlineto{\pgfqpoint{0.000000in}{0.000000in}}%
\pgfpathclose%
\pgfusepath{fill}%
\end{pgfscope}%
\begin{pgfscope}%
\pgfsetbuttcap%
\pgfsetmiterjoin%
\definecolor{currentfill}{rgb}{1.000000,1.000000,1.000000}%
\pgfsetfillcolor{currentfill}%
\pgfsetlinewidth{0.000000pt}%
\definecolor{currentstroke}{rgb}{0.000000,0.000000,0.000000}%
\pgfsetstrokecolor{currentstroke}%
\pgfsetstrokeopacity{0.000000}%
\pgfsetdash{}{0pt}%
\pgfpathmoveto{\pgfqpoint{0.448229in}{0.413680in}}%
\pgfpathlineto{\pgfqpoint{6.155857in}{0.413680in}}%
\pgfpathlineto{\pgfqpoint{6.155857in}{3.988523in}}%
\pgfpathlineto{\pgfqpoint{0.448229in}{3.988523in}}%
\pgfpathlineto{\pgfqpoint{0.448229in}{0.413680in}}%
\pgfpathclose%
\pgfusepath{fill}%
\end{pgfscope}%
\begin{pgfscope}%
\pgfsetbuttcap%
\pgfsetroundjoin%
\definecolor{currentfill}{rgb}{0.000000,0.000000,0.000000}%
\pgfsetfillcolor{currentfill}%
\pgfsetlinewidth{0.501875pt}%
\definecolor{currentstroke}{rgb}{0.000000,0.000000,0.000000}%
\pgfsetstrokecolor{currentstroke}%
\pgfsetdash{}{0pt}%
\pgfsys@defobject{currentmarker}{\pgfqpoint{0.000000in}{0.000000in}}{\pgfqpoint{0.000000in}{0.041667in}}{%
\pgfpathmoveto{\pgfqpoint{0.000000in}{0.000000in}}%
\pgfpathlineto{\pgfqpoint{0.000000in}{0.041667in}}%
\pgfusepath{stroke,fill}%
}%
\begin{pgfscope}%
\pgfsys@transformshift{0.448229in}{0.413680in}%
\pgfsys@useobject{currentmarker}{}%
\end{pgfscope}%
\end{pgfscope}%
\begin{pgfscope}%
\pgfsetbuttcap%
\pgfsetroundjoin%
\definecolor{currentfill}{rgb}{0.000000,0.000000,0.000000}%
\pgfsetfillcolor{currentfill}%
\pgfsetlinewidth{0.501875pt}%
\definecolor{currentstroke}{rgb}{0.000000,0.000000,0.000000}%
\pgfsetstrokecolor{currentstroke}%
\pgfsetdash{}{0pt}%
\pgfsys@defobject{currentmarker}{\pgfqpoint{0.000000in}{-0.041667in}}{\pgfqpoint{0.000000in}{0.000000in}}{%
\pgfpathmoveto{\pgfqpoint{0.000000in}{0.000000in}}%
\pgfpathlineto{\pgfqpoint{0.000000in}{-0.041667in}}%
\pgfusepath{stroke,fill}%
}%
\begin{pgfscope}%
\pgfsys@transformshift{0.448229in}{3.988523in}%
\pgfsys@useobject{currentmarker}{}%
\end{pgfscope}%
\end{pgfscope}%
\begin{pgfscope}%
\definecolor{textcolor}{rgb}{0.000000,0.000000,0.000000}%
\pgfsetstrokecolor{textcolor}%
\pgfsetfillcolor{textcolor}%
\pgftext[x=0.448229in,y=0.365069in,,top]{\color{textcolor}{\rmfamily\fontsize{10.000000}{12.000000}\selectfont\catcode`\^=\active\def^{\ifmmode\sp\else\^{}\fi}\catcode`\%=\active\def%{\%}$\mathdefault{0.0}$}}%
\end{pgfscope}%
\begin{pgfscope}%
\pgfsetbuttcap%
\pgfsetroundjoin%
\definecolor{currentfill}{rgb}{0.000000,0.000000,0.000000}%
\pgfsetfillcolor{currentfill}%
\pgfsetlinewidth{0.501875pt}%
\definecolor{currentstroke}{rgb}{0.000000,0.000000,0.000000}%
\pgfsetstrokecolor{currentstroke}%
\pgfsetdash{}{0pt}%
\pgfsys@defobject{currentmarker}{\pgfqpoint{0.000000in}{0.000000in}}{\pgfqpoint{0.000000in}{0.041667in}}{%
\pgfpathmoveto{\pgfqpoint{0.000000in}{0.000000in}}%
\pgfpathlineto{\pgfqpoint{0.000000in}{0.041667in}}%
\pgfusepath{stroke,fill}%
}%
\begin{pgfscope}%
\pgfsys@transformshift{1.589754in}{0.413680in}%
\pgfsys@useobject{currentmarker}{}%
\end{pgfscope}%
\end{pgfscope}%
\begin{pgfscope}%
\pgfsetbuttcap%
\pgfsetroundjoin%
\definecolor{currentfill}{rgb}{0.000000,0.000000,0.000000}%
\pgfsetfillcolor{currentfill}%
\pgfsetlinewidth{0.501875pt}%
\definecolor{currentstroke}{rgb}{0.000000,0.000000,0.000000}%
\pgfsetstrokecolor{currentstroke}%
\pgfsetdash{}{0pt}%
\pgfsys@defobject{currentmarker}{\pgfqpoint{0.000000in}{-0.041667in}}{\pgfqpoint{0.000000in}{0.000000in}}{%
\pgfpathmoveto{\pgfqpoint{0.000000in}{0.000000in}}%
\pgfpathlineto{\pgfqpoint{0.000000in}{-0.041667in}}%
\pgfusepath{stroke,fill}%
}%
\begin{pgfscope}%
\pgfsys@transformshift{1.589754in}{3.988523in}%
\pgfsys@useobject{currentmarker}{}%
\end{pgfscope}%
\end{pgfscope}%
\begin{pgfscope}%
\definecolor{textcolor}{rgb}{0.000000,0.000000,0.000000}%
\pgfsetstrokecolor{textcolor}%
\pgfsetfillcolor{textcolor}%
\pgftext[x=1.589754in,y=0.365069in,,top]{\color{textcolor}{\rmfamily\fontsize{10.000000}{12.000000}\selectfont\catcode`\^=\active\def^{\ifmmode\sp\else\^{}\fi}\catcode`\%=\active\def%{\%}$\mathdefault{0.2}$}}%
\end{pgfscope}%
\begin{pgfscope}%
\pgfsetbuttcap%
\pgfsetroundjoin%
\definecolor{currentfill}{rgb}{0.000000,0.000000,0.000000}%
\pgfsetfillcolor{currentfill}%
\pgfsetlinewidth{0.501875pt}%
\definecolor{currentstroke}{rgb}{0.000000,0.000000,0.000000}%
\pgfsetstrokecolor{currentstroke}%
\pgfsetdash{}{0pt}%
\pgfsys@defobject{currentmarker}{\pgfqpoint{0.000000in}{0.000000in}}{\pgfqpoint{0.000000in}{0.041667in}}{%
\pgfpathmoveto{\pgfqpoint{0.000000in}{0.000000in}}%
\pgfpathlineto{\pgfqpoint{0.000000in}{0.041667in}}%
\pgfusepath{stroke,fill}%
}%
\begin{pgfscope}%
\pgfsys@transformshift{2.731280in}{0.413680in}%
\pgfsys@useobject{currentmarker}{}%
\end{pgfscope}%
\end{pgfscope}%
\begin{pgfscope}%
\pgfsetbuttcap%
\pgfsetroundjoin%
\definecolor{currentfill}{rgb}{0.000000,0.000000,0.000000}%
\pgfsetfillcolor{currentfill}%
\pgfsetlinewidth{0.501875pt}%
\definecolor{currentstroke}{rgb}{0.000000,0.000000,0.000000}%
\pgfsetstrokecolor{currentstroke}%
\pgfsetdash{}{0pt}%
\pgfsys@defobject{currentmarker}{\pgfqpoint{0.000000in}{-0.041667in}}{\pgfqpoint{0.000000in}{0.000000in}}{%
\pgfpathmoveto{\pgfqpoint{0.000000in}{0.000000in}}%
\pgfpathlineto{\pgfqpoint{0.000000in}{-0.041667in}}%
\pgfusepath{stroke,fill}%
}%
\begin{pgfscope}%
\pgfsys@transformshift{2.731280in}{3.988523in}%
\pgfsys@useobject{currentmarker}{}%
\end{pgfscope}%
\end{pgfscope}%
\begin{pgfscope}%
\definecolor{textcolor}{rgb}{0.000000,0.000000,0.000000}%
\pgfsetstrokecolor{textcolor}%
\pgfsetfillcolor{textcolor}%
\pgftext[x=2.731280in,y=0.365069in,,top]{\color{textcolor}{\rmfamily\fontsize{10.000000}{12.000000}\selectfont\catcode`\^=\active\def^{\ifmmode\sp\else\^{}\fi}\catcode`\%=\active\def%{\%}$\mathdefault{0.4}$}}%
\end{pgfscope}%
\begin{pgfscope}%
\pgfsetbuttcap%
\pgfsetroundjoin%
\definecolor{currentfill}{rgb}{0.000000,0.000000,0.000000}%
\pgfsetfillcolor{currentfill}%
\pgfsetlinewidth{0.501875pt}%
\definecolor{currentstroke}{rgb}{0.000000,0.000000,0.000000}%
\pgfsetstrokecolor{currentstroke}%
\pgfsetdash{}{0pt}%
\pgfsys@defobject{currentmarker}{\pgfqpoint{0.000000in}{0.000000in}}{\pgfqpoint{0.000000in}{0.041667in}}{%
\pgfpathmoveto{\pgfqpoint{0.000000in}{0.000000in}}%
\pgfpathlineto{\pgfqpoint{0.000000in}{0.041667in}}%
\pgfusepath{stroke,fill}%
}%
\begin{pgfscope}%
\pgfsys@transformshift{3.872806in}{0.413680in}%
\pgfsys@useobject{currentmarker}{}%
\end{pgfscope}%
\end{pgfscope}%
\begin{pgfscope}%
\pgfsetbuttcap%
\pgfsetroundjoin%
\definecolor{currentfill}{rgb}{0.000000,0.000000,0.000000}%
\pgfsetfillcolor{currentfill}%
\pgfsetlinewidth{0.501875pt}%
\definecolor{currentstroke}{rgb}{0.000000,0.000000,0.000000}%
\pgfsetstrokecolor{currentstroke}%
\pgfsetdash{}{0pt}%
\pgfsys@defobject{currentmarker}{\pgfqpoint{0.000000in}{-0.041667in}}{\pgfqpoint{0.000000in}{0.000000in}}{%
\pgfpathmoveto{\pgfqpoint{0.000000in}{0.000000in}}%
\pgfpathlineto{\pgfqpoint{0.000000in}{-0.041667in}}%
\pgfusepath{stroke,fill}%
}%
\begin{pgfscope}%
\pgfsys@transformshift{3.872806in}{3.988523in}%
\pgfsys@useobject{currentmarker}{}%
\end{pgfscope}%
\end{pgfscope}%
\begin{pgfscope}%
\definecolor{textcolor}{rgb}{0.000000,0.000000,0.000000}%
\pgfsetstrokecolor{textcolor}%
\pgfsetfillcolor{textcolor}%
\pgftext[x=3.872806in,y=0.365069in,,top]{\color{textcolor}{\rmfamily\fontsize{10.000000}{12.000000}\selectfont\catcode`\^=\active\def^{\ifmmode\sp\else\^{}\fi}\catcode`\%=\active\def%{\%}$\mathdefault{0.6}$}}%
\end{pgfscope}%
\begin{pgfscope}%
\pgfsetbuttcap%
\pgfsetroundjoin%
\definecolor{currentfill}{rgb}{0.000000,0.000000,0.000000}%
\pgfsetfillcolor{currentfill}%
\pgfsetlinewidth{0.501875pt}%
\definecolor{currentstroke}{rgb}{0.000000,0.000000,0.000000}%
\pgfsetstrokecolor{currentstroke}%
\pgfsetdash{}{0pt}%
\pgfsys@defobject{currentmarker}{\pgfqpoint{0.000000in}{0.000000in}}{\pgfqpoint{0.000000in}{0.041667in}}{%
\pgfpathmoveto{\pgfqpoint{0.000000in}{0.000000in}}%
\pgfpathlineto{\pgfqpoint{0.000000in}{0.041667in}}%
\pgfusepath{stroke,fill}%
}%
\begin{pgfscope}%
\pgfsys@transformshift{5.014331in}{0.413680in}%
\pgfsys@useobject{currentmarker}{}%
\end{pgfscope}%
\end{pgfscope}%
\begin{pgfscope}%
\pgfsetbuttcap%
\pgfsetroundjoin%
\definecolor{currentfill}{rgb}{0.000000,0.000000,0.000000}%
\pgfsetfillcolor{currentfill}%
\pgfsetlinewidth{0.501875pt}%
\definecolor{currentstroke}{rgb}{0.000000,0.000000,0.000000}%
\pgfsetstrokecolor{currentstroke}%
\pgfsetdash{}{0pt}%
\pgfsys@defobject{currentmarker}{\pgfqpoint{0.000000in}{-0.041667in}}{\pgfqpoint{0.000000in}{0.000000in}}{%
\pgfpathmoveto{\pgfqpoint{0.000000in}{0.000000in}}%
\pgfpathlineto{\pgfqpoint{0.000000in}{-0.041667in}}%
\pgfusepath{stroke,fill}%
}%
\begin{pgfscope}%
\pgfsys@transformshift{5.014331in}{3.988523in}%
\pgfsys@useobject{currentmarker}{}%
\end{pgfscope}%
\end{pgfscope}%
\begin{pgfscope}%
\definecolor{textcolor}{rgb}{0.000000,0.000000,0.000000}%
\pgfsetstrokecolor{textcolor}%
\pgfsetfillcolor{textcolor}%
\pgftext[x=5.014331in,y=0.365069in,,top]{\color{textcolor}{\rmfamily\fontsize{10.000000}{12.000000}\selectfont\catcode`\^=\active\def^{\ifmmode\sp\else\^{}\fi}\catcode`\%=\active\def%{\%}$\mathdefault{0.8}$}}%
\end{pgfscope}%
\begin{pgfscope}%
\pgfsetbuttcap%
\pgfsetroundjoin%
\definecolor{currentfill}{rgb}{0.000000,0.000000,0.000000}%
\pgfsetfillcolor{currentfill}%
\pgfsetlinewidth{0.501875pt}%
\definecolor{currentstroke}{rgb}{0.000000,0.000000,0.000000}%
\pgfsetstrokecolor{currentstroke}%
\pgfsetdash{}{0pt}%
\pgfsys@defobject{currentmarker}{\pgfqpoint{0.000000in}{0.000000in}}{\pgfqpoint{0.000000in}{0.041667in}}{%
\pgfpathmoveto{\pgfqpoint{0.000000in}{0.000000in}}%
\pgfpathlineto{\pgfqpoint{0.000000in}{0.041667in}}%
\pgfusepath{stroke,fill}%
}%
\begin{pgfscope}%
\pgfsys@transformshift{6.155857in}{0.413680in}%
\pgfsys@useobject{currentmarker}{}%
\end{pgfscope}%
\end{pgfscope}%
\begin{pgfscope}%
\pgfsetbuttcap%
\pgfsetroundjoin%
\definecolor{currentfill}{rgb}{0.000000,0.000000,0.000000}%
\pgfsetfillcolor{currentfill}%
\pgfsetlinewidth{0.501875pt}%
\definecolor{currentstroke}{rgb}{0.000000,0.000000,0.000000}%
\pgfsetstrokecolor{currentstroke}%
\pgfsetdash{}{0pt}%
\pgfsys@defobject{currentmarker}{\pgfqpoint{0.000000in}{-0.041667in}}{\pgfqpoint{0.000000in}{0.000000in}}{%
\pgfpathmoveto{\pgfqpoint{0.000000in}{0.000000in}}%
\pgfpathlineto{\pgfqpoint{0.000000in}{-0.041667in}}%
\pgfusepath{stroke,fill}%
}%
\begin{pgfscope}%
\pgfsys@transformshift{6.155857in}{3.988523in}%
\pgfsys@useobject{currentmarker}{}%
\end{pgfscope}%
\end{pgfscope}%
\begin{pgfscope}%
\definecolor{textcolor}{rgb}{0.000000,0.000000,0.000000}%
\pgfsetstrokecolor{textcolor}%
\pgfsetfillcolor{textcolor}%
\pgftext[x=6.155857in,y=0.365069in,,top]{\color{textcolor}{\rmfamily\fontsize{10.000000}{12.000000}\selectfont\catcode`\^=\active\def^{\ifmmode\sp\else\^{}\fi}\catcode`\%=\active\def%{\%}$\mathdefault{1.0}$}}%
\end{pgfscope}%
\begin{pgfscope}%
\pgfsetbuttcap%
\pgfsetroundjoin%
\definecolor{currentfill}{rgb}{0.000000,0.000000,0.000000}%
\pgfsetfillcolor{currentfill}%
\pgfsetlinewidth{0.501875pt}%
\definecolor{currentstroke}{rgb}{0.000000,0.000000,0.000000}%
\pgfsetstrokecolor{currentstroke}%
\pgfsetdash{}{0pt}%
\pgfsys@defobject{currentmarker}{\pgfqpoint{0.000000in}{0.000000in}}{\pgfqpoint{0.000000in}{0.020833in}}{%
\pgfpathmoveto{\pgfqpoint{0.000000in}{0.000000in}}%
\pgfpathlineto{\pgfqpoint{0.000000in}{0.020833in}}%
\pgfusepath{stroke,fill}%
}%
\begin{pgfscope}%
\pgfsys@transformshift{0.733610in}{0.413680in}%
\pgfsys@useobject{currentmarker}{}%
\end{pgfscope}%
\end{pgfscope}%
\begin{pgfscope}%
\pgfsetbuttcap%
\pgfsetroundjoin%
\definecolor{currentfill}{rgb}{0.000000,0.000000,0.000000}%
\pgfsetfillcolor{currentfill}%
\pgfsetlinewidth{0.501875pt}%
\definecolor{currentstroke}{rgb}{0.000000,0.000000,0.000000}%
\pgfsetstrokecolor{currentstroke}%
\pgfsetdash{}{0pt}%
\pgfsys@defobject{currentmarker}{\pgfqpoint{0.000000in}{-0.020833in}}{\pgfqpoint{0.000000in}{0.000000in}}{%
\pgfpathmoveto{\pgfqpoint{0.000000in}{0.000000in}}%
\pgfpathlineto{\pgfqpoint{0.000000in}{-0.020833in}}%
\pgfusepath{stroke,fill}%
}%
\begin{pgfscope}%
\pgfsys@transformshift{0.733610in}{3.988523in}%
\pgfsys@useobject{currentmarker}{}%
\end{pgfscope}%
\end{pgfscope}%
\begin{pgfscope}%
\pgfsetbuttcap%
\pgfsetroundjoin%
\definecolor{currentfill}{rgb}{0.000000,0.000000,0.000000}%
\pgfsetfillcolor{currentfill}%
\pgfsetlinewidth{0.501875pt}%
\definecolor{currentstroke}{rgb}{0.000000,0.000000,0.000000}%
\pgfsetstrokecolor{currentstroke}%
\pgfsetdash{}{0pt}%
\pgfsys@defobject{currentmarker}{\pgfqpoint{0.000000in}{0.000000in}}{\pgfqpoint{0.000000in}{0.020833in}}{%
\pgfpathmoveto{\pgfqpoint{0.000000in}{0.000000in}}%
\pgfpathlineto{\pgfqpoint{0.000000in}{0.020833in}}%
\pgfusepath{stroke,fill}%
}%
\begin{pgfscope}%
\pgfsys@transformshift{1.018992in}{0.413680in}%
\pgfsys@useobject{currentmarker}{}%
\end{pgfscope}%
\end{pgfscope}%
\begin{pgfscope}%
\pgfsetbuttcap%
\pgfsetroundjoin%
\definecolor{currentfill}{rgb}{0.000000,0.000000,0.000000}%
\pgfsetfillcolor{currentfill}%
\pgfsetlinewidth{0.501875pt}%
\definecolor{currentstroke}{rgb}{0.000000,0.000000,0.000000}%
\pgfsetstrokecolor{currentstroke}%
\pgfsetdash{}{0pt}%
\pgfsys@defobject{currentmarker}{\pgfqpoint{0.000000in}{-0.020833in}}{\pgfqpoint{0.000000in}{0.000000in}}{%
\pgfpathmoveto{\pgfqpoint{0.000000in}{0.000000in}}%
\pgfpathlineto{\pgfqpoint{0.000000in}{-0.020833in}}%
\pgfusepath{stroke,fill}%
}%
\begin{pgfscope}%
\pgfsys@transformshift{1.018992in}{3.988523in}%
\pgfsys@useobject{currentmarker}{}%
\end{pgfscope}%
\end{pgfscope}%
\begin{pgfscope}%
\pgfsetbuttcap%
\pgfsetroundjoin%
\definecolor{currentfill}{rgb}{0.000000,0.000000,0.000000}%
\pgfsetfillcolor{currentfill}%
\pgfsetlinewidth{0.501875pt}%
\definecolor{currentstroke}{rgb}{0.000000,0.000000,0.000000}%
\pgfsetstrokecolor{currentstroke}%
\pgfsetdash{}{0pt}%
\pgfsys@defobject{currentmarker}{\pgfqpoint{0.000000in}{0.000000in}}{\pgfqpoint{0.000000in}{0.020833in}}{%
\pgfpathmoveto{\pgfqpoint{0.000000in}{0.000000in}}%
\pgfpathlineto{\pgfqpoint{0.000000in}{0.020833in}}%
\pgfusepath{stroke,fill}%
}%
\begin{pgfscope}%
\pgfsys@transformshift{1.304373in}{0.413680in}%
\pgfsys@useobject{currentmarker}{}%
\end{pgfscope}%
\end{pgfscope}%
\begin{pgfscope}%
\pgfsetbuttcap%
\pgfsetroundjoin%
\definecolor{currentfill}{rgb}{0.000000,0.000000,0.000000}%
\pgfsetfillcolor{currentfill}%
\pgfsetlinewidth{0.501875pt}%
\definecolor{currentstroke}{rgb}{0.000000,0.000000,0.000000}%
\pgfsetstrokecolor{currentstroke}%
\pgfsetdash{}{0pt}%
\pgfsys@defobject{currentmarker}{\pgfqpoint{0.000000in}{-0.020833in}}{\pgfqpoint{0.000000in}{0.000000in}}{%
\pgfpathmoveto{\pgfqpoint{0.000000in}{0.000000in}}%
\pgfpathlineto{\pgfqpoint{0.000000in}{-0.020833in}}%
\pgfusepath{stroke,fill}%
}%
\begin{pgfscope}%
\pgfsys@transformshift{1.304373in}{3.988523in}%
\pgfsys@useobject{currentmarker}{}%
\end{pgfscope}%
\end{pgfscope}%
\begin{pgfscope}%
\pgfsetbuttcap%
\pgfsetroundjoin%
\definecolor{currentfill}{rgb}{0.000000,0.000000,0.000000}%
\pgfsetfillcolor{currentfill}%
\pgfsetlinewidth{0.501875pt}%
\definecolor{currentstroke}{rgb}{0.000000,0.000000,0.000000}%
\pgfsetstrokecolor{currentstroke}%
\pgfsetdash{}{0pt}%
\pgfsys@defobject{currentmarker}{\pgfqpoint{0.000000in}{0.000000in}}{\pgfqpoint{0.000000in}{0.020833in}}{%
\pgfpathmoveto{\pgfqpoint{0.000000in}{0.000000in}}%
\pgfpathlineto{\pgfqpoint{0.000000in}{0.020833in}}%
\pgfusepath{stroke,fill}%
}%
\begin{pgfscope}%
\pgfsys@transformshift{1.875136in}{0.413680in}%
\pgfsys@useobject{currentmarker}{}%
\end{pgfscope}%
\end{pgfscope}%
\begin{pgfscope}%
\pgfsetbuttcap%
\pgfsetroundjoin%
\definecolor{currentfill}{rgb}{0.000000,0.000000,0.000000}%
\pgfsetfillcolor{currentfill}%
\pgfsetlinewidth{0.501875pt}%
\definecolor{currentstroke}{rgb}{0.000000,0.000000,0.000000}%
\pgfsetstrokecolor{currentstroke}%
\pgfsetdash{}{0pt}%
\pgfsys@defobject{currentmarker}{\pgfqpoint{0.000000in}{-0.020833in}}{\pgfqpoint{0.000000in}{0.000000in}}{%
\pgfpathmoveto{\pgfqpoint{0.000000in}{0.000000in}}%
\pgfpathlineto{\pgfqpoint{0.000000in}{-0.020833in}}%
\pgfusepath{stroke,fill}%
}%
\begin{pgfscope}%
\pgfsys@transformshift{1.875136in}{3.988523in}%
\pgfsys@useobject{currentmarker}{}%
\end{pgfscope}%
\end{pgfscope}%
\begin{pgfscope}%
\pgfsetbuttcap%
\pgfsetroundjoin%
\definecolor{currentfill}{rgb}{0.000000,0.000000,0.000000}%
\pgfsetfillcolor{currentfill}%
\pgfsetlinewidth{0.501875pt}%
\definecolor{currentstroke}{rgb}{0.000000,0.000000,0.000000}%
\pgfsetstrokecolor{currentstroke}%
\pgfsetdash{}{0pt}%
\pgfsys@defobject{currentmarker}{\pgfqpoint{0.000000in}{0.000000in}}{\pgfqpoint{0.000000in}{0.020833in}}{%
\pgfpathmoveto{\pgfqpoint{0.000000in}{0.000000in}}%
\pgfpathlineto{\pgfqpoint{0.000000in}{0.020833in}}%
\pgfusepath{stroke,fill}%
}%
\begin{pgfscope}%
\pgfsys@transformshift{2.160517in}{0.413680in}%
\pgfsys@useobject{currentmarker}{}%
\end{pgfscope}%
\end{pgfscope}%
\begin{pgfscope}%
\pgfsetbuttcap%
\pgfsetroundjoin%
\definecolor{currentfill}{rgb}{0.000000,0.000000,0.000000}%
\pgfsetfillcolor{currentfill}%
\pgfsetlinewidth{0.501875pt}%
\definecolor{currentstroke}{rgb}{0.000000,0.000000,0.000000}%
\pgfsetstrokecolor{currentstroke}%
\pgfsetdash{}{0pt}%
\pgfsys@defobject{currentmarker}{\pgfqpoint{0.000000in}{-0.020833in}}{\pgfqpoint{0.000000in}{0.000000in}}{%
\pgfpathmoveto{\pgfqpoint{0.000000in}{0.000000in}}%
\pgfpathlineto{\pgfqpoint{0.000000in}{-0.020833in}}%
\pgfusepath{stroke,fill}%
}%
\begin{pgfscope}%
\pgfsys@transformshift{2.160517in}{3.988523in}%
\pgfsys@useobject{currentmarker}{}%
\end{pgfscope}%
\end{pgfscope}%
\begin{pgfscope}%
\pgfsetbuttcap%
\pgfsetroundjoin%
\definecolor{currentfill}{rgb}{0.000000,0.000000,0.000000}%
\pgfsetfillcolor{currentfill}%
\pgfsetlinewidth{0.501875pt}%
\definecolor{currentstroke}{rgb}{0.000000,0.000000,0.000000}%
\pgfsetstrokecolor{currentstroke}%
\pgfsetdash{}{0pt}%
\pgfsys@defobject{currentmarker}{\pgfqpoint{0.000000in}{0.000000in}}{\pgfqpoint{0.000000in}{0.020833in}}{%
\pgfpathmoveto{\pgfqpoint{0.000000in}{0.000000in}}%
\pgfpathlineto{\pgfqpoint{0.000000in}{0.020833in}}%
\pgfusepath{stroke,fill}%
}%
\begin{pgfscope}%
\pgfsys@transformshift{2.445899in}{0.413680in}%
\pgfsys@useobject{currentmarker}{}%
\end{pgfscope}%
\end{pgfscope}%
\begin{pgfscope}%
\pgfsetbuttcap%
\pgfsetroundjoin%
\definecolor{currentfill}{rgb}{0.000000,0.000000,0.000000}%
\pgfsetfillcolor{currentfill}%
\pgfsetlinewidth{0.501875pt}%
\definecolor{currentstroke}{rgb}{0.000000,0.000000,0.000000}%
\pgfsetstrokecolor{currentstroke}%
\pgfsetdash{}{0pt}%
\pgfsys@defobject{currentmarker}{\pgfqpoint{0.000000in}{-0.020833in}}{\pgfqpoint{0.000000in}{0.000000in}}{%
\pgfpathmoveto{\pgfqpoint{0.000000in}{0.000000in}}%
\pgfpathlineto{\pgfqpoint{0.000000in}{-0.020833in}}%
\pgfusepath{stroke,fill}%
}%
\begin{pgfscope}%
\pgfsys@transformshift{2.445899in}{3.988523in}%
\pgfsys@useobject{currentmarker}{}%
\end{pgfscope}%
\end{pgfscope}%
\begin{pgfscope}%
\pgfsetbuttcap%
\pgfsetroundjoin%
\definecolor{currentfill}{rgb}{0.000000,0.000000,0.000000}%
\pgfsetfillcolor{currentfill}%
\pgfsetlinewidth{0.501875pt}%
\definecolor{currentstroke}{rgb}{0.000000,0.000000,0.000000}%
\pgfsetstrokecolor{currentstroke}%
\pgfsetdash{}{0pt}%
\pgfsys@defobject{currentmarker}{\pgfqpoint{0.000000in}{0.000000in}}{\pgfqpoint{0.000000in}{0.020833in}}{%
\pgfpathmoveto{\pgfqpoint{0.000000in}{0.000000in}}%
\pgfpathlineto{\pgfqpoint{0.000000in}{0.020833in}}%
\pgfusepath{stroke,fill}%
}%
\begin{pgfscope}%
\pgfsys@transformshift{3.016661in}{0.413680in}%
\pgfsys@useobject{currentmarker}{}%
\end{pgfscope}%
\end{pgfscope}%
\begin{pgfscope}%
\pgfsetbuttcap%
\pgfsetroundjoin%
\definecolor{currentfill}{rgb}{0.000000,0.000000,0.000000}%
\pgfsetfillcolor{currentfill}%
\pgfsetlinewidth{0.501875pt}%
\definecolor{currentstroke}{rgb}{0.000000,0.000000,0.000000}%
\pgfsetstrokecolor{currentstroke}%
\pgfsetdash{}{0pt}%
\pgfsys@defobject{currentmarker}{\pgfqpoint{0.000000in}{-0.020833in}}{\pgfqpoint{0.000000in}{0.000000in}}{%
\pgfpathmoveto{\pgfqpoint{0.000000in}{0.000000in}}%
\pgfpathlineto{\pgfqpoint{0.000000in}{-0.020833in}}%
\pgfusepath{stroke,fill}%
}%
\begin{pgfscope}%
\pgfsys@transformshift{3.016661in}{3.988523in}%
\pgfsys@useobject{currentmarker}{}%
\end{pgfscope}%
\end{pgfscope}%
\begin{pgfscope}%
\pgfsetbuttcap%
\pgfsetroundjoin%
\definecolor{currentfill}{rgb}{0.000000,0.000000,0.000000}%
\pgfsetfillcolor{currentfill}%
\pgfsetlinewidth{0.501875pt}%
\definecolor{currentstroke}{rgb}{0.000000,0.000000,0.000000}%
\pgfsetstrokecolor{currentstroke}%
\pgfsetdash{}{0pt}%
\pgfsys@defobject{currentmarker}{\pgfqpoint{0.000000in}{0.000000in}}{\pgfqpoint{0.000000in}{0.020833in}}{%
\pgfpathmoveto{\pgfqpoint{0.000000in}{0.000000in}}%
\pgfpathlineto{\pgfqpoint{0.000000in}{0.020833in}}%
\pgfusepath{stroke,fill}%
}%
\begin{pgfscope}%
\pgfsys@transformshift{3.302043in}{0.413680in}%
\pgfsys@useobject{currentmarker}{}%
\end{pgfscope}%
\end{pgfscope}%
\begin{pgfscope}%
\pgfsetbuttcap%
\pgfsetroundjoin%
\definecolor{currentfill}{rgb}{0.000000,0.000000,0.000000}%
\pgfsetfillcolor{currentfill}%
\pgfsetlinewidth{0.501875pt}%
\definecolor{currentstroke}{rgb}{0.000000,0.000000,0.000000}%
\pgfsetstrokecolor{currentstroke}%
\pgfsetdash{}{0pt}%
\pgfsys@defobject{currentmarker}{\pgfqpoint{0.000000in}{-0.020833in}}{\pgfqpoint{0.000000in}{0.000000in}}{%
\pgfpathmoveto{\pgfqpoint{0.000000in}{0.000000in}}%
\pgfpathlineto{\pgfqpoint{0.000000in}{-0.020833in}}%
\pgfusepath{stroke,fill}%
}%
\begin{pgfscope}%
\pgfsys@transformshift{3.302043in}{3.988523in}%
\pgfsys@useobject{currentmarker}{}%
\end{pgfscope}%
\end{pgfscope}%
\begin{pgfscope}%
\pgfsetbuttcap%
\pgfsetroundjoin%
\definecolor{currentfill}{rgb}{0.000000,0.000000,0.000000}%
\pgfsetfillcolor{currentfill}%
\pgfsetlinewidth{0.501875pt}%
\definecolor{currentstroke}{rgb}{0.000000,0.000000,0.000000}%
\pgfsetstrokecolor{currentstroke}%
\pgfsetdash{}{0pt}%
\pgfsys@defobject{currentmarker}{\pgfqpoint{0.000000in}{0.000000in}}{\pgfqpoint{0.000000in}{0.020833in}}{%
\pgfpathmoveto{\pgfqpoint{0.000000in}{0.000000in}}%
\pgfpathlineto{\pgfqpoint{0.000000in}{0.020833in}}%
\pgfusepath{stroke,fill}%
}%
\begin{pgfscope}%
\pgfsys@transformshift{3.587424in}{0.413680in}%
\pgfsys@useobject{currentmarker}{}%
\end{pgfscope}%
\end{pgfscope}%
\begin{pgfscope}%
\pgfsetbuttcap%
\pgfsetroundjoin%
\definecolor{currentfill}{rgb}{0.000000,0.000000,0.000000}%
\pgfsetfillcolor{currentfill}%
\pgfsetlinewidth{0.501875pt}%
\definecolor{currentstroke}{rgb}{0.000000,0.000000,0.000000}%
\pgfsetstrokecolor{currentstroke}%
\pgfsetdash{}{0pt}%
\pgfsys@defobject{currentmarker}{\pgfqpoint{0.000000in}{-0.020833in}}{\pgfqpoint{0.000000in}{0.000000in}}{%
\pgfpathmoveto{\pgfqpoint{0.000000in}{0.000000in}}%
\pgfpathlineto{\pgfqpoint{0.000000in}{-0.020833in}}%
\pgfusepath{stroke,fill}%
}%
\begin{pgfscope}%
\pgfsys@transformshift{3.587424in}{3.988523in}%
\pgfsys@useobject{currentmarker}{}%
\end{pgfscope}%
\end{pgfscope}%
\begin{pgfscope}%
\pgfsetbuttcap%
\pgfsetroundjoin%
\definecolor{currentfill}{rgb}{0.000000,0.000000,0.000000}%
\pgfsetfillcolor{currentfill}%
\pgfsetlinewidth{0.501875pt}%
\definecolor{currentstroke}{rgb}{0.000000,0.000000,0.000000}%
\pgfsetstrokecolor{currentstroke}%
\pgfsetdash{}{0pt}%
\pgfsys@defobject{currentmarker}{\pgfqpoint{0.000000in}{0.000000in}}{\pgfqpoint{0.000000in}{0.020833in}}{%
\pgfpathmoveto{\pgfqpoint{0.000000in}{0.000000in}}%
\pgfpathlineto{\pgfqpoint{0.000000in}{0.020833in}}%
\pgfusepath{stroke,fill}%
}%
\begin{pgfscope}%
\pgfsys@transformshift{4.158187in}{0.413680in}%
\pgfsys@useobject{currentmarker}{}%
\end{pgfscope}%
\end{pgfscope}%
\begin{pgfscope}%
\pgfsetbuttcap%
\pgfsetroundjoin%
\definecolor{currentfill}{rgb}{0.000000,0.000000,0.000000}%
\pgfsetfillcolor{currentfill}%
\pgfsetlinewidth{0.501875pt}%
\definecolor{currentstroke}{rgb}{0.000000,0.000000,0.000000}%
\pgfsetstrokecolor{currentstroke}%
\pgfsetdash{}{0pt}%
\pgfsys@defobject{currentmarker}{\pgfqpoint{0.000000in}{-0.020833in}}{\pgfqpoint{0.000000in}{0.000000in}}{%
\pgfpathmoveto{\pgfqpoint{0.000000in}{0.000000in}}%
\pgfpathlineto{\pgfqpoint{0.000000in}{-0.020833in}}%
\pgfusepath{stroke,fill}%
}%
\begin{pgfscope}%
\pgfsys@transformshift{4.158187in}{3.988523in}%
\pgfsys@useobject{currentmarker}{}%
\end{pgfscope}%
\end{pgfscope}%
\begin{pgfscope}%
\pgfsetbuttcap%
\pgfsetroundjoin%
\definecolor{currentfill}{rgb}{0.000000,0.000000,0.000000}%
\pgfsetfillcolor{currentfill}%
\pgfsetlinewidth{0.501875pt}%
\definecolor{currentstroke}{rgb}{0.000000,0.000000,0.000000}%
\pgfsetstrokecolor{currentstroke}%
\pgfsetdash{}{0pt}%
\pgfsys@defobject{currentmarker}{\pgfqpoint{0.000000in}{0.000000in}}{\pgfqpoint{0.000000in}{0.020833in}}{%
\pgfpathmoveto{\pgfqpoint{0.000000in}{0.000000in}}%
\pgfpathlineto{\pgfqpoint{0.000000in}{0.020833in}}%
\pgfusepath{stroke,fill}%
}%
\begin{pgfscope}%
\pgfsys@transformshift{4.443568in}{0.413680in}%
\pgfsys@useobject{currentmarker}{}%
\end{pgfscope}%
\end{pgfscope}%
\begin{pgfscope}%
\pgfsetbuttcap%
\pgfsetroundjoin%
\definecolor{currentfill}{rgb}{0.000000,0.000000,0.000000}%
\pgfsetfillcolor{currentfill}%
\pgfsetlinewidth{0.501875pt}%
\definecolor{currentstroke}{rgb}{0.000000,0.000000,0.000000}%
\pgfsetstrokecolor{currentstroke}%
\pgfsetdash{}{0pt}%
\pgfsys@defobject{currentmarker}{\pgfqpoint{0.000000in}{-0.020833in}}{\pgfqpoint{0.000000in}{0.000000in}}{%
\pgfpathmoveto{\pgfqpoint{0.000000in}{0.000000in}}%
\pgfpathlineto{\pgfqpoint{0.000000in}{-0.020833in}}%
\pgfusepath{stroke,fill}%
}%
\begin{pgfscope}%
\pgfsys@transformshift{4.443568in}{3.988523in}%
\pgfsys@useobject{currentmarker}{}%
\end{pgfscope}%
\end{pgfscope}%
\begin{pgfscope}%
\pgfsetbuttcap%
\pgfsetroundjoin%
\definecolor{currentfill}{rgb}{0.000000,0.000000,0.000000}%
\pgfsetfillcolor{currentfill}%
\pgfsetlinewidth{0.501875pt}%
\definecolor{currentstroke}{rgb}{0.000000,0.000000,0.000000}%
\pgfsetstrokecolor{currentstroke}%
\pgfsetdash{}{0pt}%
\pgfsys@defobject{currentmarker}{\pgfqpoint{0.000000in}{0.000000in}}{\pgfqpoint{0.000000in}{0.020833in}}{%
\pgfpathmoveto{\pgfqpoint{0.000000in}{0.000000in}}%
\pgfpathlineto{\pgfqpoint{0.000000in}{0.020833in}}%
\pgfusepath{stroke,fill}%
}%
\begin{pgfscope}%
\pgfsys@transformshift{4.728950in}{0.413680in}%
\pgfsys@useobject{currentmarker}{}%
\end{pgfscope}%
\end{pgfscope}%
\begin{pgfscope}%
\pgfsetbuttcap%
\pgfsetroundjoin%
\definecolor{currentfill}{rgb}{0.000000,0.000000,0.000000}%
\pgfsetfillcolor{currentfill}%
\pgfsetlinewidth{0.501875pt}%
\definecolor{currentstroke}{rgb}{0.000000,0.000000,0.000000}%
\pgfsetstrokecolor{currentstroke}%
\pgfsetdash{}{0pt}%
\pgfsys@defobject{currentmarker}{\pgfqpoint{0.000000in}{-0.020833in}}{\pgfqpoint{0.000000in}{0.000000in}}{%
\pgfpathmoveto{\pgfqpoint{0.000000in}{0.000000in}}%
\pgfpathlineto{\pgfqpoint{0.000000in}{-0.020833in}}%
\pgfusepath{stroke,fill}%
}%
\begin{pgfscope}%
\pgfsys@transformshift{4.728950in}{3.988523in}%
\pgfsys@useobject{currentmarker}{}%
\end{pgfscope}%
\end{pgfscope}%
\begin{pgfscope}%
\pgfsetbuttcap%
\pgfsetroundjoin%
\definecolor{currentfill}{rgb}{0.000000,0.000000,0.000000}%
\pgfsetfillcolor{currentfill}%
\pgfsetlinewidth{0.501875pt}%
\definecolor{currentstroke}{rgb}{0.000000,0.000000,0.000000}%
\pgfsetstrokecolor{currentstroke}%
\pgfsetdash{}{0pt}%
\pgfsys@defobject{currentmarker}{\pgfqpoint{0.000000in}{0.000000in}}{\pgfqpoint{0.000000in}{0.020833in}}{%
\pgfpathmoveto{\pgfqpoint{0.000000in}{0.000000in}}%
\pgfpathlineto{\pgfqpoint{0.000000in}{0.020833in}}%
\pgfusepath{stroke,fill}%
}%
\begin{pgfscope}%
\pgfsys@transformshift{5.299713in}{0.413680in}%
\pgfsys@useobject{currentmarker}{}%
\end{pgfscope}%
\end{pgfscope}%
\begin{pgfscope}%
\pgfsetbuttcap%
\pgfsetroundjoin%
\definecolor{currentfill}{rgb}{0.000000,0.000000,0.000000}%
\pgfsetfillcolor{currentfill}%
\pgfsetlinewidth{0.501875pt}%
\definecolor{currentstroke}{rgb}{0.000000,0.000000,0.000000}%
\pgfsetstrokecolor{currentstroke}%
\pgfsetdash{}{0pt}%
\pgfsys@defobject{currentmarker}{\pgfqpoint{0.000000in}{-0.020833in}}{\pgfqpoint{0.000000in}{0.000000in}}{%
\pgfpathmoveto{\pgfqpoint{0.000000in}{0.000000in}}%
\pgfpathlineto{\pgfqpoint{0.000000in}{-0.020833in}}%
\pgfusepath{stroke,fill}%
}%
\begin{pgfscope}%
\pgfsys@transformshift{5.299713in}{3.988523in}%
\pgfsys@useobject{currentmarker}{}%
\end{pgfscope}%
\end{pgfscope}%
\begin{pgfscope}%
\pgfsetbuttcap%
\pgfsetroundjoin%
\definecolor{currentfill}{rgb}{0.000000,0.000000,0.000000}%
\pgfsetfillcolor{currentfill}%
\pgfsetlinewidth{0.501875pt}%
\definecolor{currentstroke}{rgb}{0.000000,0.000000,0.000000}%
\pgfsetstrokecolor{currentstroke}%
\pgfsetdash{}{0pt}%
\pgfsys@defobject{currentmarker}{\pgfqpoint{0.000000in}{0.000000in}}{\pgfqpoint{0.000000in}{0.020833in}}{%
\pgfpathmoveto{\pgfqpoint{0.000000in}{0.000000in}}%
\pgfpathlineto{\pgfqpoint{0.000000in}{0.020833in}}%
\pgfusepath{stroke,fill}%
}%
\begin{pgfscope}%
\pgfsys@transformshift{5.585094in}{0.413680in}%
\pgfsys@useobject{currentmarker}{}%
\end{pgfscope}%
\end{pgfscope}%
\begin{pgfscope}%
\pgfsetbuttcap%
\pgfsetroundjoin%
\definecolor{currentfill}{rgb}{0.000000,0.000000,0.000000}%
\pgfsetfillcolor{currentfill}%
\pgfsetlinewidth{0.501875pt}%
\definecolor{currentstroke}{rgb}{0.000000,0.000000,0.000000}%
\pgfsetstrokecolor{currentstroke}%
\pgfsetdash{}{0pt}%
\pgfsys@defobject{currentmarker}{\pgfqpoint{0.000000in}{-0.020833in}}{\pgfqpoint{0.000000in}{0.000000in}}{%
\pgfpathmoveto{\pgfqpoint{0.000000in}{0.000000in}}%
\pgfpathlineto{\pgfqpoint{0.000000in}{-0.020833in}}%
\pgfusepath{stroke,fill}%
}%
\begin{pgfscope}%
\pgfsys@transformshift{5.585094in}{3.988523in}%
\pgfsys@useobject{currentmarker}{}%
\end{pgfscope}%
\end{pgfscope}%
\begin{pgfscope}%
\pgfsetbuttcap%
\pgfsetroundjoin%
\definecolor{currentfill}{rgb}{0.000000,0.000000,0.000000}%
\pgfsetfillcolor{currentfill}%
\pgfsetlinewidth{0.501875pt}%
\definecolor{currentstroke}{rgb}{0.000000,0.000000,0.000000}%
\pgfsetstrokecolor{currentstroke}%
\pgfsetdash{}{0pt}%
\pgfsys@defobject{currentmarker}{\pgfqpoint{0.000000in}{0.000000in}}{\pgfqpoint{0.000000in}{0.020833in}}{%
\pgfpathmoveto{\pgfqpoint{0.000000in}{0.000000in}}%
\pgfpathlineto{\pgfqpoint{0.000000in}{0.020833in}}%
\pgfusepath{stroke,fill}%
}%
\begin{pgfscope}%
\pgfsys@transformshift{5.870475in}{0.413680in}%
\pgfsys@useobject{currentmarker}{}%
\end{pgfscope}%
\end{pgfscope}%
\begin{pgfscope}%
\pgfsetbuttcap%
\pgfsetroundjoin%
\definecolor{currentfill}{rgb}{0.000000,0.000000,0.000000}%
\pgfsetfillcolor{currentfill}%
\pgfsetlinewidth{0.501875pt}%
\definecolor{currentstroke}{rgb}{0.000000,0.000000,0.000000}%
\pgfsetstrokecolor{currentstroke}%
\pgfsetdash{}{0pt}%
\pgfsys@defobject{currentmarker}{\pgfqpoint{0.000000in}{-0.020833in}}{\pgfqpoint{0.000000in}{0.000000in}}{%
\pgfpathmoveto{\pgfqpoint{0.000000in}{0.000000in}}%
\pgfpathlineto{\pgfqpoint{0.000000in}{-0.020833in}}%
\pgfusepath{stroke,fill}%
}%
\begin{pgfscope}%
\pgfsys@transformshift{5.870475in}{3.988523in}%
\pgfsys@useobject{currentmarker}{}%
\end{pgfscope}%
\end{pgfscope}%
\begin{pgfscope}%
\definecolor{textcolor}{rgb}{0.000000,0.000000,0.000000}%
\pgfsetstrokecolor{textcolor}%
\pgfsetfillcolor{textcolor}%
\pgftext[x=3.302043in,y=0.179757in,,top]{\color{textcolor}{\rmfamily\fontsize{10.000000}{12.000000}\selectfont\catcode`\^=\active\def^{\ifmmode\sp\else\^{}\fi}\catcode`\%=\active\def%{\%}$t$}}%
\end{pgfscope}%
\begin{pgfscope}%
\pgfsetbuttcap%
\pgfsetroundjoin%
\definecolor{currentfill}{rgb}{0.000000,0.000000,0.000000}%
\pgfsetfillcolor{currentfill}%
\pgfsetlinewidth{0.501875pt}%
\definecolor{currentstroke}{rgb}{0.000000,0.000000,0.000000}%
\pgfsetstrokecolor{currentstroke}%
\pgfsetdash{}{0pt}%
\pgfsys@defobject{currentmarker}{\pgfqpoint{0.000000in}{0.000000in}}{\pgfqpoint{0.041667in}{0.000000in}}{%
\pgfpathmoveto{\pgfqpoint{0.000000in}{0.000000in}}%
\pgfpathlineto{\pgfqpoint{0.041667in}{0.000000in}}%
\pgfusepath{stroke,fill}%
}%
\begin{pgfscope}%
\pgfsys@transformshift{0.448229in}{0.635932in}%
\pgfsys@useobject{currentmarker}{}%
\end{pgfscope}%
\end{pgfscope}%
\begin{pgfscope}%
\pgfsetbuttcap%
\pgfsetroundjoin%
\definecolor{currentfill}{rgb}{0.000000,0.000000,0.000000}%
\pgfsetfillcolor{currentfill}%
\pgfsetlinewidth{0.501875pt}%
\definecolor{currentstroke}{rgb}{0.000000,0.000000,0.000000}%
\pgfsetstrokecolor{currentstroke}%
\pgfsetdash{}{0pt}%
\pgfsys@defobject{currentmarker}{\pgfqpoint{-0.041667in}{0.000000in}}{\pgfqpoint{-0.000000in}{0.000000in}}{%
\pgfpathmoveto{\pgfqpoint{-0.000000in}{0.000000in}}%
\pgfpathlineto{\pgfqpoint{-0.041667in}{0.000000in}}%
\pgfusepath{stroke,fill}%
}%
\begin{pgfscope}%
\pgfsys@transformshift{6.155857in}{0.635932in}%
\pgfsys@useobject{currentmarker}{}%
\end{pgfscope}%
\end{pgfscope}%
\begin{pgfscope}%
\definecolor{textcolor}{rgb}{0.000000,0.000000,0.000000}%
\pgfsetstrokecolor{textcolor}%
\pgfsetfillcolor{textcolor}%
\pgftext[x=0.235312in, y=0.587234in, left, base]{\color{textcolor}{\rmfamily\fontsize{10.000000}{12.000000}\selectfont\catcode`\^=\active\def^{\ifmmode\sp\else\^{}\fi}\catcode`\%=\active\def%{\%}$\mathdefault{1.0}$}}%
\end{pgfscope}%
\begin{pgfscope}%
\pgfsetbuttcap%
\pgfsetroundjoin%
\definecolor{currentfill}{rgb}{0.000000,0.000000,0.000000}%
\pgfsetfillcolor{currentfill}%
\pgfsetlinewidth{0.501875pt}%
\definecolor{currentstroke}{rgb}{0.000000,0.000000,0.000000}%
\pgfsetstrokecolor{currentstroke}%
\pgfsetdash{}{0pt}%
\pgfsys@defobject{currentmarker}{\pgfqpoint{0.000000in}{0.000000in}}{\pgfqpoint{0.041667in}{0.000000in}}{%
\pgfpathmoveto{\pgfqpoint{0.000000in}{0.000000in}}%
\pgfpathlineto{\pgfqpoint{0.041667in}{0.000000in}}%
\pgfusepath{stroke,fill}%
}%
\begin{pgfscope}%
\pgfsys@transformshift{0.448229in}{1.235169in}%
\pgfsys@useobject{currentmarker}{}%
\end{pgfscope}%
\end{pgfscope}%
\begin{pgfscope}%
\pgfsetbuttcap%
\pgfsetroundjoin%
\definecolor{currentfill}{rgb}{0.000000,0.000000,0.000000}%
\pgfsetfillcolor{currentfill}%
\pgfsetlinewidth{0.501875pt}%
\definecolor{currentstroke}{rgb}{0.000000,0.000000,0.000000}%
\pgfsetstrokecolor{currentstroke}%
\pgfsetdash{}{0pt}%
\pgfsys@defobject{currentmarker}{\pgfqpoint{-0.041667in}{0.000000in}}{\pgfqpoint{-0.000000in}{0.000000in}}{%
\pgfpathmoveto{\pgfqpoint{-0.000000in}{0.000000in}}%
\pgfpathlineto{\pgfqpoint{-0.041667in}{0.000000in}}%
\pgfusepath{stroke,fill}%
}%
\begin{pgfscope}%
\pgfsys@transformshift{6.155857in}{1.235169in}%
\pgfsys@useobject{currentmarker}{}%
\end{pgfscope}%
\end{pgfscope}%
\begin{pgfscope}%
\definecolor{textcolor}{rgb}{0.000000,0.000000,0.000000}%
\pgfsetstrokecolor{textcolor}%
\pgfsetfillcolor{textcolor}%
\pgftext[x=0.235312in, y=1.186471in, left, base]{\color{textcolor}{\rmfamily\fontsize{10.000000}{12.000000}\selectfont\catcode`\^=\active\def^{\ifmmode\sp\else\^{}\fi}\catcode`\%=\active\def%{\%}$\mathdefault{1.5}$}}%
\end{pgfscope}%
\begin{pgfscope}%
\pgfsetbuttcap%
\pgfsetroundjoin%
\definecolor{currentfill}{rgb}{0.000000,0.000000,0.000000}%
\pgfsetfillcolor{currentfill}%
\pgfsetlinewidth{0.501875pt}%
\definecolor{currentstroke}{rgb}{0.000000,0.000000,0.000000}%
\pgfsetstrokecolor{currentstroke}%
\pgfsetdash{}{0pt}%
\pgfsys@defobject{currentmarker}{\pgfqpoint{0.000000in}{0.000000in}}{\pgfqpoint{0.041667in}{0.000000in}}{%
\pgfpathmoveto{\pgfqpoint{0.000000in}{0.000000in}}%
\pgfpathlineto{\pgfqpoint{0.041667in}{0.000000in}}%
\pgfusepath{stroke,fill}%
}%
\begin{pgfscope}%
\pgfsys@transformshift{0.448229in}{1.834407in}%
\pgfsys@useobject{currentmarker}{}%
\end{pgfscope}%
\end{pgfscope}%
\begin{pgfscope}%
\pgfsetbuttcap%
\pgfsetroundjoin%
\definecolor{currentfill}{rgb}{0.000000,0.000000,0.000000}%
\pgfsetfillcolor{currentfill}%
\pgfsetlinewidth{0.501875pt}%
\definecolor{currentstroke}{rgb}{0.000000,0.000000,0.000000}%
\pgfsetstrokecolor{currentstroke}%
\pgfsetdash{}{0pt}%
\pgfsys@defobject{currentmarker}{\pgfqpoint{-0.041667in}{0.000000in}}{\pgfqpoint{-0.000000in}{0.000000in}}{%
\pgfpathmoveto{\pgfqpoint{-0.000000in}{0.000000in}}%
\pgfpathlineto{\pgfqpoint{-0.041667in}{0.000000in}}%
\pgfusepath{stroke,fill}%
}%
\begin{pgfscope}%
\pgfsys@transformshift{6.155857in}{1.834407in}%
\pgfsys@useobject{currentmarker}{}%
\end{pgfscope}%
\end{pgfscope}%
\begin{pgfscope}%
\definecolor{textcolor}{rgb}{0.000000,0.000000,0.000000}%
\pgfsetstrokecolor{textcolor}%
\pgfsetfillcolor{textcolor}%
\pgftext[x=0.235312in, y=1.785709in, left, base]{\color{textcolor}{\rmfamily\fontsize{10.000000}{12.000000}\selectfont\catcode`\^=\active\def^{\ifmmode\sp\else\^{}\fi}\catcode`\%=\active\def%{\%}$\mathdefault{2.0}$}}%
\end{pgfscope}%
\begin{pgfscope}%
\pgfsetbuttcap%
\pgfsetroundjoin%
\definecolor{currentfill}{rgb}{0.000000,0.000000,0.000000}%
\pgfsetfillcolor{currentfill}%
\pgfsetlinewidth{0.501875pt}%
\definecolor{currentstroke}{rgb}{0.000000,0.000000,0.000000}%
\pgfsetstrokecolor{currentstroke}%
\pgfsetdash{}{0pt}%
\pgfsys@defobject{currentmarker}{\pgfqpoint{0.000000in}{0.000000in}}{\pgfqpoint{0.041667in}{0.000000in}}{%
\pgfpathmoveto{\pgfqpoint{0.000000in}{0.000000in}}%
\pgfpathlineto{\pgfqpoint{0.041667in}{0.000000in}}%
\pgfusepath{stroke,fill}%
}%
\begin{pgfscope}%
\pgfsys@transformshift{0.448229in}{2.433644in}%
\pgfsys@useobject{currentmarker}{}%
\end{pgfscope}%
\end{pgfscope}%
\begin{pgfscope}%
\pgfsetbuttcap%
\pgfsetroundjoin%
\definecolor{currentfill}{rgb}{0.000000,0.000000,0.000000}%
\pgfsetfillcolor{currentfill}%
\pgfsetlinewidth{0.501875pt}%
\definecolor{currentstroke}{rgb}{0.000000,0.000000,0.000000}%
\pgfsetstrokecolor{currentstroke}%
\pgfsetdash{}{0pt}%
\pgfsys@defobject{currentmarker}{\pgfqpoint{-0.041667in}{0.000000in}}{\pgfqpoint{-0.000000in}{0.000000in}}{%
\pgfpathmoveto{\pgfqpoint{-0.000000in}{0.000000in}}%
\pgfpathlineto{\pgfqpoint{-0.041667in}{0.000000in}}%
\pgfusepath{stroke,fill}%
}%
\begin{pgfscope}%
\pgfsys@transformshift{6.155857in}{2.433644in}%
\pgfsys@useobject{currentmarker}{}%
\end{pgfscope}%
\end{pgfscope}%
\begin{pgfscope}%
\definecolor{textcolor}{rgb}{0.000000,0.000000,0.000000}%
\pgfsetstrokecolor{textcolor}%
\pgfsetfillcolor{textcolor}%
\pgftext[x=0.235312in, y=2.384946in, left, base]{\color{textcolor}{\rmfamily\fontsize{10.000000}{12.000000}\selectfont\catcode`\^=\active\def^{\ifmmode\sp\else\^{}\fi}\catcode`\%=\active\def%{\%}$\mathdefault{2.5}$}}%
\end{pgfscope}%
\begin{pgfscope}%
\pgfsetbuttcap%
\pgfsetroundjoin%
\definecolor{currentfill}{rgb}{0.000000,0.000000,0.000000}%
\pgfsetfillcolor{currentfill}%
\pgfsetlinewidth{0.501875pt}%
\definecolor{currentstroke}{rgb}{0.000000,0.000000,0.000000}%
\pgfsetstrokecolor{currentstroke}%
\pgfsetdash{}{0pt}%
\pgfsys@defobject{currentmarker}{\pgfqpoint{0.000000in}{0.000000in}}{\pgfqpoint{0.041667in}{0.000000in}}{%
\pgfpathmoveto{\pgfqpoint{0.000000in}{0.000000in}}%
\pgfpathlineto{\pgfqpoint{0.041667in}{0.000000in}}%
\pgfusepath{stroke,fill}%
}%
\begin{pgfscope}%
\pgfsys@transformshift{0.448229in}{3.032881in}%
\pgfsys@useobject{currentmarker}{}%
\end{pgfscope}%
\end{pgfscope}%
\begin{pgfscope}%
\pgfsetbuttcap%
\pgfsetroundjoin%
\definecolor{currentfill}{rgb}{0.000000,0.000000,0.000000}%
\pgfsetfillcolor{currentfill}%
\pgfsetlinewidth{0.501875pt}%
\definecolor{currentstroke}{rgb}{0.000000,0.000000,0.000000}%
\pgfsetstrokecolor{currentstroke}%
\pgfsetdash{}{0pt}%
\pgfsys@defobject{currentmarker}{\pgfqpoint{-0.041667in}{0.000000in}}{\pgfqpoint{-0.000000in}{0.000000in}}{%
\pgfpathmoveto{\pgfqpoint{-0.000000in}{0.000000in}}%
\pgfpathlineto{\pgfqpoint{-0.041667in}{0.000000in}}%
\pgfusepath{stroke,fill}%
}%
\begin{pgfscope}%
\pgfsys@transformshift{6.155857in}{3.032881in}%
\pgfsys@useobject{currentmarker}{}%
\end{pgfscope}%
\end{pgfscope}%
\begin{pgfscope}%
\definecolor{textcolor}{rgb}{0.000000,0.000000,0.000000}%
\pgfsetstrokecolor{textcolor}%
\pgfsetfillcolor{textcolor}%
\pgftext[x=0.235312in, y=2.984183in, left, base]{\color{textcolor}{\rmfamily\fontsize{10.000000}{12.000000}\selectfont\catcode`\^=\active\def^{\ifmmode\sp\else\^{}\fi}\catcode`\%=\active\def%{\%}$\mathdefault{3.0}$}}%
\end{pgfscope}%
\begin{pgfscope}%
\pgfsetbuttcap%
\pgfsetroundjoin%
\definecolor{currentfill}{rgb}{0.000000,0.000000,0.000000}%
\pgfsetfillcolor{currentfill}%
\pgfsetlinewidth{0.501875pt}%
\definecolor{currentstroke}{rgb}{0.000000,0.000000,0.000000}%
\pgfsetstrokecolor{currentstroke}%
\pgfsetdash{}{0pt}%
\pgfsys@defobject{currentmarker}{\pgfqpoint{0.000000in}{0.000000in}}{\pgfqpoint{0.041667in}{0.000000in}}{%
\pgfpathmoveto{\pgfqpoint{0.000000in}{0.000000in}}%
\pgfpathlineto{\pgfqpoint{0.041667in}{0.000000in}}%
\pgfusepath{stroke,fill}%
}%
\begin{pgfscope}%
\pgfsys@transformshift{0.448229in}{3.632118in}%
\pgfsys@useobject{currentmarker}{}%
\end{pgfscope}%
\end{pgfscope}%
\begin{pgfscope}%
\pgfsetbuttcap%
\pgfsetroundjoin%
\definecolor{currentfill}{rgb}{0.000000,0.000000,0.000000}%
\pgfsetfillcolor{currentfill}%
\pgfsetlinewidth{0.501875pt}%
\definecolor{currentstroke}{rgb}{0.000000,0.000000,0.000000}%
\pgfsetstrokecolor{currentstroke}%
\pgfsetdash{}{0pt}%
\pgfsys@defobject{currentmarker}{\pgfqpoint{-0.041667in}{0.000000in}}{\pgfqpoint{-0.000000in}{0.000000in}}{%
\pgfpathmoveto{\pgfqpoint{-0.000000in}{0.000000in}}%
\pgfpathlineto{\pgfqpoint{-0.041667in}{0.000000in}}%
\pgfusepath{stroke,fill}%
}%
\begin{pgfscope}%
\pgfsys@transformshift{6.155857in}{3.632118in}%
\pgfsys@useobject{currentmarker}{}%
\end{pgfscope}%
\end{pgfscope}%
\begin{pgfscope}%
\definecolor{textcolor}{rgb}{0.000000,0.000000,0.000000}%
\pgfsetstrokecolor{textcolor}%
\pgfsetfillcolor{textcolor}%
\pgftext[x=0.235312in, y=3.583420in, left, base]{\color{textcolor}{\rmfamily\fontsize{10.000000}{12.000000}\selectfont\catcode`\^=\active\def^{\ifmmode\sp\else\^{}\fi}\catcode`\%=\active\def%{\%}$\mathdefault{3.5}$}}%
\end{pgfscope}%
\begin{pgfscope}%
\pgfsetbuttcap%
\pgfsetroundjoin%
\definecolor{currentfill}{rgb}{0.000000,0.000000,0.000000}%
\pgfsetfillcolor{currentfill}%
\pgfsetlinewidth{0.501875pt}%
\definecolor{currentstroke}{rgb}{0.000000,0.000000,0.000000}%
\pgfsetstrokecolor{currentstroke}%
\pgfsetdash{}{0pt}%
\pgfsys@defobject{currentmarker}{\pgfqpoint{0.000000in}{0.000000in}}{\pgfqpoint{0.020833in}{0.000000in}}{%
\pgfpathmoveto{\pgfqpoint{0.000000in}{0.000000in}}%
\pgfpathlineto{\pgfqpoint{0.020833in}{0.000000in}}%
\pgfusepath{stroke,fill}%
}%
\begin{pgfscope}%
\pgfsys@transformshift{0.448229in}{0.516085in}%
\pgfsys@useobject{currentmarker}{}%
\end{pgfscope}%
\end{pgfscope}%
\begin{pgfscope}%
\pgfsetbuttcap%
\pgfsetroundjoin%
\definecolor{currentfill}{rgb}{0.000000,0.000000,0.000000}%
\pgfsetfillcolor{currentfill}%
\pgfsetlinewidth{0.501875pt}%
\definecolor{currentstroke}{rgb}{0.000000,0.000000,0.000000}%
\pgfsetstrokecolor{currentstroke}%
\pgfsetdash{}{0pt}%
\pgfsys@defobject{currentmarker}{\pgfqpoint{-0.020833in}{0.000000in}}{\pgfqpoint{-0.000000in}{0.000000in}}{%
\pgfpathmoveto{\pgfqpoint{-0.000000in}{0.000000in}}%
\pgfpathlineto{\pgfqpoint{-0.020833in}{0.000000in}}%
\pgfusepath{stroke,fill}%
}%
\begin{pgfscope}%
\pgfsys@transformshift{6.155857in}{0.516085in}%
\pgfsys@useobject{currentmarker}{}%
\end{pgfscope}%
\end{pgfscope}%
\begin{pgfscope}%
\pgfsetbuttcap%
\pgfsetroundjoin%
\definecolor{currentfill}{rgb}{0.000000,0.000000,0.000000}%
\pgfsetfillcolor{currentfill}%
\pgfsetlinewidth{0.501875pt}%
\definecolor{currentstroke}{rgb}{0.000000,0.000000,0.000000}%
\pgfsetstrokecolor{currentstroke}%
\pgfsetdash{}{0pt}%
\pgfsys@defobject{currentmarker}{\pgfqpoint{0.000000in}{0.000000in}}{\pgfqpoint{0.020833in}{0.000000in}}{%
\pgfpathmoveto{\pgfqpoint{0.000000in}{0.000000in}}%
\pgfpathlineto{\pgfqpoint{0.020833in}{0.000000in}}%
\pgfusepath{stroke,fill}%
}%
\begin{pgfscope}%
\pgfsys@transformshift{0.448229in}{0.755780in}%
\pgfsys@useobject{currentmarker}{}%
\end{pgfscope}%
\end{pgfscope}%
\begin{pgfscope}%
\pgfsetbuttcap%
\pgfsetroundjoin%
\definecolor{currentfill}{rgb}{0.000000,0.000000,0.000000}%
\pgfsetfillcolor{currentfill}%
\pgfsetlinewidth{0.501875pt}%
\definecolor{currentstroke}{rgb}{0.000000,0.000000,0.000000}%
\pgfsetstrokecolor{currentstroke}%
\pgfsetdash{}{0pt}%
\pgfsys@defobject{currentmarker}{\pgfqpoint{-0.020833in}{0.000000in}}{\pgfqpoint{-0.000000in}{0.000000in}}{%
\pgfpathmoveto{\pgfqpoint{-0.000000in}{0.000000in}}%
\pgfpathlineto{\pgfqpoint{-0.020833in}{0.000000in}}%
\pgfusepath{stroke,fill}%
}%
\begin{pgfscope}%
\pgfsys@transformshift{6.155857in}{0.755780in}%
\pgfsys@useobject{currentmarker}{}%
\end{pgfscope}%
\end{pgfscope}%
\begin{pgfscope}%
\pgfsetbuttcap%
\pgfsetroundjoin%
\definecolor{currentfill}{rgb}{0.000000,0.000000,0.000000}%
\pgfsetfillcolor{currentfill}%
\pgfsetlinewidth{0.501875pt}%
\definecolor{currentstroke}{rgb}{0.000000,0.000000,0.000000}%
\pgfsetstrokecolor{currentstroke}%
\pgfsetdash{}{0pt}%
\pgfsys@defobject{currentmarker}{\pgfqpoint{0.000000in}{0.000000in}}{\pgfqpoint{0.020833in}{0.000000in}}{%
\pgfpathmoveto{\pgfqpoint{0.000000in}{0.000000in}}%
\pgfpathlineto{\pgfqpoint{0.020833in}{0.000000in}}%
\pgfusepath{stroke,fill}%
}%
\begin{pgfscope}%
\pgfsys@transformshift{0.448229in}{0.875627in}%
\pgfsys@useobject{currentmarker}{}%
\end{pgfscope}%
\end{pgfscope}%
\begin{pgfscope}%
\pgfsetbuttcap%
\pgfsetroundjoin%
\definecolor{currentfill}{rgb}{0.000000,0.000000,0.000000}%
\pgfsetfillcolor{currentfill}%
\pgfsetlinewidth{0.501875pt}%
\definecolor{currentstroke}{rgb}{0.000000,0.000000,0.000000}%
\pgfsetstrokecolor{currentstroke}%
\pgfsetdash{}{0pt}%
\pgfsys@defobject{currentmarker}{\pgfqpoint{-0.020833in}{0.000000in}}{\pgfqpoint{-0.000000in}{0.000000in}}{%
\pgfpathmoveto{\pgfqpoint{-0.000000in}{0.000000in}}%
\pgfpathlineto{\pgfqpoint{-0.020833in}{0.000000in}}%
\pgfusepath{stroke,fill}%
}%
\begin{pgfscope}%
\pgfsys@transformshift{6.155857in}{0.875627in}%
\pgfsys@useobject{currentmarker}{}%
\end{pgfscope}%
\end{pgfscope}%
\begin{pgfscope}%
\pgfsetbuttcap%
\pgfsetroundjoin%
\definecolor{currentfill}{rgb}{0.000000,0.000000,0.000000}%
\pgfsetfillcolor{currentfill}%
\pgfsetlinewidth{0.501875pt}%
\definecolor{currentstroke}{rgb}{0.000000,0.000000,0.000000}%
\pgfsetstrokecolor{currentstroke}%
\pgfsetdash{}{0pt}%
\pgfsys@defobject{currentmarker}{\pgfqpoint{0.000000in}{0.000000in}}{\pgfqpoint{0.020833in}{0.000000in}}{%
\pgfpathmoveto{\pgfqpoint{0.000000in}{0.000000in}}%
\pgfpathlineto{\pgfqpoint{0.020833in}{0.000000in}}%
\pgfusepath{stroke,fill}%
}%
\begin{pgfscope}%
\pgfsys@transformshift{0.448229in}{0.995474in}%
\pgfsys@useobject{currentmarker}{}%
\end{pgfscope}%
\end{pgfscope}%
\begin{pgfscope}%
\pgfsetbuttcap%
\pgfsetroundjoin%
\definecolor{currentfill}{rgb}{0.000000,0.000000,0.000000}%
\pgfsetfillcolor{currentfill}%
\pgfsetlinewidth{0.501875pt}%
\definecolor{currentstroke}{rgb}{0.000000,0.000000,0.000000}%
\pgfsetstrokecolor{currentstroke}%
\pgfsetdash{}{0pt}%
\pgfsys@defobject{currentmarker}{\pgfqpoint{-0.020833in}{0.000000in}}{\pgfqpoint{-0.000000in}{0.000000in}}{%
\pgfpathmoveto{\pgfqpoint{-0.000000in}{0.000000in}}%
\pgfpathlineto{\pgfqpoint{-0.020833in}{0.000000in}}%
\pgfusepath{stroke,fill}%
}%
\begin{pgfscope}%
\pgfsys@transformshift{6.155857in}{0.995474in}%
\pgfsys@useobject{currentmarker}{}%
\end{pgfscope}%
\end{pgfscope}%
\begin{pgfscope}%
\pgfsetbuttcap%
\pgfsetroundjoin%
\definecolor{currentfill}{rgb}{0.000000,0.000000,0.000000}%
\pgfsetfillcolor{currentfill}%
\pgfsetlinewidth{0.501875pt}%
\definecolor{currentstroke}{rgb}{0.000000,0.000000,0.000000}%
\pgfsetstrokecolor{currentstroke}%
\pgfsetdash{}{0pt}%
\pgfsys@defobject{currentmarker}{\pgfqpoint{0.000000in}{0.000000in}}{\pgfqpoint{0.020833in}{0.000000in}}{%
\pgfpathmoveto{\pgfqpoint{0.000000in}{0.000000in}}%
\pgfpathlineto{\pgfqpoint{0.020833in}{0.000000in}}%
\pgfusepath{stroke,fill}%
}%
\begin{pgfscope}%
\pgfsys@transformshift{0.448229in}{1.115322in}%
\pgfsys@useobject{currentmarker}{}%
\end{pgfscope}%
\end{pgfscope}%
\begin{pgfscope}%
\pgfsetbuttcap%
\pgfsetroundjoin%
\definecolor{currentfill}{rgb}{0.000000,0.000000,0.000000}%
\pgfsetfillcolor{currentfill}%
\pgfsetlinewidth{0.501875pt}%
\definecolor{currentstroke}{rgb}{0.000000,0.000000,0.000000}%
\pgfsetstrokecolor{currentstroke}%
\pgfsetdash{}{0pt}%
\pgfsys@defobject{currentmarker}{\pgfqpoint{-0.020833in}{0.000000in}}{\pgfqpoint{-0.000000in}{0.000000in}}{%
\pgfpathmoveto{\pgfqpoint{-0.000000in}{0.000000in}}%
\pgfpathlineto{\pgfqpoint{-0.020833in}{0.000000in}}%
\pgfusepath{stroke,fill}%
}%
\begin{pgfscope}%
\pgfsys@transformshift{6.155857in}{1.115322in}%
\pgfsys@useobject{currentmarker}{}%
\end{pgfscope}%
\end{pgfscope}%
\begin{pgfscope}%
\pgfsetbuttcap%
\pgfsetroundjoin%
\definecolor{currentfill}{rgb}{0.000000,0.000000,0.000000}%
\pgfsetfillcolor{currentfill}%
\pgfsetlinewidth{0.501875pt}%
\definecolor{currentstroke}{rgb}{0.000000,0.000000,0.000000}%
\pgfsetstrokecolor{currentstroke}%
\pgfsetdash{}{0pt}%
\pgfsys@defobject{currentmarker}{\pgfqpoint{0.000000in}{0.000000in}}{\pgfqpoint{0.020833in}{0.000000in}}{%
\pgfpathmoveto{\pgfqpoint{0.000000in}{0.000000in}}%
\pgfpathlineto{\pgfqpoint{0.020833in}{0.000000in}}%
\pgfusepath{stroke,fill}%
}%
\begin{pgfscope}%
\pgfsys@transformshift{0.448229in}{1.355017in}%
\pgfsys@useobject{currentmarker}{}%
\end{pgfscope}%
\end{pgfscope}%
\begin{pgfscope}%
\pgfsetbuttcap%
\pgfsetroundjoin%
\definecolor{currentfill}{rgb}{0.000000,0.000000,0.000000}%
\pgfsetfillcolor{currentfill}%
\pgfsetlinewidth{0.501875pt}%
\definecolor{currentstroke}{rgb}{0.000000,0.000000,0.000000}%
\pgfsetstrokecolor{currentstroke}%
\pgfsetdash{}{0pt}%
\pgfsys@defobject{currentmarker}{\pgfqpoint{-0.020833in}{0.000000in}}{\pgfqpoint{-0.000000in}{0.000000in}}{%
\pgfpathmoveto{\pgfqpoint{-0.000000in}{0.000000in}}%
\pgfpathlineto{\pgfqpoint{-0.020833in}{0.000000in}}%
\pgfusepath{stroke,fill}%
}%
\begin{pgfscope}%
\pgfsys@transformshift{6.155857in}{1.355017in}%
\pgfsys@useobject{currentmarker}{}%
\end{pgfscope}%
\end{pgfscope}%
\begin{pgfscope}%
\pgfsetbuttcap%
\pgfsetroundjoin%
\definecolor{currentfill}{rgb}{0.000000,0.000000,0.000000}%
\pgfsetfillcolor{currentfill}%
\pgfsetlinewidth{0.501875pt}%
\definecolor{currentstroke}{rgb}{0.000000,0.000000,0.000000}%
\pgfsetstrokecolor{currentstroke}%
\pgfsetdash{}{0pt}%
\pgfsys@defobject{currentmarker}{\pgfqpoint{0.000000in}{0.000000in}}{\pgfqpoint{0.020833in}{0.000000in}}{%
\pgfpathmoveto{\pgfqpoint{0.000000in}{0.000000in}}%
\pgfpathlineto{\pgfqpoint{0.020833in}{0.000000in}}%
\pgfusepath{stroke,fill}%
}%
\begin{pgfscope}%
\pgfsys@transformshift{0.448229in}{1.474864in}%
\pgfsys@useobject{currentmarker}{}%
\end{pgfscope}%
\end{pgfscope}%
\begin{pgfscope}%
\pgfsetbuttcap%
\pgfsetroundjoin%
\definecolor{currentfill}{rgb}{0.000000,0.000000,0.000000}%
\pgfsetfillcolor{currentfill}%
\pgfsetlinewidth{0.501875pt}%
\definecolor{currentstroke}{rgb}{0.000000,0.000000,0.000000}%
\pgfsetstrokecolor{currentstroke}%
\pgfsetdash{}{0pt}%
\pgfsys@defobject{currentmarker}{\pgfqpoint{-0.020833in}{0.000000in}}{\pgfqpoint{-0.000000in}{0.000000in}}{%
\pgfpathmoveto{\pgfqpoint{-0.000000in}{0.000000in}}%
\pgfpathlineto{\pgfqpoint{-0.020833in}{0.000000in}}%
\pgfusepath{stroke,fill}%
}%
\begin{pgfscope}%
\pgfsys@transformshift{6.155857in}{1.474864in}%
\pgfsys@useobject{currentmarker}{}%
\end{pgfscope}%
\end{pgfscope}%
\begin{pgfscope}%
\pgfsetbuttcap%
\pgfsetroundjoin%
\definecolor{currentfill}{rgb}{0.000000,0.000000,0.000000}%
\pgfsetfillcolor{currentfill}%
\pgfsetlinewidth{0.501875pt}%
\definecolor{currentstroke}{rgb}{0.000000,0.000000,0.000000}%
\pgfsetstrokecolor{currentstroke}%
\pgfsetdash{}{0pt}%
\pgfsys@defobject{currentmarker}{\pgfqpoint{0.000000in}{0.000000in}}{\pgfqpoint{0.020833in}{0.000000in}}{%
\pgfpathmoveto{\pgfqpoint{0.000000in}{0.000000in}}%
\pgfpathlineto{\pgfqpoint{0.020833in}{0.000000in}}%
\pgfusepath{stroke,fill}%
}%
\begin{pgfscope}%
\pgfsys@transformshift{0.448229in}{1.594712in}%
\pgfsys@useobject{currentmarker}{}%
\end{pgfscope}%
\end{pgfscope}%
\begin{pgfscope}%
\pgfsetbuttcap%
\pgfsetroundjoin%
\definecolor{currentfill}{rgb}{0.000000,0.000000,0.000000}%
\pgfsetfillcolor{currentfill}%
\pgfsetlinewidth{0.501875pt}%
\definecolor{currentstroke}{rgb}{0.000000,0.000000,0.000000}%
\pgfsetstrokecolor{currentstroke}%
\pgfsetdash{}{0pt}%
\pgfsys@defobject{currentmarker}{\pgfqpoint{-0.020833in}{0.000000in}}{\pgfqpoint{-0.000000in}{0.000000in}}{%
\pgfpathmoveto{\pgfqpoint{-0.000000in}{0.000000in}}%
\pgfpathlineto{\pgfqpoint{-0.020833in}{0.000000in}}%
\pgfusepath{stroke,fill}%
}%
\begin{pgfscope}%
\pgfsys@transformshift{6.155857in}{1.594712in}%
\pgfsys@useobject{currentmarker}{}%
\end{pgfscope}%
\end{pgfscope}%
\begin{pgfscope}%
\pgfsetbuttcap%
\pgfsetroundjoin%
\definecolor{currentfill}{rgb}{0.000000,0.000000,0.000000}%
\pgfsetfillcolor{currentfill}%
\pgfsetlinewidth{0.501875pt}%
\definecolor{currentstroke}{rgb}{0.000000,0.000000,0.000000}%
\pgfsetstrokecolor{currentstroke}%
\pgfsetdash{}{0pt}%
\pgfsys@defobject{currentmarker}{\pgfqpoint{0.000000in}{0.000000in}}{\pgfqpoint{0.020833in}{0.000000in}}{%
\pgfpathmoveto{\pgfqpoint{0.000000in}{0.000000in}}%
\pgfpathlineto{\pgfqpoint{0.020833in}{0.000000in}}%
\pgfusepath{stroke,fill}%
}%
\begin{pgfscope}%
\pgfsys@transformshift{0.448229in}{1.714559in}%
\pgfsys@useobject{currentmarker}{}%
\end{pgfscope}%
\end{pgfscope}%
\begin{pgfscope}%
\pgfsetbuttcap%
\pgfsetroundjoin%
\definecolor{currentfill}{rgb}{0.000000,0.000000,0.000000}%
\pgfsetfillcolor{currentfill}%
\pgfsetlinewidth{0.501875pt}%
\definecolor{currentstroke}{rgb}{0.000000,0.000000,0.000000}%
\pgfsetstrokecolor{currentstroke}%
\pgfsetdash{}{0pt}%
\pgfsys@defobject{currentmarker}{\pgfqpoint{-0.020833in}{0.000000in}}{\pgfqpoint{-0.000000in}{0.000000in}}{%
\pgfpathmoveto{\pgfqpoint{-0.000000in}{0.000000in}}%
\pgfpathlineto{\pgfqpoint{-0.020833in}{0.000000in}}%
\pgfusepath{stroke,fill}%
}%
\begin{pgfscope}%
\pgfsys@transformshift{6.155857in}{1.714559in}%
\pgfsys@useobject{currentmarker}{}%
\end{pgfscope}%
\end{pgfscope}%
\begin{pgfscope}%
\pgfsetbuttcap%
\pgfsetroundjoin%
\definecolor{currentfill}{rgb}{0.000000,0.000000,0.000000}%
\pgfsetfillcolor{currentfill}%
\pgfsetlinewidth{0.501875pt}%
\definecolor{currentstroke}{rgb}{0.000000,0.000000,0.000000}%
\pgfsetstrokecolor{currentstroke}%
\pgfsetdash{}{0pt}%
\pgfsys@defobject{currentmarker}{\pgfqpoint{0.000000in}{0.000000in}}{\pgfqpoint{0.020833in}{0.000000in}}{%
\pgfpathmoveto{\pgfqpoint{0.000000in}{0.000000in}}%
\pgfpathlineto{\pgfqpoint{0.020833in}{0.000000in}}%
\pgfusepath{stroke,fill}%
}%
\begin{pgfscope}%
\pgfsys@transformshift{0.448229in}{1.954254in}%
\pgfsys@useobject{currentmarker}{}%
\end{pgfscope}%
\end{pgfscope}%
\begin{pgfscope}%
\pgfsetbuttcap%
\pgfsetroundjoin%
\definecolor{currentfill}{rgb}{0.000000,0.000000,0.000000}%
\pgfsetfillcolor{currentfill}%
\pgfsetlinewidth{0.501875pt}%
\definecolor{currentstroke}{rgb}{0.000000,0.000000,0.000000}%
\pgfsetstrokecolor{currentstroke}%
\pgfsetdash{}{0pt}%
\pgfsys@defobject{currentmarker}{\pgfqpoint{-0.020833in}{0.000000in}}{\pgfqpoint{-0.000000in}{0.000000in}}{%
\pgfpathmoveto{\pgfqpoint{-0.000000in}{0.000000in}}%
\pgfpathlineto{\pgfqpoint{-0.020833in}{0.000000in}}%
\pgfusepath{stroke,fill}%
}%
\begin{pgfscope}%
\pgfsys@transformshift{6.155857in}{1.954254in}%
\pgfsys@useobject{currentmarker}{}%
\end{pgfscope}%
\end{pgfscope}%
\begin{pgfscope}%
\pgfsetbuttcap%
\pgfsetroundjoin%
\definecolor{currentfill}{rgb}{0.000000,0.000000,0.000000}%
\pgfsetfillcolor{currentfill}%
\pgfsetlinewidth{0.501875pt}%
\definecolor{currentstroke}{rgb}{0.000000,0.000000,0.000000}%
\pgfsetstrokecolor{currentstroke}%
\pgfsetdash{}{0pt}%
\pgfsys@defobject{currentmarker}{\pgfqpoint{0.000000in}{0.000000in}}{\pgfqpoint{0.020833in}{0.000000in}}{%
\pgfpathmoveto{\pgfqpoint{0.000000in}{0.000000in}}%
\pgfpathlineto{\pgfqpoint{0.020833in}{0.000000in}}%
\pgfusepath{stroke,fill}%
}%
\begin{pgfscope}%
\pgfsys@transformshift{0.448229in}{2.074101in}%
\pgfsys@useobject{currentmarker}{}%
\end{pgfscope}%
\end{pgfscope}%
\begin{pgfscope}%
\pgfsetbuttcap%
\pgfsetroundjoin%
\definecolor{currentfill}{rgb}{0.000000,0.000000,0.000000}%
\pgfsetfillcolor{currentfill}%
\pgfsetlinewidth{0.501875pt}%
\definecolor{currentstroke}{rgb}{0.000000,0.000000,0.000000}%
\pgfsetstrokecolor{currentstroke}%
\pgfsetdash{}{0pt}%
\pgfsys@defobject{currentmarker}{\pgfqpoint{-0.020833in}{0.000000in}}{\pgfqpoint{-0.000000in}{0.000000in}}{%
\pgfpathmoveto{\pgfqpoint{-0.000000in}{0.000000in}}%
\pgfpathlineto{\pgfqpoint{-0.020833in}{0.000000in}}%
\pgfusepath{stroke,fill}%
}%
\begin{pgfscope}%
\pgfsys@transformshift{6.155857in}{2.074101in}%
\pgfsys@useobject{currentmarker}{}%
\end{pgfscope}%
\end{pgfscope}%
\begin{pgfscope}%
\pgfsetbuttcap%
\pgfsetroundjoin%
\definecolor{currentfill}{rgb}{0.000000,0.000000,0.000000}%
\pgfsetfillcolor{currentfill}%
\pgfsetlinewidth{0.501875pt}%
\definecolor{currentstroke}{rgb}{0.000000,0.000000,0.000000}%
\pgfsetstrokecolor{currentstroke}%
\pgfsetdash{}{0pt}%
\pgfsys@defobject{currentmarker}{\pgfqpoint{0.000000in}{0.000000in}}{\pgfqpoint{0.020833in}{0.000000in}}{%
\pgfpathmoveto{\pgfqpoint{0.000000in}{0.000000in}}%
\pgfpathlineto{\pgfqpoint{0.020833in}{0.000000in}}%
\pgfusepath{stroke,fill}%
}%
\begin{pgfscope}%
\pgfsys@transformshift{0.448229in}{2.193949in}%
\pgfsys@useobject{currentmarker}{}%
\end{pgfscope}%
\end{pgfscope}%
\begin{pgfscope}%
\pgfsetbuttcap%
\pgfsetroundjoin%
\definecolor{currentfill}{rgb}{0.000000,0.000000,0.000000}%
\pgfsetfillcolor{currentfill}%
\pgfsetlinewidth{0.501875pt}%
\definecolor{currentstroke}{rgb}{0.000000,0.000000,0.000000}%
\pgfsetstrokecolor{currentstroke}%
\pgfsetdash{}{0pt}%
\pgfsys@defobject{currentmarker}{\pgfqpoint{-0.020833in}{0.000000in}}{\pgfqpoint{-0.000000in}{0.000000in}}{%
\pgfpathmoveto{\pgfqpoint{-0.000000in}{0.000000in}}%
\pgfpathlineto{\pgfqpoint{-0.020833in}{0.000000in}}%
\pgfusepath{stroke,fill}%
}%
\begin{pgfscope}%
\pgfsys@transformshift{6.155857in}{2.193949in}%
\pgfsys@useobject{currentmarker}{}%
\end{pgfscope}%
\end{pgfscope}%
\begin{pgfscope}%
\pgfsetbuttcap%
\pgfsetroundjoin%
\definecolor{currentfill}{rgb}{0.000000,0.000000,0.000000}%
\pgfsetfillcolor{currentfill}%
\pgfsetlinewidth{0.501875pt}%
\definecolor{currentstroke}{rgb}{0.000000,0.000000,0.000000}%
\pgfsetstrokecolor{currentstroke}%
\pgfsetdash{}{0pt}%
\pgfsys@defobject{currentmarker}{\pgfqpoint{0.000000in}{0.000000in}}{\pgfqpoint{0.020833in}{0.000000in}}{%
\pgfpathmoveto{\pgfqpoint{0.000000in}{0.000000in}}%
\pgfpathlineto{\pgfqpoint{0.020833in}{0.000000in}}%
\pgfusepath{stroke,fill}%
}%
\begin{pgfscope}%
\pgfsys@transformshift{0.448229in}{2.313796in}%
\pgfsys@useobject{currentmarker}{}%
\end{pgfscope}%
\end{pgfscope}%
\begin{pgfscope}%
\pgfsetbuttcap%
\pgfsetroundjoin%
\definecolor{currentfill}{rgb}{0.000000,0.000000,0.000000}%
\pgfsetfillcolor{currentfill}%
\pgfsetlinewidth{0.501875pt}%
\definecolor{currentstroke}{rgb}{0.000000,0.000000,0.000000}%
\pgfsetstrokecolor{currentstroke}%
\pgfsetdash{}{0pt}%
\pgfsys@defobject{currentmarker}{\pgfqpoint{-0.020833in}{0.000000in}}{\pgfqpoint{-0.000000in}{0.000000in}}{%
\pgfpathmoveto{\pgfqpoint{-0.000000in}{0.000000in}}%
\pgfpathlineto{\pgfqpoint{-0.020833in}{0.000000in}}%
\pgfusepath{stroke,fill}%
}%
\begin{pgfscope}%
\pgfsys@transformshift{6.155857in}{2.313796in}%
\pgfsys@useobject{currentmarker}{}%
\end{pgfscope}%
\end{pgfscope}%
\begin{pgfscope}%
\pgfsetbuttcap%
\pgfsetroundjoin%
\definecolor{currentfill}{rgb}{0.000000,0.000000,0.000000}%
\pgfsetfillcolor{currentfill}%
\pgfsetlinewidth{0.501875pt}%
\definecolor{currentstroke}{rgb}{0.000000,0.000000,0.000000}%
\pgfsetstrokecolor{currentstroke}%
\pgfsetdash{}{0pt}%
\pgfsys@defobject{currentmarker}{\pgfqpoint{0.000000in}{0.000000in}}{\pgfqpoint{0.020833in}{0.000000in}}{%
\pgfpathmoveto{\pgfqpoint{0.000000in}{0.000000in}}%
\pgfpathlineto{\pgfqpoint{0.020833in}{0.000000in}}%
\pgfusepath{stroke,fill}%
}%
\begin{pgfscope}%
\pgfsys@transformshift{0.448229in}{2.553491in}%
\pgfsys@useobject{currentmarker}{}%
\end{pgfscope}%
\end{pgfscope}%
\begin{pgfscope}%
\pgfsetbuttcap%
\pgfsetroundjoin%
\definecolor{currentfill}{rgb}{0.000000,0.000000,0.000000}%
\pgfsetfillcolor{currentfill}%
\pgfsetlinewidth{0.501875pt}%
\definecolor{currentstroke}{rgb}{0.000000,0.000000,0.000000}%
\pgfsetstrokecolor{currentstroke}%
\pgfsetdash{}{0pt}%
\pgfsys@defobject{currentmarker}{\pgfqpoint{-0.020833in}{0.000000in}}{\pgfqpoint{-0.000000in}{0.000000in}}{%
\pgfpathmoveto{\pgfqpoint{-0.000000in}{0.000000in}}%
\pgfpathlineto{\pgfqpoint{-0.020833in}{0.000000in}}%
\pgfusepath{stroke,fill}%
}%
\begin{pgfscope}%
\pgfsys@transformshift{6.155857in}{2.553491in}%
\pgfsys@useobject{currentmarker}{}%
\end{pgfscope}%
\end{pgfscope}%
\begin{pgfscope}%
\pgfsetbuttcap%
\pgfsetroundjoin%
\definecolor{currentfill}{rgb}{0.000000,0.000000,0.000000}%
\pgfsetfillcolor{currentfill}%
\pgfsetlinewidth{0.501875pt}%
\definecolor{currentstroke}{rgb}{0.000000,0.000000,0.000000}%
\pgfsetstrokecolor{currentstroke}%
\pgfsetdash{}{0pt}%
\pgfsys@defobject{currentmarker}{\pgfqpoint{0.000000in}{0.000000in}}{\pgfqpoint{0.020833in}{0.000000in}}{%
\pgfpathmoveto{\pgfqpoint{0.000000in}{0.000000in}}%
\pgfpathlineto{\pgfqpoint{0.020833in}{0.000000in}}%
\pgfusepath{stroke,fill}%
}%
\begin{pgfscope}%
\pgfsys@transformshift{0.448229in}{2.673339in}%
\pgfsys@useobject{currentmarker}{}%
\end{pgfscope}%
\end{pgfscope}%
\begin{pgfscope}%
\pgfsetbuttcap%
\pgfsetroundjoin%
\definecolor{currentfill}{rgb}{0.000000,0.000000,0.000000}%
\pgfsetfillcolor{currentfill}%
\pgfsetlinewidth{0.501875pt}%
\definecolor{currentstroke}{rgb}{0.000000,0.000000,0.000000}%
\pgfsetstrokecolor{currentstroke}%
\pgfsetdash{}{0pt}%
\pgfsys@defobject{currentmarker}{\pgfqpoint{-0.020833in}{0.000000in}}{\pgfqpoint{-0.000000in}{0.000000in}}{%
\pgfpathmoveto{\pgfqpoint{-0.000000in}{0.000000in}}%
\pgfpathlineto{\pgfqpoint{-0.020833in}{0.000000in}}%
\pgfusepath{stroke,fill}%
}%
\begin{pgfscope}%
\pgfsys@transformshift{6.155857in}{2.673339in}%
\pgfsys@useobject{currentmarker}{}%
\end{pgfscope}%
\end{pgfscope}%
\begin{pgfscope}%
\pgfsetbuttcap%
\pgfsetroundjoin%
\definecolor{currentfill}{rgb}{0.000000,0.000000,0.000000}%
\pgfsetfillcolor{currentfill}%
\pgfsetlinewidth{0.501875pt}%
\definecolor{currentstroke}{rgb}{0.000000,0.000000,0.000000}%
\pgfsetstrokecolor{currentstroke}%
\pgfsetdash{}{0pt}%
\pgfsys@defobject{currentmarker}{\pgfqpoint{0.000000in}{0.000000in}}{\pgfqpoint{0.020833in}{0.000000in}}{%
\pgfpathmoveto{\pgfqpoint{0.000000in}{0.000000in}}%
\pgfpathlineto{\pgfqpoint{0.020833in}{0.000000in}}%
\pgfusepath{stroke,fill}%
}%
\begin{pgfscope}%
\pgfsys@transformshift{0.448229in}{2.793186in}%
\pgfsys@useobject{currentmarker}{}%
\end{pgfscope}%
\end{pgfscope}%
\begin{pgfscope}%
\pgfsetbuttcap%
\pgfsetroundjoin%
\definecolor{currentfill}{rgb}{0.000000,0.000000,0.000000}%
\pgfsetfillcolor{currentfill}%
\pgfsetlinewidth{0.501875pt}%
\definecolor{currentstroke}{rgb}{0.000000,0.000000,0.000000}%
\pgfsetstrokecolor{currentstroke}%
\pgfsetdash{}{0pt}%
\pgfsys@defobject{currentmarker}{\pgfqpoint{-0.020833in}{0.000000in}}{\pgfqpoint{-0.000000in}{0.000000in}}{%
\pgfpathmoveto{\pgfqpoint{-0.000000in}{0.000000in}}%
\pgfpathlineto{\pgfqpoint{-0.020833in}{0.000000in}}%
\pgfusepath{stroke,fill}%
}%
\begin{pgfscope}%
\pgfsys@transformshift{6.155857in}{2.793186in}%
\pgfsys@useobject{currentmarker}{}%
\end{pgfscope}%
\end{pgfscope}%
\begin{pgfscope}%
\pgfsetbuttcap%
\pgfsetroundjoin%
\definecolor{currentfill}{rgb}{0.000000,0.000000,0.000000}%
\pgfsetfillcolor{currentfill}%
\pgfsetlinewidth{0.501875pt}%
\definecolor{currentstroke}{rgb}{0.000000,0.000000,0.000000}%
\pgfsetstrokecolor{currentstroke}%
\pgfsetdash{}{0pt}%
\pgfsys@defobject{currentmarker}{\pgfqpoint{0.000000in}{0.000000in}}{\pgfqpoint{0.020833in}{0.000000in}}{%
\pgfpathmoveto{\pgfqpoint{0.000000in}{0.000000in}}%
\pgfpathlineto{\pgfqpoint{0.020833in}{0.000000in}}%
\pgfusepath{stroke,fill}%
}%
\begin{pgfscope}%
\pgfsys@transformshift{0.448229in}{2.913033in}%
\pgfsys@useobject{currentmarker}{}%
\end{pgfscope}%
\end{pgfscope}%
\begin{pgfscope}%
\pgfsetbuttcap%
\pgfsetroundjoin%
\definecolor{currentfill}{rgb}{0.000000,0.000000,0.000000}%
\pgfsetfillcolor{currentfill}%
\pgfsetlinewidth{0.501875pt}%
\definecolor{currentstroke}{rgb}{0.000000,0.000000,0.000000}%
\pgfsetstrokecolor{currentstroke}%
\pgfsetdash{}{0pt}%
\pgfsys@defobject{currentmarker}{\pgfqpoint{-0.020833in}{0.000000in}}{\pgfqpoint{-0.000000in}{0.000000in}}{%
\pgfpathmoveto{\pgfqpoint{-0.000000in}{0.000000in}}%
\pgfpathlineto{\pgfqpoint{-0.020833in}{0.000000in}}%
\pgfusepath{stroke,fill}%
}%
\begin{pgfscope}%
\pgfsys@transformshift{6.155857in}{2.913033in}%
\pgfsys@useobject{currentmarker}{}%
\end{pgfscope}%
\end{pgfscope}%
\begin{pgfscope}%
\pgfsetbuttcap%
\pgfsetroundjoin%
\definecolor{currentfill}{rgb}{0.000000,0.000000,0.000000}%
\pgfsetfillcolor{currentfill}%
\pgfsetlinewidth{0.501875pt}%
\definecolor{currentstroke}{rgb}{0.000000,0.000000,0.000000}%
\pgfsetstrokecolor{currentstroke}%
\pgfsetdash{}{0pt}%
\pgfsys@defobject{currentmarker}{\pgfqpoint{0.000000in}{0.000000in}}{\pgfqpoint{0.020833in}{0.000000in}}{%
\pgfpathmoveto{\pgfqpoint{0.000000in}{0.000000in}}%
\pgfpathlineto{\pgfqpoint{0.020833in}{0.000000in}}%
\pgfusepath{stroke,fill}%
}%
\begin{pgfscope}%
\pgfsys@transformshift{0.448229in}{3.152728in}%
\pgfsys@useobject{currentmarker}{}%
\end{pgfscope}%
\end{pgfscope}%
\begin{pgfscope}%
\pgfsetbuttcap%
\pgfsetroundjoin%
\definecolor{currentfill}{rgb}{0.000000,0.000000,0.000000}%
\pgfsetfillcolor{currentfill}%
\pgfsetlinewidth{0.501875pt}%
\definecolor{currentstroke}{rgb}{0.000000,0.000000,0.000000}%
\pgfsetstrokecolor{currentstroke}%
\pgfsetdash{}{0pt}%
\pgfsys@defobject{currentmarker}{\pgfqpoint{-0.020833in}{0.000000in}}{\pgfqpoint{-0.000000in}{0.000000in}}{%
\pgfpathmoveto{\pgfqpoint{-0.000000in}{0.000000in}}%
\pgfpathlineto{\pgfqpoint{-0.020833in}{0.000000in}}%
\pgfusepath{stroke,fill}%
}%
\begin{pgfscope}%
\pgfsys@transformshift{6.155857in}{3.152728in}%
\pgfsys@useobject{currentmarker}{}%
\end{pgfscope}%
\end{pgfscope}%
\begin{pgfscope}%
\pgfsetbuttcap%
\pgfsetroundjoin%
\definecolor{currentfill}{rgb}{0.000000,0.000000,0.000000}%
\pgfsetfillcolor{currentfill}%
\pgfsetlinewidth{0.501875pt}%
\definecolor{currentstroke}{rgb}{0.000000,0.000000,0.000000}%
\pgfsetstrokecolor{currentstroke}%
\pgfsetdash{}{0pt}%
\pgfsys@defobject{currentmarker}{\pgfqpoint{0.000000in}{0.000000in}}{\pgfqpoint{0.020833in}{0.000000in}}{%
\pgfpathmoveto{\pgfqpoint{0.000000in}{0.000000in}}%
\pgfpathlineto{\pgfqpoint{0.020833in}{0.000000in}}%
\pgfusepath{stroke,fill}%
}%
\begin{pgfscope}%
\pgfsys@transformshift{0.448229in}{3.272576in}%
\pgfsys@useobject{currentmarker}{}%
\end{pgfscope}%
\end{pgfscope}%
\begin{pgfscope}%
\pgfsetbuttcap%
\pgfsetroundjoin%
\definecolor{currentfill}{rgb}{0.000000,0.000000,0.000000}%
\pgfsetfillcolor{currentfill}%
\pgfsetlinewidth{0.501875pt}%
\definecolor{currentstroke}{rgb}{0.000000,0.000000,0.000000}%
\pgfsetstrokecolor{currentstroke}%
\pgfsetdash{}{0pt}%
\pgfsys@defobject{currentmarker}{\pgfqpoint{-0.020833in}{0.000000in}}{\pgfqpoint{-0.000000in}{0.000000in}}{%
\pgfpathmoveto{\pgfqpoint{-0.000000in}{0.000000in}}%
\pgfpathlineto{\pgfqpoint{-0.020833in}{0.000000in}}%
\pgfusepath{stroke,fill}%
}%
\begin{pgfscope}%
\pgfsys@transformshift{6.155857in}{3.272576in}%
\pgfsys@useobject{currentmarker}{}%
\end{pgfscope}%
\end{pgfscope}%
\begin{pgfscope}%
\pgfsetbuttcap%
\pgfsetroundjoin%
\definecolor{currentfill}{rgb}{0.000000,0.000000,0.000000}%
\pgfsetfillcolor{currentfill}%
\pgfsetlinewidth{0.501875pt}%
\definecolor{currentstroke}{rgb}{0.000000,0.000000,0.000000}%
\pgfsetstrokecolor{currentstroke}%
\pgfsetdash{}{0pt}%
\pgfsys@defobject{currentmarker}{\pgfqpoint{0.000000in}{0.000000in}}{\pgfqpoint{0.020833in}{0.000000in}}{%
\pgfpathmoveto{\pgfqpoint{0.000000in}{0.000000in}}%
\pgfpathlineto{\pgfqpoint{0.020833in}{0.000000in}}%
\pgfusepath{stroke,fill}%
}%
\begin{pgfscope}%
\pgfsys@transformshift{0.448229in}{3.392423in}%
\pgfsys@useobject{currentmarker}{}%
\end{pgfscope}%
\end{pgfscope}%
\begin{pgfscope}%
\pgfsetbuttcap%
\pgfsetroundjoin%
\definecolor{currentfill}{rgb}{0.000000,0.000000,0.000000}%
\pgfsetfillcolor{currentfill}%
\pgfsetlinewidth{0.501875pt}%
\definecolor{currentstroke}{rgb}{0.000000,0.000000,0.000000}%
\pgfsetstrokecolor{currentstroke}%
\pgfsetdash{}{0pt}%
\pgfsys@defobject{currentmarker}{\pgfqpoint{-0.020833in}{0.000000in}}{\pgfqpoint{-0.000000in}{0.000000in}}{%
\pgfpathmoveto{\pgfqpoint{-0.000000in}{0.000000in}}%
\pgfpathlineto{\pgfqpoint{-0.020833in}{0.000000in}}%
\pgfusepath{stroke,fill}%
}%
\begin{pgfscope}%
\pgfsys@transformshift{6.155857in}{3.392423in}%
\pgfsys@useobject{currentmarker}{}%
\end{pgfscope}%
\end{pgfscope}%
\begin{pgfscope}%
\pgfsetbuttcap%
\pgfsetroundjoin%
\definecolor{currentfill}{rgb}{0.000000,0.000000,0.000000}%
\pgfsetfillcolor{currentfill}%
\pgfsetlinewidth{0.501875pt}%
\definecolor{currentstroke}{rgb}{0.000000,0.000000,0.000000}%
\pgfsetstrokecolor{currentstroke}%
\pgfsetdash{}{0pt}%
\pgfsys@defobject{currentmarker}{\pgfqpoint{0.000000in}{0.000000in}}{\pgfqpoint{0.020833in}{0.000000in}}{%
\pgfpathmoveto{\pgfqpoint{0.000000in}{0.000000in}}%
\pgfpathlineto{\pgfqpoint{0.020833in}{0.000000in}}%
\pgfusepath{stroke,fill}%
}%
\begin{pgfscope}%
\pgfsys@transformshift{0.448229in}{3.512271in}%
\pgfsys@useobject{currentmarker}{}%
\end{pgfscope}%
\end{pgfscope}%
\begin{pgfscope}%
\pgfsetbuttcap%
\pgfsetroundjoin%
\definecolor{currentfill}{rgb}{0.000000,0.000000,0.000000}%
\pgfsetfillcolor{currentfill}%
\pgfsetlinewidth{0.501875pt}%
\definecolor{currentstroke}{rgb}{0.000000,0.000000,0.000000}%
\pgfsetstrokecolor{currentstroke}%
\pgfsetdash{}{0pt}%
\pgfsys@defobject{currentmarker}{\pgfqpoint{-0.020833in}{0.000000in}}{\pgfqpoint{-0.000000in}{0.000000in}}{%
\pgfpathmoveto{\pgfqpoint{-0.000000in}{0.000000in}}%
\pgfpathlineto{\pgfqpoint{-0.020833in}{0.000000in}}%
\pgfusepath{stroke,fill}%
}%
\begin{pgfscope}%
\pgfsys@transformshift{6.155857in}{3.512271in}%
\pgfsys@useobject{currentmarker}{}%
\end{pgfscope}%
\end{pgfscope}%
\begin{pgfscope}%
\pgfsetbuttcap%
\pgfsetroundjoin%
\definecolor{currentfill}{rgb}{0.000000,0.000000,0.000000}%
\pgfsetfillcolor{currentfill}%
\pgfsetlinewidth{0.501875pt}%
\definecolor{currentstroke}{rgb}{0.000000,0.000000,0.000000}%
\pgfsetstrokecolor{currentstroke}%
\pgfsetdash{}{0pt}%
\pgfsys@defobject{currentmarker}{\pgfqpoint{0.000000in}{0.000000in}}{\pgfqpoint{0.020833in}{0.000000in}}{%
\pgfpathmoveto{\pgfqpoint{0.000000in}{0.000000in}}%
\pgfpathlineto{\pgfqpoint{0.020833in}{0.000000in}}%
\pgfusepath{stroke,fill}%
}%
\begin{pgfscope}%
\pgfsys@transformshift{0.448229in}{3.751965in}%
\pgfsys@useobject{currentmarker}{}%
\end{pgfscope}%
\end{pgfscope}%
\begin{pgfscope}%
\pgfsetbuttcap%
\pgfsetroundjoin%
\definecolor{currentfill}{rgb}{0.000000,0.000000,0.000000}%
\pgfsetfillcolor{currentfill}%
\pgfsetlinewidth{0.501875pt}%
\definecolor{currentstroke}{rgb}{0.000000,0.000000,0.000000}%
\pgfsetstrokecolor{currentstroke}%
\pgfsetdash{}{0pt}%
\pgfsys@defobject{currentmarker}{\pgfqpoint{-0.020833in}{0.000000in}}{\pgfqpoint{-0.000000in}{0.000000in}}{%
\pgfpathmoveto{\pgfqpoint{-0.000000in}{0.000000in}}%
\pgfpathlineto{\pgfqpoint{-0.020833in}{0.000000in}}%
\pgfusepath{stroke,fill}%
}%
\begin{pgfscope}%
\pgfsys@transformshift{6.155857in}{3.751965in}%
\pgfsys@useobject{currentmarker}{}%
\end{pgfscope}%
\end{pgfscope}%
\begin{pgfscope}%
\pgfsetbuttcap%
\pgfsetroundjoin%
\definecolor{currentfill}{rgb}{0.000000,0.000000,0.000000}%
\pgfsetfillcolor{currentfill}%
\pgfsetlinewidth{0.501875pt}%
\definecolor{currentstroke}{rgb}{0.000000,0.000000,0.000000}%
\pgfsetstrokecolor{currentstroke}%
\pgfsetdash{}{0pt}%
\pgfsys@defobject{currentmarker}{\pgfqpoint{0.000000in}{0.000000in}}{\pgfqpoint{0.020833in}{0.000000in}}{%
\pgfpathmoveto{\pgfqpoint{0.000000in}{0.000000in}}%
\pgfpathlineto{\pgfqpoint{0.020833in}{0.000000in}}%
\pgfusepath{stroke,fill}%
}%
\begin{pgfscope}%
\pgfsys@transformshift{0.448229in}{3.871813in}%
\pgfsys@useobject{currentmarker}{}%
\end{pgfscope}%
\end{pgfscope}%
\begin{pgfscope}%
\pgfsetbuttcap%
\pgfsetroundjoin%
\definecolor{currentfill}{rgb}{0.000000,0.000000,0.000000}%
\pgfsetfillcolor{currentfill}%
\pgfsetlinewidth{0.501875pt}%
\definecolor{currentstroke}{rgb}{0.000000,0.000000,0.000000}%
\pgfsetstrokecolor{currentstroke}%
\pgfsetdash{}{0pt}%
\pgfsys@defobject{currentmarker}{\pgfqpoint{-0.020833in}{0.000000in}}{\pgfqpoint{-0.000000in}{0.000000in}}{%
\pgfpathmoveto{\pgfqpoint{-0.000000in}{0.000000in}}%
\pgfpathlineto{\pgfqpoint{-0.020833in}{0.000000in}}%
\pgfusepath{stroke,fill}%
}%
\begin{pgfscope}%
\pgfsys@transformshift{6.155857in}{3.871813in}%
\pgfsys@useobject{currentmarker}{}%
\end{pgfscope}%
\end{pgfscope}%
\begin{pgfscope}%
\definecolor{textcolor}{rgb}{0.000000,0.000000,0.000000}%
\pgfsetstrokecolor{textcolor}%
\pgfsetfillcolor{textcolor}%
\pgftext[x=0.179757in,y=2.201101in,,bottom,rotate=90.000000]{\color{textcolor}{\rmfamily\fontsize{10.000000}{12.000000}\selectfont\catcode`\^=\active\def^{\ifmmode\sp\else\^{}\fi}\catcode`\%=\active\def%{\%}$X_t$}}%
\end{pgfscope}%
\begin{pgfscope}%
\pgfpathrectangle{\pgfqpoint{0.448229in}{0.413680in}}{\pgfqpoint{5.707628in}{3.574842in}}%
\pgfusepath{clip}%
\pgfsetrectcap%
\pgfsetroundjoin%
\pgfsetlinewidth{1.003750pt}%
\definecolor{currentstroke}{rgb}{0.000000,0.000000,0.000000}%
\pgfsetstrokecolor{currentstroke}%
\pgfsetdash{}{0pt}%
\pgfpathmoveto{\pgfqpoint{0.448229in}{0.635932in}}%
\pgfpathlineto{\pgfqpoint{0.470524in}{0.641937in}}%
\pgfpathlineto{\pgfqpoint{0.492820in}{0.665092in}}%
\pgfpathlineto{\pgfqpoint{0.515115in}{0.631548in}}%
\pgfpathlineto{\pgfqpoint{0.537411in}{0.661890in}}%
\pgfpathlineto{\pgfqpoint{0.559706in}{0.650127in}}%
\pgfpathlineto{\pgfqpoint{0.582001in}{0.595492in}}%
\pgfpathlineto{\pgfqpoint{0.604297in}{0.576173in}}%
\pgfpathlineto{\pgfqpoint{0.626592in}{0.601646in}}%
\pgfpathlineto{\pgfqpoint{0.648888in}{0.617824in}}%
\pgfpathlineto{\pgfqpoint{0.671183in}{0.580147in}}%
\pgfpathlineto{\pgfqpoint{0.693479in}{0.606443in}}%
\pgfpathlineto{\pgfqpoint{0.715774in}{0.607224in}}%
\pgfpathlineto{\pgfqpoint{0.738069in}{0.626925in}}%
\pgfpathlineto{\pgfqpoint{0.760365in}{0.643455in}}%
\pgfpathlineto{\pgfqpoint{0.782660in}{0.604618in}}%
\pgfpathlineto{\pgfqpoint{0.804956in}{0.621531in}}%
\pgfpathlineto{\pgfqpoint{0.827251in}{0.603264in}}%
\pgfpathlineto{\pgfqpoint{0.849546in}{0.665197in}}%
\pgfpathlineto{\pgfqpoint{0.871842in}{0.697565in}}%
\pgfpathlineto{\pgfqpoint{0.894137in}{0.703704in}}%
\pgfpathlineto{\pgfqpoint{0.916433in}{0.669602in}}%
\pgfpathlineto{\pgfqpoint{0.938728in}{0.675925in}}%
\pgfpathlineto{\pgfqpoint{0.961024in}{0.686391in}}%
\pgfpathlineto{\pgfqpoint{0.983319in}{0.644204in}}%
\pgfpathlineto{\pgfqpoint{1.005614in}{0.735459in}}%
\pgfpathlineto{\pgfqpoint{1.027910in}{0.756051in}}%
\pgfpathlineto{\pgfqpoint{1.050205in}{0.832423in}}%
\pgfpathlineto{\pgfqpoint{1.072501in}{0.789359in}}%
\pgfpathlineto{\pgfqpoint{1.094796in}{0.865133in}}%
\pgfpathlineto{\pgfqpoint{1.117092in}{0.810371in}}%
\pgfpathlineto{\pgfqpoint{1.139387in}{0.787320in}}%
\pgfpathlineto{\pgfqpoint{1.161682in}{0.815098in}}%
\pgfpathlineto{\pgfqpoint{1.183978in}{0.821890in}}%
\pgfpathlineto{\pgfqpoint{1.206273in}{0.727943in}}%
\pgfpathlineto{\pgfqpoint{1.228569in}{0.688387in}}%
\pgfpathlineto{\pgfqpoint{1.250864in}{0.725968in}}%
\pgfpathlineto{\pgfqpoint{1.273159in}{0.817551in}}%
\pgfpathlineto{\pgfqpoint{1.295455in}{0.862438in}}%
\pgfpathlineto{\pgfqpoint{1.317750in}{0.855028in}}%
\pgfpathlineto{\pgfqpoint{1.340046in}{0.895731in}}%
\pgfpathlineto{\pgfqpoint{1.362341in}{0.849345in}}%
\pgfpathlineto{\pgfqpoint{1.384637in}{0.876936in}}%
\pgfpathlineto{\pgfqpoint{1.406932in}{0.822306in}}%
\pgfpathlineto{\pgfqpoint{1.429227in}{0.826641in}}%
\pgfpathlineto{\pgfqpoint{1.451523in}{0.825715in}}%
\pgfpathlineto{\pgfqpoint{1.473818in}{0.865826in}}%
\pgfpathlineto{\pgfqpoint{1.496114in}{0.857211in}}%
\pgfpathlineto{\pgfqpoint{1.518409in}{0.916601in}}%
\pgfpathlineto{\pgfqpoint{1.540705in}{0.912568in}}%
\pgfpathlineto{\pgfqpoint{1.585295in}{0.999965in}}%
\pgfpathlineto{\pgfqpoint{1.607591in}{0.917665in}}%
\pgfpathlineto{\pgfqpoint{1.652182in}{0.931245in}}%
\pgfpathlineto{\pgfqpoint{1.674477in}{0.900323in}}%
\pgfpathlineto{\pgfqpoint{1.696772in}{0.913584in}}%
\pgfpathlineto{\pgfqpoint{1.719068in}{0.852775in}}%
\pgfpathlineto{\pgfqpoint{1.763659in}{0.910752in}}%
\pgfpathlineto{\pgfqpoint{1.785954in}{0.929742in}}%
\pgfpathlineto{\pgfqpoint{1.808250in}{1.055342in}}%
\pgfpathlineto{\pgfqpoint{1.830545in}{1.078855in}}%
\pgfpathlineto{\pgfqpoint{1.852840in}{1.087951in}}%
\pgfpathlineto{\pgfqpoint{1.875136in}{1.056135in}}%
\pgfpathlineto{\pgfqpoint{1.897431in}{1.076829in}}%
\pgfpathlineto{\pgfqpoint{1.919727in}{1.033737in}}%
\pgfpathlineto{\pgfqpoint{1.942022in}{1.104144in}}%
\pgfpathlineto{\pgfqpoint{1.964318in}{1.076343in}}%
\pgfpathlineto{\pgfqpoint{1.986613in}{1.090119in}}%
\pgfpathlineto{\pgfqpoint{2.008908in}{1.202183in}}%
\pgfpathlineto{\pgfqpoint{2.031204in}{1.273833in}}%
\pgfpathlineto{\pgfqpoint{2.053499in}{1.294147in}}%
\pgfpathlineto{\pgfqpoint{2.075795in}{1.382600in}}%
\pgfpathlineto{\pgfqpoint{2.098090in}{1.385935in}}%
\pgfpathlineto{\pgfqpoint{2.120385in}{1.441803in}}%
\pgfpathlineto{\pgfqpoint{2.142681in}{1.449865in}}%
\pgfpathlineto{\pgfqpoint{2.164976in}{1.400420in}}%
\pgfpathlineto{\pgfqpoint{2.187272in}{1.433804in}}%
\pgfpathlineto{\pgfqpoint{2.209567in}{1.500011in}}%
\pgfpathlineto{\pgfqpoint{2.231863in}{1.484809in}}%
\pgfpathlineto{\pgfqpoint{2.254158in}{1.422578in}}%
\pgfpathlineto{\pgfqpoint{2.276453in}{1.396766in}}%
\pgfpathlineto{\pgfqpoint{2.298749in}{1.262570in}}%
\pgfpathlineto{\pgfqpoint{2.321044in}{1.234369in}}%
\pgfpathlineto{\pgfqpoint{2.343340in}{1.180940in}}%
\pgfpathlineto{\pgfqpoint{2.365635in}{1.301044in}}%
\pgfpathlineto{\pgfqpoint{2.387931in}{1.274705in}}%
\pgfpathlineto{\pgfqpoint{2.410226in}{1.194215in}}%
\pgfpathlineto{\pgfqpoint{2.432521in}{1.248860in}}%
\pgfpathlineto{\pgfqpoint{2.454817in}{1.245138in}}%
\pgfpathlineto{\pgfqpoint{2.477112in}{1.254084in}}%
\pgfpathlineto{\pgfqpoint{2.521703in}{1.289434in}}%
\pgfpathlineto{\pgfqpoint{2.543998in}{1.274156in}}%
\pgfpathlineto{\pgfqpoint{2.566294in}{1.336258in}}%
\pgfpathlineto{\pgfqpoint{2.588589in}{1.374169in}}%
\pgfpathlineto{\pgfqpoint{2.610885in}{1.350921in}}%
\pgfpathlineto{\pgfqpoint{2.633180in}{1.347212in}}%
\pgfpathlineto{\pgfqpoint{2.655476in}{1.298025in}}%
\pgfpathlineto{\pgfqpoint{2.677771in}{1.320853in}}%
\pgfpathlineto{\pgfqpoint{2.700066in}{1.371407in}}%
\pgfpathlineto{\pgfqpoint{2.744657in}{1.463222in}}%
\pgfpathlineto{\pgfqpoint{2.766953in}{1.462134in}}%
\pgfpathlineto{\pgfqpoint{2.789248in}{1.571819in}}%
\pgfpathlineto{\pgfqpoint{2.811544in}{1.565306in}}%
\pgfpathlineto{\pgfqpoint{2.833839in}{1.574598in}}%
\pgfpathlineto{\pgfqpoint{2.856134in}{1.559504in}}%
\pgfpathlineto{\pgfqpoint{2.878430in}{1.572184in}}%
\pgfpathlineto{\pgfqpoint{2.900725in}{1.563453in}}%
\pgfpathlineto{\pgfqpoint{2.923021in}{1.710601in}}%
\pgfpathlineto{\pgfqpoint{2.945316in}{1.748941in}}%
\pgfpathlineto{\pgfqpoint{2.967611in}{1.638024in}}%
\pgfpathlineto{\pgfqpoint{2.989907in}{1.615821in}}%
\pgfpathlineto{\pgfqpoint{3.012202in}{1.528650in}}%
\pgfpathlineto{\pgfqpoint{3.034498in}{1.594309in}}%
\pgfpathlineto{\pgfqpoint{3.056793in}{1.696686in}}%
\pgfpathlineto{\pgfqpoint{3.079089in}{1.672692in}}%
\pgfpathlineto{\pgfqpoint{3.101384in}{1.723763in}}%
\pgfpathlineto{\pgfqpoint{3.123679in}{1.716574in}}%
\pgfpathlineto{\pgfqpoint{3.145975in}{1.753726in}}%
\pgfpathlineto{\pgfqpoint{3.168270in}{1.890951in}}%
\pgfpathlineto{\pgfqpoint{3.190566in}{1.872613in}}%
\pgfpathlineto{\pgfqpoint{3.212861in}{1.788580in}}%
\pgfpathlineto{\pgfqpoint{3.235157in}{1.850183in}}%
\pgfpathlineto{\pgfqpoint{3.257452in}{1.897842in}}%
\pgfpathlineto{\pgfqpoint{3.279747in}{1.958801in}}%
\pgfpathlineto{\pgfqpoint{3.302043in}{2.049902in}}%
\pgfpathlineto{\pgfqpoint{3.324338in}{2.003638in}}%
\pgfpathlineto{\pgfqpoint{3.346634in}{1.858985in}}%
\pgfpathlineto{\pgfqpoint{3.368929in}{1.947613in}}%
\pgfpathlineto{\pgfqpoint{3.391224in}{1.924247in}}%
\pgfpathlineto{\pgfqpoint{3.413520in}{1.840365in}}%
\pgfpathlineto{\pgfqpoint{3.435815in}{1.762151in}}%
\pgfpathlineto{\pgfqpoint{3.458111in}{1.853831in}}%
\pgfpathlineto{\pgfqpoint{3.480406in}{1.875845in}}%
\pgfpathlineto{\pgfqpoint{3.502702in}{1.907463in}}%
\pgfpathlineto{\pgfqpoint{3.524997in}{1.982667in}}%
\pgfpathlineto{\pgfqpoint{3.547292in}{1.937925in}}%
\pgfpathlineto{\pgfqpoint{3.569588in}{1.859802in}}%
\pgfpathlineto{\pgfqpoint{3.591883in}{1.804164in}}%
\pgfpathlineto{\pgfqpoint{3.614179in}{1.731692in}}%
\pgfpathlineto{\pgfqpoint{3.636474in}{1.718700in}}%
\pgfpathlineto{\pgfqpoint{3.658770in}{1.663656in}}%
\pgfpathlineto{\pgfqpoint{3.681065in}{1.673977in}}%
\pgfpathlineto{\pgfqpoint{3.703360in}{1.648164in}}%
\pgfpathlineto{\pgfqpoint{3.725656in}{1.650196in}}%
\pgfpathlineto{\pgfqpoint{3.747951in}{1.703318in}}%
\pgfpathlineto{\pgfqpoint{3.770247in}{1.812850in}}%
\pgfpathlineto{\pgfqpoint{3.792542in}{1.939741in}}%
\pgfpathlineto{\pgfqpoint{3.814838in}{1.957342in}}%
\pgfpathlineto{\pgfqpoint{3.837133in}{1.982806in}}%
\pgfpathlineto{\pgfqpoint{3.859428in}{1.951265in}}%
\pgfpathlineto{\pgfqpoint{3.881724in}{2.040428in}}%
\pgfpathlineto{\pgfqpoint{3.904019in}{2.177408in}}%
\pgfpathlineto{\pgfqpoint{3.926315in}{2.136259in}}%
\pgfpathlineto{\pgfqpoint{3.948610in}{2.210941in}}%
\pgfpathlineto{\pgfqpoint{3.970905in}{2.191922in}}%
\pgfpathlineto{\pgfqpoint{3.993201in}{2.186613in}}%
\pgfpathlineto{\pgfqpoint{4.015496in}{2.211957in}}%
\pgfpathlineto{\pgfqpoint{4.037792in}{2.186772in}}%
\pgfpathlineto{\pgfqpoint{4.060087in}{2.274442in}}%
\pgfpathlineto{\pgfqpoint{4.082383in}{2.154148in}}%
\pgfpathlineto{\pgfqpoint{4.104678in}{2.104849in}}%
\pgfpathlineto{\pgfqpoint{4.126973in}{2.104013in}}%
\pgfpathlineto{\pgfqpoint{4.149269in}{2.027190in}}%
\pgfpathlineto{\pgfqpoint{4.171564in}{2.065509in}}%
\pgfpathlineto{\pgfqpoint{4.193860in}{1.952576in}}%
\pgfpathlineto{\pgfqpoint{4.216155in}{1.965542in}}%
\pgfpathlineto{\pgfqpoint{4.238451in}{1.862555in}}%
\pgfpathlineto{\pgfqpoint{4.260746in}{1.859838in}}%
\pgfpathlineto{\pgfqpoint{4.283041in}{1.907032in}}%
\pgfpathlineto{\pgfqpoint{4.305337in}{2.092141in}}%
\pgfpathlineto{\pgfqpoint{4.327632in}{2.192463in}}%
\pgfpathlineto{\pgfqpoint{4.349928in}{2.345198in}}%
\pgfpathlineto{\pgfqpoint{4.372223in}{2.361357in}}%
\pgfpathlineto{\pgfqpoint{4.394518in}{2.571905in}}%
\pgfpathlineto{\pgfqpoint{4.439109in}{2.545112in}}%
\pgfpathlineto{\pgfqpoint{4.461405in}{2.433287in}}%
\pgfpathlineto{\pgfqpoint{4.483700in}{2.464496in}}%
\pgfpathlineto{\pgfqpoint{4.505996in}{2.576531in}}%
\pgfpathlineto{\pgfqpoint{4.528291in}{2.525500in}}%
\pgfpathlineto{\pgfqpoint{4.550586in}{2.500394in}}%
\pgfpathlineto{\pgfqpoint{4.572882in}{2.466681in}}%
\pgfpathlineto{\pgfqpoint{4.595177in}{2.445731in}}%
\pgfpathlineto{\pgfqpoint{4.617473in}{2.400851in}}%
\pgfpathlineto{\pgfqpoint{4.639768in}{2.288156in}}%
\pgfpathlineto{\pgfqpoint{4.662064in}{2.424063in}}%
\pgfpathlineto{\pgfqpoint{4.684359in}{2.255206in}}%
\pgfpathlineto{\pgfqpoint{4.706654in}{2.184115in}}%
\pgfpathlineto{\pgfqpoint{4.728950in}{2.223284in}}%
\pgfpathlineto{\pgfqpoint{4.751245in}{2.322599in}}%
\pgfpathlineto{\pgfqpoint{4.773541in}{2.313666in}}%
\pgfpathlineto{\pgfqpoint{4.795836in}{2.298081in}}%
\pgfpathlineto{\pgfqpoint{4.818131in}{2.312317in}}%
\pgfpathlineto{\pgfqpoint{4.840427in}{2.298699in}}%
\pgfpathlineto{\pgfqpoint{4.862722in}{2.388244in}}%
\pgfpathlineto{\pgfqpoint{4.885018in}{2.354003in}}%
\pgfpathlineto{\pgfqpoint{4.907313in}{2.275242in}}%
\pgfpathlineto{\pgfqpoint{4.929609in}{2.243207in}}%
\pgfpathlineto{\pgfqpoint{4.951904in}{2.261296in}}%
\pgfpathlineto{\pgfqpoint{4.974199in}{2.381925in}}%
\pgfpathlineto{\pgfqpoint{4.996495in}{2.543477in}}%
\pgfpathlineto{\pgfqpoint{5.018790in}{2.515452in}}%
\pgfpathlineto{\pgfqpoint{5.041086in}{2.387857in}}%
\pgfpathlineto{\pgfqpoint{5.063381in}{2.272658in}}%
\pgfpathlineto{\pgfqpoint{5.085677in}{2.477730in}}%
\pgfpathlineto{\pgfqpoint{5.107972in}{2.449973in}}%
\pgfpathlineto{\pgfqpoint{5.130267in}{2.503853in}}%
\pgfpathlineto{\pgfqpoint{5.152563in}{2.704339in}}%
\pgfpathlineto{\pgfqpoint{5.174858in}{2.646589in}}%
\pgfpathlineto{\pgfqpoint{5.197154in}{2.638056in}}%
\pgfpathlineto{\pgfqpoint{5.219449in}{2.591271in}}%
\pgfpathlineto{\pgfqpoint{5.241744in}{2.495524in}}%
\pgfpathlineto{\pgfqpoint{5.264040in}{2.583650in}}%
\pgfpathlineto{\pgfqpoint{5.286335in}{2.701330in}}%
\pgfpathlineto{\pgfqpoint{5.308631in}{2.526243in}}%
\pgfpathlineto{\pgfqpoint{5.330926in}{2.633375in}}%
\pgfpathlineto{\pgfqpoint{5.353222in}{2.620757in}}%
\pgfpathlineto{\pgfqpoint{5.397812in}{2.681032in}}%
\pgfpathlineto{\pgfqpoint{5.420108in}{2.561927in}}%
\pgfpathlineto{\pgfqpoint{5.442403in}{2.542315in}}%
\pgfpathlineto{\pgfqpoint{5.464699in}{2.457950in}}%
\pgfpathlineto{\pgfqpoint{5.486994in}{2.655230in}}%
\pgfpathlineto{\pgfqpoint{5.509290in}{2.894506in}}%
\pgfpathlineto{\pgfqpoint{5.531585in}{2.915463in}}%
\pgfpathlineto{\pgfqpoint{5.553880in}{2.801134in}}%
\pgfpathlineto{\pgfqpoint{5.576176in}{2.936815in}}%
\pgfpathlineto{\pgfqpoint{5.598471in}{2.928009in}}%
\pgfpathlineto{\pgfqpoint{5.620767in}{3.136311in}}%
\pgfpathlineto{\pgfqpoint{5.643062in}{3.384487in}}%
\pgfpathlineto{\pgfqpoint{5.665357in}{3.610560in}}%
\pgfpathlineto{\pgfqpoint{5.687653in}{3.573444in}}%
\pgfpathlineto{\pgfqpoint{5.709948in}{3.460160in}}%
\pgfpathlineto{\pgfqpoint{5.732244in}{3.391156in}}%
\pgfpathlineto{\pgfqpoint{5.754539in}{3.392767in}}%
\pgfpathlineto{\pgfqpoint{5.776835in}{3.293370in}}%
\pgfpathlineto{\pgfqpoint{5.799130in}{3.453586in}}%
\pgfpathlineto{\pgfqpoint{5.821425in}{3.550990in}}%
\pgfpathlineto{\pgfqpoint{5.843721in}{3.560695in}}%
\pgfpathlineto{\pgfqpoint{5.866016in}{3.523367in}}%
\pgfpathlineto{\pgfqpoint{5.888312in}{3.621398in}}%
\pgfpathlineto{\pgfqpoint{5.910607in}{3.820517in}}%
\pgfpathlineto{\pgfqpoint{5.932903in}{3.578133in}}%
\pgfpathlineto{\pgfqpoint{5.955198in}{3.622787in}}%
\pgfpathlineto{\pgfqpoint{5.977493in}{3.597868in}}%
\pgfpathlineto{\pgfqpoint{5.999789in}{3.516143in}}%
\pgfpathlineto{\pgfqpoint{6.022084in}{3.486694in}}%
\pgfpathlineto{\pgfqpoint{6.044380in}{3.517674in}}%
\pgfpathlineto{\pgfqpoint{6.066675in}{3.777968in}}%
\pgfpathlineto{\pgfqpoint{6.088970in}{3.800601in}}%
\pgfpathlineto{\pgfqpoint{6.111266in}{3.757135in}}%
\pgfpathlineto{\pgfqpoint{6.133561in}{3.608846in}}%
\pgfpathlineto{\pgfqpoint{6.155857in}{3.725603in}}%
\pgfpathlineto{\pgfqpoint{6.155857in}{3.725603in}}%
\pgfusepath{stroke}%
\end{pgfscope}%
\begin{pgfscope}%
\pgfpathrectangle{\pgfqpoint{0.448229in}{0.413680in}}{\pgfqpoint{5.707628in}{3.574842in}}%
\pgfusepath{clip}%
\pgfsetbuttcap%
\pgfsetroundjoin%
\pgfsetlinewidth{1.003750pt}%
\definecolor{currentstroke}{rgb}{0.690196,0.188235,0.376471}%
\pgfsetstrokecolor{currentstroke}%
\pgfsetdash{{3.700000pt}{1.600000pt}}{0.000000pt}%
\pgfpathmoveto{\pgfqpoint{0.448229in}{0.635932in}}%
\pgfpathlineto{\pgfqpoint{0.470524in}{0.642507in}}%
\pgfpathlineto{\pgfqpoint{0.492820in}{0.666042in}}%
\pgfpathlineto{\pgfqpoint{0.515115in}{0.632605in}}%
\pgfpathlineto{\pgfqpoint{0.537411in}{0.663178in}}%
\pgfpathlineto{\pgfqpoint{0.559706in}{0.651943in}}%
\pgfpathlineto{\pgfqpoint{0.582001in}{0.596549in}}%
\pgfpathlineto{\pgfqpoint{0.604297in}{0.577615in}}%
\pgfpathlineto{\pgfqpoint{0.626592in}{0.603396in}}%
\pgfpathlineto{\pgfqpoint{0.648888in}{0.620056in}}%
\pgfpathlineto{\pgfqpoint{0.671183in}{0.582269in}}%
\pgfpathlineto{\pgfqpoint{0.693479in}{0.608875in}}%
\pgfpathlineto{\pgfqpoint{0.715774in}{0.610229in}}%
\pgfpathlineto{\pgfqpoint{0.738069in}{0.630389in}}%
\pgfpathlineto{\pgfqpoint{0.760365in}{0.647435in}}%
\pgfpathlineto{\pgfqpoint{0.782660in}{0.608420in}}%
\pgfpathlineto{\pgfqpoint{0.804956in}{0.625838in}}%
\pgfpathlineto{\pgfqpoint{0.827251in}{0.607942in}}%
\pgfpathlineto{\pgfqpoint{0.849546in}{0.669100in}}%
\pgfpathlineto{\pgfqpoint{0.871842in}{0.701751in}}%
\pgfpathlineto{\pgfqpoint{0.894137in}{0.708513in}}%
\pgfpathlineto{\pgfqpoint{0.916433in}{0.674433in}}%
\pgfpathlineto{\pgfqpoint{0.938728in}{0.681369in}}%
\pgfpathlineto{\pgfqpoint{0.961024in}{0.692443in}}%
\pgfpathlineto{\pgfqpoint{0.983319in}{0.649933in}}%
\pgfpathlineto{\pgfqpoint{1.005614in}{0.738912in}}%
\pgfpathlineto{\pgfqpoint{1.027910in}{0.760032in}}%
\pgfpathlineto{\pgfqpoint{1.050205in}{0.835144in}}%
\pgfpathlineto{\pgfqpoint{1.072501in}{0.791998in}}%
\pgfpathlineto{\pgfqpoint{1.094796in}{0.866530in}}%
\pgfpathlineto{\pgfqpoint{1.117092in}{0.811333in}}%
\pgfpathlineto{\pgfqpoint{1.139387in}{0.788741in}}%
\pgfpathlineto{\pgfqpoint{1.161682in}{0.816926in}}%
\pgfpathlineto{\pgfqpoint{1.183978in}{0.824384in}}%
\pgfpathlineto{\pgfqpoint{1.206273in}{0.727599in}}%
\pgfpathlineto{\pgfqpoint{1.228569in}{0.688065in}}%
\pgfpathlineto{\pgfqpoint{1.250864in}{0.725694in}}%
\pgfpathlineto{\pgfqpoint{1.273159in}{0.814778in}}%
\pgfpathlineto{\pgfqpoint{1.295455in}{0.859534in}}%
\pgfpathlineto{\pgfqpoint{1.317750in}{0.852815in}}%
\pgfpathlineto{\pgfqpoint{1.340046in}{0.893573in}}%
\pgfpathlineto{\pgfqpoint{1.362341in}{0.847214in}}%
\pgfpathlineto{\pgfqpoint{1.384637in}{0.875186in}}%
\pgfpathlineto{\pgfqpoint{1.406932in}{0.820262in}}%
\pgfpathlineto{\pgfqpoint{1.429227in}{0.825259in}}%
\pgfpathlineto{\pgfqpoint{1.451523in}{0.825011in}}%
\pgfpathlineto{\pgfqpoint{1.473818in}{0.865211in}}%
\pgfpathlineto{\pgfqpoint{1.496114in}{0.857271in}}%
\pgfpathlineto{\pgfqpoint{1.518409in}{0.916149in}}%
\pgfpathlineto{\pgfqpoint{1.540705in}{0.912833in}}%
\pgfpathlineto{\pgfqpoint{1.585295in}{1.000458in}}%
\pgfpathlineto{\pgfqpoint{1.607591in}{0.916648in}}%
\pgfpathlineto{\pgfqpoint{1.652182in}{0.931638in}}%
\pgfpathlineto{\pgfqpoint{1.674477in}{0.901113in}}%
\pgfpathlineto{\pgfqpoint{1.696772in}{0.915037in}}%
\pgfpathlineto{\pgfqpoint{1.719068in}{0.853600in}}%
\pgfpathlineto{\pgfqpoint{1.763659in}{0.912435in}}%
\pgfpathlineto{\pgfqpoint{1.785954in}{0.932045in}}%
\pgfpathlineto{\pgfqpoint{1.808250in}{1.053555in}}%
\pgfpathlineto{\pgfqpoint{1.830545in}{1.077662in}}%
\pgfpathlineto{\pgfqpoint{1.852840in}{1.087527in}}%
\pgfpathlineto{\pgfqpoint{1.875136in}{1.056214in}}%
\pgfpathlineto{\pgfqpoint{1.897431in}{1.077568in}}%
\pgfpathlineto{\pgfqpoint{1.919727in}{1.034681in}}%
\pgfpathlineto{\pgfqpoint{1.942022in}{1.104400in}}%
\pgfpathlineto{\pgfqpoint{1.964318in}{1.077175in}}%
\pgfpathlineto{\pgfqpoint{1.986613in}{1.091700in}}%
\pgfpathlineto{\pgfqpoint{2.008908in}{1.201039in}}%
\pgfpathlineto{\pgfqpoint{2.031204in}{1.272089in}}%
\pgfpathlineto{\pgfqpoint{2.053499in}{1.293168in}}%
\pgfpathlineto{\pgfqpoint{2.075795in}{1.380440in}}%
\pgfpathlineto{\pgfqpoint{2.098090in}{1.384717in}}%
\pgfpathlineto{\pgfqpoint{2.120385in}{1.440715in}}%
\pgfpathlineto{\pgfqpoint{2.142681in}{1.449735in}}%
\pgfpathlineto{\pgfqpoint{2.164976in}{1.400658in}}%
\pgfpathlineto{\pgfqpoint{2.187272in}{1.434724in}}%
\pgfpathlineto{\pgfqpoint{2.209567in}{1.500862in}}%
\pgfpathlineto{\pgfqpoint{2.231863in}{1.486605in}}%
\pgfpathlineto{\pgfqpoint{2.254158in}{1.424354in}}%
\pgfpathlineto{\pgfqpoint{2.276453in}{1.399319in}}%
\pgfpathlineto{\pgfqpoint{2.298749in}{1.261083in}}%
\pgfpathlineto{\pgfqpoint{2.321044in}{1.233576in}}%
\pgfpathlineto{\pgfqpoint{2.343340in}{1.180237in}}%
\pgfpathlineto{\pgfqpoint{2.365635in}{1.297189in}}%
\pgfpathlineto{\pgfqpoint{2.387931in}{1.271625in}}%
\pgfpathlineto{\pgfqpoint{2.410226in}{1.190352in}}%
\pgfpathlineto{\pgfqpoint{2.432521in}{1.244903in}}%
\pgfpathlineto{\pgfqpoint{2.454817in}{1.242067in}}%
\pgfpathlineto{\pgfqpoint{2.477112in}{1.251857in}}%
\pgfpathlineto{\pgfqpoint{2.521703in}{1.288782in}}%
\pgfpathlineto{\pgfqpoint{2.543998in}{1.274351in}}%
\pgfpathlineto{\pgfqpoint{2.566294in}{1.336328in}}%
\pgfpathlineto{\pgfqpoint{2.588589in}{1.374795in}}%
\pgfpathlineto{\pgfqpoint{2.610885in}{1.352345in}}%
\pgfpathlineto{\pgfqpoint{2.633180in}{1.349565in}}%
\pgfpathlineto{\pgfqpoint{2.655476in}{1.300605in}}%
\pgfpathlineto{\pgfqpoint{2.677771in}{1.324236in}}%
\pgfpathlineto{\pgfqpoint{2.700066in}{1.375134in}}%
\pgfpathlineto{\pgfqpoint{2.744657in}{1.467987in}}%
\pgfpathlineto{\pgfqpoint{2.766953in}{1.467887in}}%
\pgfpathlineto{\pgfqpoint{2.789248in}{1.576000in}}%
\pgfpathlineto{\pgfqpoint{2.811544in}{1.570507in}}%
\pgfpathlineto{\pgfqpoint{2.833839in}{1.580844in}}%
\pgfpathlineto{\pgfqpoint{2.856134in}{1.566699in}}%
\pgfpathlineto{\pgfqpoint{2.878430in}{1.580423in}}%
\pgfpathlineto{\pgfqpoint{2.900725in}{1.572687in}}%
\pgfpathlineto{\pgfqpoint{2.923021in}{1.716627in}}%
\pgfpathlineto{\pgfqpoint{2.945316in}{1.755860in}}%
\pgfpathlineto{\pgfqpoint{2.967611in}{1.642986in}}%
\pgfpathlineto{\pgfqpoint{2.989907in}{1.621696in}}%
\pgfpathlineto{\pgfqpoint{3.012202in}{1.533560in}}%
\pgfpathlineto{\pgfqpoint{3.034498in}{1.599385in}}%
\pgfpathlineto{\pgfqpoint{3.056793in}{1.700697in}}%
\pgfpathlineto{\pgfqpoint{3.079089in}{1.677638in}}%
\pgfpathlineto{\pgfqpoint{3.101384in}{1.729339in}}%
\pgfpathlineto{\pgfqpoint{3.123679in}{1.723240in}}%
\pgfpathlineto{\pgfqpoint{3.145975in}{1.761317in}}%
\pgfpathlineto{\pgfqpoint{3.168270in}{1.896202in}}%
\pgfpathlineto{\pgfqpoint{3.190566in}{1.878957in}}%
\pgfpathlineto{\pgfqpoint{3.212861in}{1.794409in}}%
\pgfpathlineto{\pgfqpoint{3.235157in}{1.856520in}}%
\pgfpathlineto{\pgfqpoint{3.257452in}{1.905019in}}%
\pgfpathlineto{\pgfqpoint{3.279747in}{1.966616in}}%
\pgfpathlineto{\pgfqpoint{3.302043in}{2.057623in}}%
\pgfpathlineto{\pgfqpoint{3.324338in}{2.012085in}}%
\pgfpathlineto{\pgfqpoint{3.346634in}{1.863962in}}%
\pgfpathlineto{\pgfqpoint{3.368929in}{1.952371in}}%
\pgfpathlineto{\pgfqpoint{3.391224in}{1.930079in}}%
\pgfpathlineto{\pgfqpoint{3.413520in}{1.845766in}}%
\pgfpathlineto{\pgfqpoint{3.435815in}{1.767248in}}%
\pgfpathlineto{\pgfqpoint{3.458111in}{1.858501in}}%
\pgfpathlineto{\pgfqpoint{3.480406in}{1.881640in}}%
\pgfpathlineto{\pgfqpoint{3.502702in}{1.914323in}}%
\pgfpathlineto{\pgfqpoint{3.524997in}{1.989820in}}%
\pgfpathlineto{\pgfqpoint{3.547292in}{1.945800in}}%
\pgfpathlineto{\pgfqpoint{3.569588in}{1.867405in}}%
\pgfpathlineto{\pgfqpoint{3.591883in}{1.812128in}}%
\pgfpathlineto{\pgfqpoint{3.614179in}{1.739435in}}%
\pgfpathlineto{\pgfqpoint{3.636474in}{1.727487in}}%
\pgfpathlineto{\pgfqpoint{3.658770in}{1.672671in}}%
\pgfpathlineto{\pgfqpoint{3.681065in}{1.684101in}}%
\pgfpathlineto{\pgfqpoint{3.703360in}{1.659117in}}%
\pgfpathlineto{\pgfqpoint{3.725656in}{1.662243in}}%
\pgfpathlineto{\pgfqpoint{3.747951in}{1.716110in}}%
\pgfpathlineto{\pgfqpoint{3.770247in}{1.824793in}}%
\pgfpathlineto{\pgfqpoint{3.792542in}{1.950198in}}%
\pgfpathlineto{\pgfqpoint{3.814838in}{1.969037in}}%
\pgfpathlineto{\pgfqpoint{3.837133in}{1.995727in}}%
\pgfpathlineto{\pgfqpoint{3.859428in}{1.965078in}}%
\pgfpathlineto{\pgfqpoint{3.881724in}{2.054412in}}%
\pgfpathlineto{\pgfqpoint{3.904019in}{2.189904in}}%
\pgfpathlineto{\pgfqpoint{3.926315in}{2.149598in}}%
\pgfpathlineto{\pgfqpoint{3.948610in}{2.224953in}}%
\pgfpathlineto{\pgfqpoint{3.970905in}{2.207134in}}%
\pgfpathlineto{\pgfqpoint{3.993201in}{2.203142in}}%
\pgfpathlineto{\pgfqpoint{4.015496in}{2.229873in}}%
\pgfpathlineto{\pgfqpoint{4.037792in}{2.205773in}}%
\pgfpathlineto{\pgfqpoint{4.060087in}{2.294022in}}%
\pgfpathlineto{\pgfqpoint{4.082383in}{2.171649in}}%
\pgfpathlineto{\pgfqpoint{4.104678in}{2.122913in}}%
\pgfpathlineto{\pgfqpoint{4.126973in}{2.123382in}}%
\pgfpathlineto{\pgfqpoint{4.149269in}{2.046176in}}%
\pgfpathlineto{\pgfqpoint{4.171564in}{2.085766in}}%
\pgfpathlineto{\pgfqpoint{4.193860in}{1.970738in}}%
\pgfpathlineto{\pgfqpoint{4.216155in}{1.985002in}}%
\pgfpathlineto{\pgfqpoint{4.238451in}{1.880293in}}%
\pgfpathlineto{\pgfqpoint{4.260746in}{1.878747in}}%
\pgfpathlineto{\pgfqpoint{4.283041in}{1.927045in}}%
\pgfpathlineto{\pgfqpoint{4.305337in}{2.108206in}}%
\pgfpathlineto{\pgfqpoint{4.327632in}{2.208579in}}%
\pgfpathlineto{\pgfqpoint{4.349928in}{2.359453in}}%
\pgfpathlineto{\pgfqpoint{4.372223in}{2.377074in}}%
\pgfpathlineto{\pgfqpoint{4.394518in}{2.582914in}}%
\pgfpathlineto{\pgfqpoint{4.439109in}{2.559029in}}%
\pgfpathlineto{\pgfqpoint{4.461405in}{2.446157in}}%
\pgfpathlineto{\pgfqpoint{4.483700in}{2.478806in}}%
\pgfpathlineto{\pgfqpoint{4.505996in}{2.590823in}}%
\pgfpathlineto{\pgfqpoint{4.528291in}{2.540678in}}%
\pgfpathlineto{\pgfqpoint{4.550586in}{2.516861in}}%
\pgfpathlineto{\pgfqpoint{4.572882in}{2.484282in}}%
\pgfpathlineto{\pgfqpoint{4.595177in}{2.464625in}}%
\pgfpathlineto{\pgfqpoint{4.617473in}{2.420600in}}%
\pgfpathlineto{\pgfqpoint{4.639768in}{2.306398in}}%
\pgfpathlineto{\pgfqpoint{4.662064in}{2.441415in}}%
\pgfpathlineto{\pgfqpoint{4.684359in}{2.268053in}}%
\pgfpathlineto{\pgfqpoint{4.706654in}{2.197104in}}%
\pgfpathlineto{\pgfqpoint{4.728950in}{2.237528in}}%
\pgfpathlineto{\pgfqpoint{4.751245in}{2.336980in}}%
\pgfpathlineto{\pgfqpoint{4.773541in}{2.329405in}}%
\pgfpathlineto{\pgfqpoint{4.795836in}{2.315103in}}%
\pgfpathlineto{\pgfqpoint{4.818131in}{2.330794in}}%
\pgfpathlineto{\pgfqpoint{4.840427in}{2.318469in}}%
\pgfpathlineto{\pgfqpoint{4.862722in}{2.408657in}}%
\pgfpathlineto{\pgfqpoint{4.885018in}{2.375428in}}%
\pgfpathlineto{\pgfqpoint{4.907313in}{2.296432in}}%
\pgfpathlineto{\pgfqpoint{4.929609in}{2.265371in}}%
\pgfpathlineto{\pgfqpoint{4.951904in}{2.284924in}}%
\pgfpathlineto{\pgfqpoint{4.974199in}{2.405427in}}%
\pgfpathlineto{\pgfqpoint{4.996495in}{2.565408in}}%
\pgfpathlineto{\pgfqpoint{5.018790in}{2.538584in}}%
\pgfpathlineto{\pgfqpoint{5.041086in}{2.408803in}}%
\pgfpathlineto{\pgfqpoint{5.063381in}{2.291912in}}%
\pgfpathlineto{\pgfqpoint{5.085677in}{2.492645in}}%
\pgfpathlineto{\pgfqpoint{5.107972in}{2.466115in}}%
\pgfpathlineto{\pgfqpoint{5.130267in}{2.521285in}}%
\pgfpathlineto{\pgfqpoint{5.152563in}{2.718099in}}%
\pgfpathlineto{\pgfqpoint{5.174858in}{2.661188in}}%
\pgfpathlineto{\pgfqpoint{5.197154in}{2.654179in}}%
\pgfpathlineto{\pgfqpoint{5.219449in}{2.608382in}}%
\pgfpathlineto{\pgfqpoint{5.241744in}{2.512172in}}%
\pgfpathlineto{\pgfqpoint{5.264040in}{2.601027in}}%
\pgfpathlineto{\pgfqpoint{5.286335in}{2.718742in}}%
\pgfpathlineto{\pgfqpoint{5.308631in}{2.539426in}}%
\pgfpathlineto{\pgfqpoint{5.330926in}{2.646706in}}%
\pgfpathlineto{\pgfqpoint{5.353222in}{2.635578in}}%
\pgfpathlineto{\pgfqpoint{5.397812in}{2.699001in}}%
\pgfpathlineto{\pgfqpoint{5.420108in}{2.578575in}}%
\pgfpathlineto{\pgfqpoint{5.442403in}{2.560330in}}%
\pgfpathlineto{\pgfqpoint{5.464699in}{2.475826in}}%
\pgfpathlineto{\pgfqpoint{5.486994in}{2.669546in}}%
\pgfpathlineto{\pgfqpoint{5.509290in}{2.902948in}}%
\pgfpathlineto{\pgfqpoint{5.531585in}{2.925585in}}%
\pgfpathlineto{\pgfqpoint{5.553880in}{2.810700in}}%
\pgfpathlineto{\pgfqpoint{5.576176in}{2.945741in}}%
\pgfpathlineto{\pgfqpoint{5.598471in}{2.938614in}}%
\pgfpathlineto{\pgfqpoint{5.620767in}{3.143262in}}%
\pgfpathlineto{\pgfqpoint{5.643062in}{3.385728in}}%
\pgfpathlineto{\pgfqpoint{5.665357in}{3.607560in}}%
\pgfpathlineto{\pgfqpoint{5.687653in}{3.572341in}}%
\pgfpathlineto{\pgfqpoint{5.709948in}{3.459526in}}%
\pgfpathlineto{\pgfqpoint{5.732244in}{3.391899in}}%
\pgfpathlineto{\pgfqpoint{5.754539in}{3.395440in}}%
\pgfpathlineto{\pgfqpoint{5.776835in}{3.296638in}}%
\pgfpathlineto{\pgfqpoint{5.799130in}{3.455632in}}%
\pgfpathlineto{\pgfqpoint{5.821425in}{3.553885in}}%
\pgfpathlineto{\pgfqpoint{5.843721in}{3.565595in}}%
\pgfpathlineto{\pgfqpoint{5.866016in}{3.530068in}}%
\pgfpathlineto{\pgfqpoint{5.888312in}{3.629099in}}%
\pgfpathlineto{\pgfqpoint{5.910607in}{3.826030in}}%
\pgfpathlineto{\pgfqpoint{5.932903in}{3.578515in}}%
\pgfpathlineto{\pgfqpoint{5.955198in}{3.624956in}}%
\pgfpathlineto{\pgfqpoint{5.977493in}{3.601995in}}%
\pgfpathlineto{\pgfqpoint{5.999789in}{3.521407in}}%
\pgfpathlineto{\pgfqpoint{6.022084in}{3.493808in}}%
\pgfpathlineto{\pgfqpoint{6.044380in}{3.526704in}}%
\pgfpathlineto{\pgfqpoint{6.066675in}{3.781588in}}%
\pgfpathlineto{\pgfqpoint{6.088970in}{3.806303in}}%
\pgfpathlineto{\pgfqpoint{6.111266in}{3.764694in}}%
\pgfpathlineto{\pgfqpoint{6.133561in}{3.615649in}}%
\pgfpathlineto{\pgfqpoint{6.155857in}{3.733030in}}%
\pgfpathlineto{\pgfqpoint{6.155857in}{3.733030in}}%
\pgfusepath{stroke}%
\end{pgfscope}%
\begin{pgfscope}%
\pgfpathrectangle{\pgfqpoint{0.448229in}{0.413680in}}{\pgfqpoint{5.707628in}{3.574842in}}%
\pgfusepath{clip}%
\pgfsetbuttcap%
\pgfsetroundjoin%
\pgfsetlinewidth{1.003750pt}%
\definecolor{currentstroke}{rgb}{0.690196,0.690196,0.690196}%
\pgfsetstrokecolor{currentstroke}%
\pgfsetdash{{3.700000pt}{1.600000pt}}{0.000000pt}%
\pgfpathmoveto{\pgfqpoint{0.448229in}{0.635932in}}%
\pgfpathlineto{\pgfqpoint{0.537411in}{0.663954in}}%
\pgfpathlineto{\pgfqpoint{0.626592in}{0.604469in}}%
\pgfpathlineto{\pgfqpoint{0.715774in}{0.612326in}}%
\pgfpathlineto{\pgfqpoint{0.804956in}{0.628903in}}%
\pgfpathlineto{\pgfqpoint{0.894137in}{0.711171in}}%
\pgfpathlineto{\pgfqpoint{0.983319in}{0.652357in}}%
\pgfpathlineto{\pgfqpoint{1.072501in}{0.792723in}}%
\pgfpathlineto{\pgfqpoint{1.161682in}{0.820930in}}%
\pgfpathlineto{\pgfqpoint{1.250864in}{0.731099in}}%
\pgfpathlineto{\pgfqpoint{1.340046in}{0.893735in}}%
\pgfpathlineto{\pgfqpoint{1.429227in}{0.825895in}}%
\pgfpathlineto{\pgfqpoint{1.518409in}{0.915728in}}%
\pgfpathlineto{\pgfqpoint{1.607591in}{0.919678in}}%
\pgfpathlineto{\pgfqpoint{1.696772in}{0.918480in}}%
\pgfpathlineto{\pgfqpoint{1.785954in}{0.937496in}}%
\pgfpathlineto{\pgfqpoint{1.875136in}{1.062381in}}%
\pgfpathlineto{\pgfqpoint{1.964318in}{1.085716in}}%
\pgfpathlineto{\pgfqpoint{2.053499in}{1.294602in}}%
\pgfpathlineto{\pgfqpoint{2.142681in}{1.447798in}}%
\pgfpathlineto{\pgfqpoint{2.231863in}{1.486332in}}%
\pgfpathlineto{\pgfqpoint{2.321044in}{1.223002in}}%
\pgfpathlineto{\pgfqpoint{2.410226in}{1.186137in}}%
\pgfpathlineto{\pgfqpoint{2.499408in}{1.264438in}}%
\pgfpathlineto{\pgfqpoint{2.588589in}{1.367938in}}%
\pgfpathlineto{\pgfqpoint{2.677771in}{1.317818in}}%
\pgfpathlineto{\pgfqpoint{2.766953in}{1.457504in}}%
\pgfpathlineto{\pgfqpoint{2.856134in}{1.556333in}}%
\pgfpathlineto{\pgfqpoint{2.945316in}{1.741653in}}%
\pgfpathlineto{\pgfqpoint{3.034498in}{1.586611in}}%
\pgfpathlineto{\pgfqpoint{3.123679in}{1.709310in}}%
\pgfpathlineto{\pgfqpoint{3.212861in}{1.784414in}}%
\pgfpathlineto{\pgfqpoint{3.302043in}{2.036351in}}%
\pgfpathlineto{\pgfqpoint{3.391224in}{1.913317in}}%
\pgfpathlineto{\pgfqpoint{3.480406in}{1.869488in}}%
\pgfpathlineto{\pgfqpoint{3.569588in}{1.858185in}}%
\pgfpathlineto{\pgfqpoint{3.658770in}{1.658505in}}%
\pgfpathlineto{\pgfqpoint{3.747951in}{1.702065in}}%
\pgfpathlineto{\pgfqpoint{3.837133in}{1.969891in}}%
\pgfpathlineto{\pgfqpoint{3.926315in}{2.123087in}}%
\pgfpathlineto{\pgfqpoint{4.015496in}{2.202624in}}%
\pgfpathlineto{\pgfqpoint{4.104678in}{2.099162in}}%
\pgfpathlineto{\pgfqpoint{4.193860in}{1.947902in}}%
\pgfpathlineto{\pgfqpoint{4.283041in}{1.906930in}}%
\pgfpathlineto{\pgfqpoint{4.372223in}{2.328778in}}%
\pgfpathlineto{\pgfqpoint{4.461405in}{2.404693in}}%
\pgfpathlineto{\pgfqpoint{4.550586in}{2.476222in}}%
\pgfpathlineto{\pgfqpoint{4.639768in}{2.263943in}}%
\pgfpathlineto{\pgfqpoint{4.728950in}{2.204399in}}%
\pgfpathlineto{\pgfqpoint{4.818131in}{2.296849in}}%
\pgfpathlineto{\pgfqpoint{4.907313in}{2.265318in}}%
\pgfpathlineto{\pgfqpoint{4.996495in}{2.526249in}}%
\pgfpathlineto{\pgfqpoint{5.085677in}{2.466197in}}%
\pgfpathlineto{\pgfqpoint{5.174858in}{2.635826in}}%
\pgfpathlineto{\pgfqpoint{5.264040in}{2.578722in}}%
\pgfpathlineto{\pgfqpoint{5.353222in}{2.621689in}}%
\pgfpathlineto{\pgfqpoint{5.442403in}{2.548461in}}%
\pgfpathlineto{\pgfqpoint{5.531585in}{2.907608in}}%
\pgfpathlineto{\pgfqpoint{5.620767in}{3.128022in}}%
\pgfpathlineto{\pgfqpoint{5.709948in}{3.444983in}}%
\pgfpathlineto{\pgfqpoint{5.799130in}{3.446256in}}%
\pgfpathlineto{\pgfqpoint{5.888312in}{3.618186in}}%
\pgfpathlineto{\pgfqpoint{5.977493in}{3.602774in}}%
\pgfpathlineto{\pgfqpoint{6.066675in}{3.787427in}}%
\pgfpathlineto{\pgfqpoint{6.155857in}{3.743125in}}%
\pgfusepath{stroke}%
\end{pgfscope}%
\begin{pgfscope}%
\pgfsetrectcap%
\pgfsetmiterjoin%
\pgfsetlinewidth{0.501875pt}%
\definecolor{currentstroke}{rgb}{0.000000,0.000000,0.000000}%
\pgfsetstrokecolor{currentstroke}%
\pgfsetdash{}{0pt}%
\pgfpathmoveto{\pgfqpoint{0.448229in}{0.413680in}}%
\pgfpathlineto{\pgfqpoint{0.448229in}{3.988523in}}%
\pgfusepath{stroke}%
\end{pgfscope}%
\begin{pgfscope}%
\pgfsetrectcap%
\pgfsetmiterjoin%
\pgfsetlinewidth{0.501875pt}%
\definecolor{currentstroke}{rgb}{0.000000,0.000000,0.000000}%
\pgfsetstrokecolor{currentstroke}%
\pgfsetdash{}{0pt}%
\pgfpathmoveto{\pgfqpoint{6.155857in}{0.413680in}}%
\pgfpathlineto{\pgfqpoint{6.155857in}{3.988523in}}%
\pgfusepath{stroke}%
\end{pgfscope}%
\begin{pgfscope}%
\pgfsetrectcap%
\pgfsetmiterjoin%
\pgfsetlinewidth{0.501875pt}%
\definecolor{currentstroke}{rgb}{0.000000,0.000000,0.000000}%
\pgfsetstrokecolor{currentstroke}%
\pgfsetdash{}{0pt}%
\pgfpathmoveto{\pgfqpoint{0.448229in}{0.413680in}}%
\pgfpathlineto{\pgfqpoint{6.155857in}{0.413680in}}%
\pgfusepath{stroke}%
\end{pgfscope}%
\begin{pgfscope}%
\pgfsetrectcap%
\pgfsetmiterjoin%
\pgfsetlinewidth{0.501875pt}%
\definecolor{currentstroke}{rgb}{0.000000,0.000000,0.000000}%
\pgfsetstrokecolor{currentstroke}%
\pgfsetdash{}{0pt}%
\pgfpathmoveto{\pgfqpoint{0.448229in}{3.988523in}}%
\pgfpathlineto{\pgfqpoint{6.155857in}{3.988523in}}%
\pgfusepath{stroke}%
\end{pgfscope}%
\begin{pgfscope}%
\pgfsetbuttcap%
\pgfsetmiterjoin%
\definecolor{currentfill}{rgb}{1.000000,1.000000,1.000000}%
\pgfsetfillcolor{currentfill}%
\pgfsetfillopacity{0.800000}%
\pgfsetlinewidth{1.003750pt}%
\definecolor{currentstroke}{rgb}{0.800000,0.800000,0.800000}%
\pgfsetstrokecolor{currentstroke}%
\pgfsetstrokeopacity{0.800000}%
\pgfsetdash{}{0pt}%
\pgfpathmoveto{\pgfqpoint{0.545451in}{3.279808in}}%
\pgfpathlineto{\pgfqpoint{2.192320in}{3.279808in}}%
\pgfpathquadraticcurveto{\pgfqpoint{2.220098in}{3.279808in}}{\pgfqpoint{2.220098in}{3.307586in}}%
\pgfpathlineto{\pgfqpoint{2.220098in}{3.891300in}}%
\pgfpathquadraticcurveto{\pgfqpoint{2.220098in}{3.919078in}}{\pgfqpoint{2.192320in}{3.919078in}}%
\pgfpathlineto{\pgfqpoint{0.545451in}{3.919078in}}%
\pgfpathquadraticcurveto{\pgfqpoint{0.517673in}{3.919078in}}{\pgfqpoint{0.517673in}{3.891300in}}%
\pgfpathlineto{\pgfqpoint{0.517673in}{3.307586in}}%
\pgfpathquadraticcurveto{\pgfqpoint{0.517673in}{3.279808in}}{\pgfqpoint{0.545451in}{3.279808in}}%
\pgfpathlineto{\pgfqpoint{0.545451in}{3.279808in}}%
\pgfpathclose%
\pgfusepath{stroke,fill}%
\end{pgfscope}%
\begin{pgfscope}%
\pgfsetrectcap%
\pgfsetroundjoin%
\pgfsetlinewidth{1.003750pt}%
\definecolor{currentstroke}{rgb}{0.000000,0.000000,0.000000}%
\pgfsetstrokecolor{currentstroke}%
\pgfsetdash{}{0pt}%
\pgfpathmoveto{\pgfqpoint{0.573229in}{3.814738in}}%
\pgfpathlineto{\pgfqpoint{0.712118in}{3.814738in}}%
\pgfpathlineto{\pgfqpoint{0.851007in}{3.814738in}}%
\pgfusepath{stroke}%
\end{pgfscope}%
\begin{pgfscope}%
\definecolor{textcolor}{rgb}{0.000000,0.000000,0.000000}%
\pgfsetstrokecolor{textcolor}%
\pgfsetfillcolor{textcolor}%
\pgftext[x=0.962118in,y=3.766127in,left,base]{\color{textcolor}{\rmfamily\fontsize{10.000000}{12.000000}\selectfont\catcode`\^=\active\def^{\ifmmode\sp\else\^{}\fi}\catcode`\%=\active\def%{\%}Exact ($Y_t$)}}%
\end{pgfscope}%
\begin{pgfscope}%
\pgfsetbuttcap%
\pgfsetroundjoin%
\pgfsetlinewidth{1.003750pt}%
\definecolor{currentstroke}{rgb}{0.690196,0.188235,0.376471}%
\pgfsetstrokecolor{currentstroke}%
\pgfsetdash{{3.700000pt}{1.600000pt}}{0.000000pt}%
\pgfpathmoveto{\pgfqpoint{0.573229in}{3.615537in}}%
\pgfpathlineto{\pgfqpoint{0.712118in}{3.615537in}}%
\pgfpathlineto{\pgfqpoint{0.851007in}{3.615537in}}%
\pgfusepath{stroke}%
\end{pgfscope}%
\begin{pgfscope}%
\definecolor{textcolor}{rgb}{0.000000,0.000000,0.000000}%
\pgfsetstrokecolor{textcolor}%
\pgfsetfillcolor{textcolor}%
\pgftext[x=0.962118in,y=3.566926in,left,base]{\color{textcolor}{\rmfamily\fontsize{10.000000}{12.000000}\selectfont\catcode`\^=\active\def^{\ifmmode\sp\else\^{}\fi}\catcode`\%=\active\def%{\%}EM ($X_t$): Fine Grid}}%
\end{pgfscope}%
\begin{pgfscope}%
\pgfsetbuttcap%
\pgfsetroundjoin%
\pgfsetlinewidth{1.003750pt}%
\definecolor{currentstroke}{rgb}{0.690196,0.690196,0.690196}%
\pgfsetstrokecolor{currentstroke}%
\pgfsetdash{{3.700000pt}{1.600000pt}}{0.000000pt}%
\pgfpathmoveto{\pgfqpoint{0.573229in}{3.416335in}}%
\pgfpathlineto{\pgfqpoint{0.712118in}{3.416335in}}%
\pgfpathlineto{\pgfqpoint{0.851007in}{3.416335in}}%
\pgfusepath{stroke}%
\end{pgfscope}%
\begin{pgfscope}%
\definecolor{textcolor}{rgb}{0.000000,0.000000,0.000000}%
\pgfsetstrokecolor{textcolor}%
\pgfsetfillcolor{textcolor}%
\pgftext[x=0.962118in,y=3.367724in,left,base]{\color{textcolor}{\rmfamily\fontsize{10.000000}{12.000000}\selectfont\catcode`\^=\active\def^{\ifmmode\sp\else\^{}\fi}\catcode`\%=\active\def%{\%}EM ($X_t$): Coarse Grid}}%
\end{pgfscope}%
\end{pgfpicture}%
\makeatother%
\endgroup%

    \caption{Comparison of Euler-Maruyama approximations against the exact analytical solution for a GBM. The simulation uses a fine grid step size of $\Delta t = 2^{-8}$ and a coarse grid step size of $\Delta t_{\text{coarse}} = 4\Delta t$. Parameters: $\mu=1.0$, $\sigma=0.5$, $X_0=1.0$.}
    \label{fig:eulerMaruyamaCoarseFineComparison}
\end{figure}

\clearpage
\subsection{Milstein}
In the Milstein scheme we use \cref{eq:SDE} to expand $a(X_s)$ as
\begin{align*}
    a(X_s) &\simeq a(X_{t_k}) + a'(X_{t_k}) (X_s - X_{t_k}) \\
    &\simeq a(X_{t_k}) + a'(X_{t_k})(b(X_{t_k}) (s-t_k) + a(X_{t_k}) (W_s - W_{t_k}),
\end{align*}
$0 \leq t_q \leq s$. As a consequence, we get
\begin{align*}
    \hat{X}_{t_{k+1}}^N &= \hat{X}_{t_k}^N + \int_{t_k}^{t_{k+1}} b(X_s) ds + \int_{t_k}^{t_{k+1}} a(X_s) dW_s \\
    &\simeq \hat{X}_{t_k}^N + \int_{t_k}^{t_{k+1}} b(X_s) ds + a(X_{t_k})(W_{t_{k+1}} - W_{t_k}) \\
    &+ a'(X_{t_k})b(X_{t_k}) \int_{t_k}^{t_{k+1}} (s - t_k) dW_s \\
    &+ a'(X_{t_k})a(X_{t_k}) \int_{t_k}^{t_{k+1}} (W_s - W_{t_k}) dW_s \\
    &\simeq \hat{X}_{t_k}^N + \int_{t_k}^{t_{k+1}} b(X_s) ds + a(X_{t_k})(W_{t_{k+1}} - W_{t_k}) \\
    &+ a'(X_{t_k})a(X_{t_k}) \int_{t_k}^{t_{k+1}} (W_s - W_{t_k}) dW_s.
\end{align*}
Next, using \ Ito's formula, we note that
\begin{equation*}
    (W_{t_{k+1}} - W_{t_k})^2 = 2 \int_{t_k}^{t_{k+1}} (W_s - W_k)dW_s + \int_{t_k}^{t_{k+1}} ds,
\end{equation*}
hence
\begin{equation*}
    \int_{t_k}^{t_{k+1}} (W_s - W_{t_k}) dW_s = \frac{1}{2} ((W_{t_{k+1}}-W_{t_k})^2 - (t_{k+1} - t_k))
\end{equation*}
and
\begin{align*}
    \hat{X}_{t_{k+1}}^N &\simeq \hat{X}_{t_k}^N + \int_{t_k}^{t_{k+1}} b(X_s) ds + a(X_{t_k})(W_{t_{k+1}} - W_{t_k}) \\
    &\quad + \frac{1}{2} a'(X_{t_k})a(X_{t_k})((W_{t_{k+1}} - W_{t_k})^2 - (t_{k+1} - t_k)) \\
    &\simeq \hat{X}_{t_k}^N + b(X_{t_k})(t_{k+1} - t_k) + a(X_{t_k})(W_{t_{k+1}} - W_{t_k}) \\
    &\quad + \frac{1}{2} a'(X_{t_k})a(X_{t_k})((W_{t_{k+1}} - W_{t_k})^2 - (t_{k+1} - t_k)).
\end{align*}
As a consequence, the Milstein scheme is written as
\begin{align*}
    \hat{X}_{t_{k+1}}^N &\simeq \hat{X}_{t_k}^N + b(\hat{X}_{t_k}^N)(t_{k+1} - t_k) + a(\hat{X}_{t_k}^N)(W_{t_{k+1}} - W_{t_k}) \\
    &\quad + \frac{1}{2} a' (\hat{X}_{t_k}^N)a(\hat{X}_{t_k}^N)((W_{t_{k+1}} - W_{t_k})^2 - (t_{k+1} - t_k)),
\end{align*}
i.e., in the Milstein scheme, we take into account the "small" difference
\begin{equation*}
    (W_{t_{k+1}} - W_{t_k})^2 - (t_{k+1} - t_k)
\end{equation*}
existing between $(\Delta W_t)^2$ and $\Delta t$. Taking $(\Delta W_t)^2$ equal to $\Delta t$ brings us back to the Euler scheme.

Let $T > 0$. Suppose that the initial data satisfies
\begin{align}
    \E \left(|X_0|^2\right) &< \infty, \label{eq:MilsteinInitialMoment} \\
    \E \left(|X_0 - Y^\delta_0|^2\right)^{1/2} &\leq K_1 \delta^{1/2}. \label{eq:MilsteinInitialError}
\end{align}
We assume higher-order regularity on the coefficients compared to the Euler scheme. Specifically, let the differential operator $\scL^j$ be defined appropriately for the diffusion terms. Assume the coefficients satisfy the global Lipschitz conditions:
\begin{align}
    |b(t, x) - b(t, y)| &\leq K_2 |x - y| \notag \\
    |a^j(t, x) - a^j(t, y)| &\leq K_2 |x - y| \label{eq:MilsteinLipschitz} \\
    |\scL^{j_1} a^{j_2}(t, x) - \scL^{j_1} a^{j_2}(t, y)| &\leq K_2 |x - y| \notag
\end{align}
and the linear growth conditions:
\begin{align}
    |b(t, x)| + |\scL^j b(t, x)| &\leq K_3 (1 + |x|) \notag \\
    |a^j(t, x)| + |\scL^j a^{j_2}(t, x)| &\leq K_3 (1 + |x|) \label{eq:MilsteinLinearGrowth} \\
    |\scL^{j_1} \scL^{j_2} a^{j_3}(t, x)| &\leq K_3 (1 + |x|) \notag
\end{align}
Finally, assume the coefficients satisfy the H\"older continuity condition in time:
\begin{align}
    |b(s, x) - b(t, x)| &\leq K_4 (1 + |x|) |s - t|^{1/2} \notag \\
    |a^j(s, x) - a^j(t, x)| &\leq K_4 (1 + |x|) |s - t|^{1/2} \label{eq:MilsteinHolder} \\
    |\scL^{j_1} a^{j_2}(s, x) - \scL^{j_1} a^{j_2}(t, x)| &\leq K_4 (1 + |x|) |s - t|^{1/2} \notag
\end{align}
for all $s, t \in [0, T]$, $x, y \in \R^d$, and indices $j, j_1, j_2, j_3 = 1, \dots, m$. The constants $K_1, \dots, K_4$ must not depend on $\delta$.

\begin{thm}{Strong Convergence of the Milstein Scheme}{StrongConvergenceMilstein}
    Under the assumptions above, for the Milstein approximation $Y^\delta$, the following strong error estimate holds:
    \begin{equation}
        \E \left(|X_T - Y^\delta(T)|\right) \leq K_5 \delta, \label{eq:MilsteinConvergenceBound}
    \end{equation}
    where the constant $K_5$ does not depend on $\delta$.
\end{thm}

\begin{Pcode}{milstein.py}
    T, N, R = 1.0, 2**8, 4
    dt = T / N
    mu, sigma, X0 = 1.0, 0.5, 1.0

    np.random.seed(4)
    dB = np.sqrt(dt) * np.random.randn(N)

    t = np.linspace(0, T, N + 1)
    B = np.r_[0, np.cumsum(dB)]
    Y = X0 * np.exp((mu - 0.5 * sigma**2) * t + sigma * B)

    X = np.r_[X0, X0 * np.cumprod(1 + mu * dt + sigma * dB + 0.5 * sigma**2 * (dB**2 - dt))]

    dB_c = dB.reshape(-1, R).sum(axis=1)
    dt_c = dt * R
    X_c = np.r_[X0, X0 * np.cumprod(1 + mu * dt_c + sigma * dB_c + 0.5 * sigma**2 * (dB_c**2 - dt_c))]
\end{Pcode}

\begin{figure}[H]
    \centering
    %% Creator: Matplotlib, PGF backend
%%
%% To include the figure in your LaTeX document, write
%%   \input{<filename>.pgf}
%%
%% Make sure the required packages are loaded in your preamble
%%   \usepackage{pgf}
%%
%% Also ensure that all the required font packages are loaded; for instance,
%% the lmodern package is sometimes necessary when using math font.
%%   \usepackage{lmodern}
%%
%% Figures using additional raster images can only be included by \input if
%% they are in the same directory as the main LaTeX file. For loading figures
%% from other directories you can use the `import` package
%%   \usepackage{import}
%%
%% and then include the figures with
%%   \import{<path to file>}{<filename>.pgf}
%%
%% Matplotlib used the following preamble
%%   \def\mathdefault#1{#1}
%%   \everymath=\expandafter{\the\everymath\displaystyle}
%%   \IfFileExists{scrextend.sty}{
%%     \usepackage[fontsize=10.000000pt]{scrextend}
%%   }{
%%     \renewcommand{\normalsize}{\fontsize{10.000000}{12.000000}\selectfont}
%%     \normalsize
%%   }
%%   \usepackage{amsmath}
%%   \usepackage{amssymb}
%%   \usepackage{amsfonts}
%%   \usepackage{libertine}
%%   \usepackage[libertine]{newtxmath}
%%   \makeatletter\@ifpackageloaded{underscore}{}{\usepackage[strings]{underscore}}\makeatother
%%
\begingroup%
\makeatletter%
\begin{pgfpicture}%
\pgfpathrectangle{\pgfpointorigin}{\pgfqpoint{6.291307in}{4.038523in}}%
\pgfusepath{use as bounding box, clip}%
\begin{pgfscope}%
\pgfsetbuttcap%
\pgfsetmiterjoin%
\definecolor{currentfill}{rgb}{1.000000,1.000000,1.000000}%
\pgfsetfillcolor{currentfill}%
\pgfsetlinewidth{0.000000pt}%
\definecolor{currentstroke}{rgb}{1.000000,1.000000,1.000000}%
\pgfsetstrokecolor{currentstroke}%
\pgfsetdash{}{0pt}%
\pgfpathmoveto{\pgfqpoint{0.000000in}{0.000000in}}%
\pgfpathlineto{\pgfqpoint{6.291307in}{0.000000in}}%
\pgfpathlineto{\pgfqpoint{6.291307in}{4.038523in}}%
\pgfpathlineto{\pgfqpoint{0.000000in}{4.038523in}}%
\pgfpathlineto{\pgfqpoint{0.000000in}{0.000000in}}%
\pgfpathclose%
\pgfusepath{fill}%
\end{pgfscope}%
\begin{pgfscope}%
\pgfsetbuttcap%
\pgfsetmiterjoin%
\definecolor{currentfill}{rgb}{1.000000,1.000000,1.000000}%
\pgfsetfillcolor{currentfill}%
\pgfsetlinewidth{0.000000pt}%
\definecolor{currentstroke}{rgb}{0.000000,0.000000,0.000000}%
\pgfsetstrokecolor{currentstroke}%
\pgfsetstrokeopacity{0.000000}%
\pgfsetdash{}{0pt}%
\pgfpathmoveto{\pgfqpoint{0.413090in}{0.413680in}}%
\pgfpathlineto{\pgfqpoint{6.159154in}{0.413680in}}%
\pgfpathlineto{\pgfqpoint{6.159154in}{3.988523in}}%
\pgfpathlineto{\pgfqpoint{0.413090in}{3.988523in}}%
\pgfpathlineto{\pgfqpoint{0.413090in}{0.413680in}}%
\pgfpathclose%
\pgfusepath{fill}%
\end{pgfscope}%
\begin{pgfscope}%
\pgfsetbuttcap%
\pgfsetroundjoin%
\definecolor{currentfill}{rgb}{0.000000,0.000000,0.000000}%
\pgfsetfillcolor{currentfill}%
\pgfsetlinewidth{0.501875pt}%
\definecolor{currentstroke}{rgb}{0.000000,0.000000,0.000000}%
\pgfsetstrokecolor{currentstroke}%
\pgfsetdash{}{0pt}%
\pgfsys@defobject{currentmarker}{\pgfqpoint{0.000000in}{0.000000in}}{\pgfqpoint{0.000000in}{0.041667in}}{%
\pgfpathmoveto{\pgfqpoint{0.000000in}{0.000000in}}%
\pgfpathlineto{\pgfqpoint{0.000000in}{0.041667in}}%
\pgfusepath{stroke,fill}%
}%
\begin{pgfscope}%
\pgfsys@transformshift{0.413090in}{0.413680in}%
\pgfsys@useobject{currentmarker}{}%
\end{pgfscope}%
\end{pgfscope}%
\begin{pgfscope}%
\pgfsetbuttcap%
\pgfsetroundjoin%
\definecolor{currentfill}{rgb}{0.000000,0.000000,0.000000}%
\pgfsetfillcolor{currentfill}%
\pgfsetlinewidth{0.501875pt}%
\definecolor{currentstroke}{rgb}{0.000000,0.000000,0.000000}%
\pgfsetstrokecolor{currentstroke}%
\pgfsetdash{}{0pt}%
\pgfsys@defobject{currentmarker}{\pgfqpoint{0.000000in}{-0.041667in}}{\pgfqpoint{0.000000in}{0.000000in}}{%
\pgfpathmoveto{\pgfqpoint{0.000000in}{0.000000in}}%
\pgfpathlineto{\pgfqpoint{0.000000in}{-0.041667in}}%
\pgfusepath{stroke,fill}%
}%
\begin{pgfscope}%
\pgfsys@transformshift{0.413090in}{3.988523in}%
\pgfsys@useobject{currentmarker}{}%
\end{pgfscope}%
\end{pgfscope}%
\begin{pgfscope}%
\definecolor{textcolor}{rgb}{0.000000,0.000000,0.000000}%
\pgfsetstrokecolor{textcolor}%
\pgfsetfillcolor{textcolor}%
\pgftext[x=0.413090in,y=0.365069in,,top]{\color{textcolor}{\rmfamily\fontsize{10.000000}{12.000000}\selectfont\catcode`\^=\active\def^{\ifmmode\sp\else\^{}\fi}\catcode`\%=\active\def%{\%}$\mathdefault{0.0}$}}%
\end{pgfscope}%
\begin{pgfscope}%
\pgfsetbuttcap%
\pgfsetroundjoin%
\definecolor{currentfill}{rgb}{0.000000,0.000000,0.000000}%
\pgfsetfillcolor{currentfill}%
\pgfsetlinewidth{0.501875pt}%
\definecolor{currentstroke}{rgb}{0.000000,0.000000,0.000000}%
\pgfsetstrokecolor{currentstroke}%
\pgfsetdash{}{0pt}%
\pgfsys@defobject{currentmarker}{\pgfqpoint{0.000000in}{0.000000in}}{\pgfqpoint{0.000000in}{0.041667in}}{%
\pgfpathmoveto{\pgfqpoint{0.000000in}{0.000000in}}%
\pgfpathlineto{\pgfqpoint{0.000000in}{0.041667in}}%
\pgfusepath{stroke,fill}%
}%
\begin{pgfscope}%
\pgfsys@transformshift{1.562303in}{0.413680in}%
\pgfsys@useobject{currentmarker}{}%
\end{pgfscope}%
\end{pgfscope}%
\begin{pgfscope}%
\pgfsetbuttcap%
\pgfsetroundjoin%
\definecolor{currentfill}{rgb}{0.000000,0.000000,0.000000}%
\pgfsetfillcolor{currentfill}%
\pgfsetlinewidth{0.501875pt}%
\definecolor{currentstroke}{rgb}{0.000000,0.000000,0.000000}%
\pgfsetstrokecolor{currentstroke}%
\pgfsetdash{}{0pt}%
\pgfsys@defobject{currentmarker}{\pgfqpoint{0.000000in}{-0.041667in}}{\pgfqpoint{0.000000in}{0.000000in}}{%
\pgfpathmoveto{\pgfqpoint{0.000000in}{0.000000in}}%
\pgfpathlineto{\pgfqpoint{0.000000in}{-0.041667in}}%
\pgfusepath{stroke,fill}%
}%
\begin{pgfscope}%
\pgfsys@transformshift{1.562303in}{3.988523in}%
\pgfsys@useobject{currentmarker}{}%
\end{pgfscope}%
\end{pgfscope}%
\begin{pgfscope}%
\definecolor{textcolor}{rgb}{0.000000,0.000000,0.000000}%
\pgfsetstrokecolor{textcolor}%
\pgfsetfillcolor{textcolor}%
\pgftext[x=1.562303in,y=0.365069in,,top]{\color{textcolor}{\rmfamily\fontsize{10.000000}{12.000000}\selectfont\catcode`\^=\active\def^{\ifmmode\sp\else\^{}\fi}\catcode`\%=\active\def%{\%}$\mathdefault{0.2}$}}%
\end{pgfscope}%
\begin{pgfscope}%
\pgfsetbuttcap%
\pgfsetroundjoin%
\definecolor{currentfill}{rgb}{0.000000,0.000000,0.000000}%
\pgfsetfillcolor{currentfill}%
\pgfsetlinewidth{0.501875pt}%
\definecolor{currentstroke}{rgb}{0.000000,0.000000,0.000000}%
\pgfsetstrokecolor{currentstroke}%
\pgfsetdash{}{0pt}%
\pgfsys@defobject{currentmarker}{\pgfqpoint{0.000000in}{0.000000in}}{\pgfqpoint{0.000000in}{0.041667in}}{%
\pgfpathmoveto{\pgfqpoint{0.000000in}{0.000000in}}%
\pgfpathlineto{\pgfqpoint{0.000000in}{0.041667in}}%
\pgfusepath{stroke,fill}%
}%
\begin{pgfscope}%
\pgfsys@transformshift{2.711516in}{0.413680in}%
\pgfsys@useobject{currentmarker}{}%
\end{pgfscope}%
\end{pgfscope}%
\begin{pgfscope}%
\pgfsetbuttcap%
\pgfsetroundjoin%
\definecolor{currentfill}{rgb}{0.000000,0.000000,0.000000}%
\pgfsetfillcolor{currentfill}%
\pgfsetlinewidth{0.501875pt}%
\definecolor{currentstroke}{rgb}{0.000000,0.000000,0.000000}%
\pgfsetstrokecolor{currentstroke}%
\pgfsetdash{}{0pt}%
\pgfsys@defobject{currentmarker}{\pgfqpoint{0.000000in}{-0.041667in}}{\pgfqpoint{0.000000in}{0.000000in}}{%
\pgfpathmoveto{\pgfqpoint{0.000000in}{0.000000in}}%
\pgfpathlineto{\pgfqpoint{0.000000in}{-0.041667in}}%
\pgfusepath{stroke,fill}%
}%
\begin{pgfscope}%
\pgfsys@transformshift{2.711516in}{3.988523in}%
\pgfsys@useobject{currentmarker}{}%
\end{pgfscope}%
\end{pgfscope}%
\begin{pgfscope}%
\definecolor{textcolor}{rgb}{0.000000,0.000000,0.000000}%
\pgfsetstrokecolor{textcolor}%
\pgfsetfillcolor{textcolor}%
\pgftext[x=2.711516in,y=0.365069in,,top]{\color{textcolor}{\rmfamily\fontsize{10.000000}{12.000000}\selectfont\catcode`\^=\active\def^{\ifmmode\sp\else\^{}\fi}\catcode`\%=\active\def%{\%}$\mathdefault{0.4}$}}%
\end{pgfscope}%
\begin{pgfscope}%
\pgfsetbuttcap%
\pgfsetroundjoin%
\definecolor{currentfill}{rgb}{0.000000,0.000000,0.000000}%
\pgfsetfillcolor{currentfill}%
\pgfsetlinewidth{0.501875pt}%
\definecolor{currentstroke}{rgb}{0.000000,0.000000,0.000000}%
\pgfsetstrokecolor{currentstroke}%
\pgfsetdash{}{0pt}%
\pgfsys@defobject{currentmarker}{\pgfqpoint{0.000000in}{0.000000in}}{\pgfqpoint{0.000000in}{0.041667in}}{%
\pgfpathmoveto{\pgfqpoint{0.000000in}{0.000000in}}%
\pgfpathlineto{\pgfqpoint{0.000000in}{0.041667in}}%
\pgfusepath{stroke,fill}%
}%
\begin{pgfscope}%
\pgfsys@transformshift{3.860728in}{0.413680in}%
\pgfsys@useobject{currentmarker}{}%
\end{pgfscope}%
\end{pgfscope}%
\begin{pgfscope}%
\pgfsetbuttcap%
\pgfsetroundjoin%
\definecolor{currentfill}{rgb}{0.000000,0.000000,0.000000}%
\pgfsetfillcolor{currentfill}%
\pgfsetlinewidth{0.501875pt}%
\definecolor{currentstroke}{rgb}{0.000000,0.000000,0.000000}%
\pgfsetstrokecolor{currentstroke}%
\pgfsetdash{}{0pt}%
\pgfsys@defobject{currentmarker}{\pgfqpoint{0.000000in}{-0.041667in}}{\pgfqpoint{0.000000in}{0.000000in}}{%
\pgfpathmoveto{\pgfqpoint{0.000000in}{0.000000in}}%
\pgfpathlineto{\pgfqpoint{0.000000in}{-0.041667in}}%
\pgfusepath{stroke,fill}%
}%
\begin{pgfscope}%
\pgfsys@transformshift{3.860728in}{3.988523in}%
\pgfsys@useobject{currentmarker}{}%
\end{pgfscope}%
\end{pgfscope}%
\begin{pgfscope}%
\definecolor{textcolor}{rgb}{0.000000,0.000000,0.000000}%
\pgfsetstrokecolor{textcolor}%
\pgfsetfillcolor{textcolor}%
\pgftext[x=3.860728in,y=0.365069in,,top]{\color{textcolor}{\rmfamily\fontsize{10.000000}{12.000000}\selectfont\catcode`\^=\active\def^{\ifmmode\sp\else\^{}\fi}\catcode`\%=\active\def%{\%}$\mathdefault{0.6}$}}%
\end{pgfscope}%
\begin{pgfscope}%
\pgfsetbuttcap%
\pgfsetroundjoin%
\definecolor{currentfill}{rgb}{0.000000,0.000000,0.000000}%
\pgfsetfillcolor{currentfill}%
\pgfsetlinewidth{0.501875pt}%
\definecolor{currentstroke}{rgb}{0.000000,0.000000,0.000000}%
\pgfsetstrokecolor{currentstroke}%
\pgfsetdash{}{0pt}%
\pgfsys@defobject{currentmarker}{\pgfqpoint{0.000000in}{0.000000in}}{\pgfqpoint{0.000000in}{0.041667in}}{%
\pgfpathmoveto{\pgfqpoint{0.000000in}{0.000000in}}%
\pgfpathlineto{\pgfqpoint{0.000000in}{0.041667in}}%
\pgfusepath{stroke,fill}%
}%
\begin{pgfscope}%
\pgfsys@transformshift{5.009941in}{0.413680in}%
\pgfsys@useobject{currentmarker}{}%
\end{pgfscope}%
\end{pgfscope}%
\begin{pgfscope}%
\pgfsetbuttcap%
\pgfsetroundjoin%
\definecolor{currentfill}{rgb}{0.000000,0.000000,0.000000}%
\pgfsetfillcolor{currentfill}%
\pgfsetlinewidth{0.501875pt}%
\definecolor{currentstroke}{rgb}{0.000000,0.000000,0.000000}%
\pgfsetstrokecolor{currentstroke}%
\pgfsetdash{}{0pt}%
\pgfsys@defobject{currentmarker}{\pgfqpoint{0.000000in}{-0.041667in}}{\pgfqpoint{0.000000in}{0.000000in}}{%
\pgfpathmoveto{\pgfqpoint{0.000000in}{0.000000in}}%
\pgfpathlineto{\pgfqpoint{0.000000in}{-0.041667in}}%
\pgfusepath{stroke,fill}%
}%
\begin{pgfscope}%
\pgfsys@transformshift{5.009941in}{3.988523in}%
\pgfsys@useobject{currentmarker}{}%
\end{pgfscope}%
\end{pgfscope}%
\begin{pgfscope}%
\definecolor{textcolor}{rgb}{0.000000,0.000000,0.000000}%
\pgfsetstrokecolor{textcolor}%
\pgfsetfillcolor{textcolor}%
\pgftext[x=5.009941in,y=0.365069in,,top]{\color{textcolor}{\rmfamily\fontsize{10.000000}{12.000000}\selectfont\catcode`\^=\active\def^{\ifmmode\sp\else\^{}\fi}\catcode`\%=\active\def%{\%}$\mathdefault{0.8}$}}%
\end{pgfscope}%
\begin{pgfscope}%
\pgfsetbuttcap%
\pgfsetroundjoin%
\definecolor{currentfill}{rgb}{0.000000,0.000000,0.000000}%
\pgfsetfillcolor{currentfill}%
\pgfsetlinewidth{0.501875pt}%
\definecolor{currentstroke}{rgb}{0.000000,0.000000,0.000000}%
\pgfsetstrokecolor{currentstroke}%
\pgfsetdash{}{0pt}%
\pgfsys@defobject{currentmarker}{\pgfqpoint{0.000000in}{0.000000in}}{\pgfqpoint{0.000000in}{0.041667in}}{%
\pgfpathmoveto{\pgfqpoint{0.000000in}{0.000000in}}%
\pgfpathlineto{\pgfqpoint{0.000000in}{0.041667in}}%
\pgfusepath{stroke,fill}%
}%
\begin{pgfscope}%
\pgfsys@transformshift{6.159154in}{0.413680in}%
\pgfsys@useobject{currentmarker}{}%
\end{pgfscope}%
\end{pgfscope}%
\begin{pgfscope}%
\pgfsetbuttcap%
\pgfsetroundjoin%
\definecolor{currentfill}{rgb}{0.000000,0.000000,0.000000}%
\pgfsetfillcolor{currentfill}%
\pgfsetlinewidth{0.501875pt}%
\definecolor{currentstroke}{rgb}{0.000000,0.000000,0.000000}%
\pgfsetstrokecolor{currentstroke}%
\pgfsetdash{}{0pt}%
\pgfsys@defobject{currentmarker}{\pgfqpoint{0.000000in}{-0.041667in}}{\pgfqpoint{0.000000in}{0.000000in}}{%
\pgfpathmoveto{\pgfqpoint{0.000000in}{0.000000in}}%
\pgfpathlineto{\pgfqpoint{0.000000in}{-0.041667in}}%
\pgfusepath{stroke,fill}%
}%
\begin{pgfscope}%
\pgfsys@transformshift{6.159154in}{3.988523in}%
\pgfsys@useobject{currentmarker}{}%
\end{pgfscope}%
\end{pgfscope}%
\begin{pgfscope}%
\definecolor{textcolor}{rgb}{0.000000,0.000000,0.000000}%
\pgfsetstrokecolor{textcolor}%
\pgfsetfillcolor{textcolor}%
\pgftext[x=6.159154in,y=0.365069in,,top]{\color{textcolor}{\rmfamily\fontsize{10.000000}{12.000000}\selectfont\catcode`\^=\active\def^{\ifmmode\sp\else\^{}\fi}\catcode`\%=\active\def%{\%}$\mathdefault{1.0}$}}%
\end{pgfscope}%
\begin{pgfscope}%
\pgfsetbuttcap%
\pgfsetroundjoin%
\definecolor{currentfill}{rgb}{0.000000,0.000000,0.000000}%
\pgfsetfillcolor{currentfill}%
\pgfsetlinewidth{0.501875pt}%
\definecolor{currentstroke}{rgb}{0.000000,0.000000,0.000000}%
\pgfsetstrokecolor{currentstroke}%
\pgfsetdash{}{0pt}%
\pgfsys@defobject{currentmarker}{\pgfqpoint{0.000000in}{0.000000in}}{\pgfqpoint{0.000000in}{0.020833in}}{%
\pgfpathmoveto{\pgfqpoint{0.000000in}{0.000000in}}%
\pgfpathlineto{\pgfqpoint{0.000000in}{0.020833in}}%
\pgfusepath{stroke,fill}%
}%
\begin{pgfscope}%
\pgfsys@transformshift{0.700393in}{0.413680in}%
\pgfsys@useobject{currentmarker}{}%
\end{pgfscope}%
\end{pgfscope}%
\begin{pgfscope}%
\pgfsetbuttcap%
\pgfsetroundjoin%
\definecolor{currentfill}{rgb}{0.000000,0.000000,0.000000}%
\pgfsetfillcolor{currentfill}%
\pgfsetlinewidth{0.501875pt}%
\definecolor{currentstroke}{rgb}{0.000000,0.000000,0.000000}%
\pgfsetstrokecolor{currentstroke}%
\pgfsetdash{}{0pt}%
\pgfsys@defobject{currentmarker}{\pgfqpoint{0.000000in}{-0.020833in}}{\pgfqpoint{0.000000in}{0.000000in}}{%
\pgfpathmoveto{\pgfqpoint{0.000000in}{0.000000in}}%
\pgfpathlineto{\pgfqpoint{0.000000in}{-0.020833in}}%
\pgfusepath{stroke,fill}%
}%
\begin{pgfscope}%
\pgfsys@transformshift{0.700393in}{3.988523in}%
\pgfsys@useobject{currentmarker}{}%
\end{pgfscope}%
\end{pgfscope}%
\begin{pgfscope}%
\pgfsetbuttcap%
\pgfsetroundjoin%
\definecolor{currentfill}{rgb}{0.000000,0.000000,0.000000}%
\pgfsetfillcolor{currentfill}%
\pgfsetlinewidth{0.501875pt}%
\definecolor{currentstroke}{rgb}{0.000000,0.000000,0.000000}%
\pgfsetstrokecolor{currentstroke}%
\pgfsetdash{}{0pt}%
\pgfsys@defobject{currentmarker}{\pgfqpoint{0.000000in}{0.000000in}}{\pgfqpoint{0.000000in}{0.020833in}}{%
\pgfpathmoveto{\pgfqpoint{0.000000in}{0.000000in}}%
\pgfpathlineto{\pgfqpoint{0.000000in}{0.020833in}}%
\pgfusepath{stroke,fill}%
}%
\begin{pgfscope}%
\pgfsys@transformshift{0.987696in}{0.413680in}%
\pgfsys@useobject{currentmarker}{}%
\end{pgfscope}%
\end{pgfscope}%
\begin{pgfscope}%
\pgfsetbuttcap%
\pgfsetroundjoin%
\definecolor{currentfill}{rgb}{0.000000,0.000000,0.000000}%
\pgfsetfillcolor{currentfill}%
\pgfsetlinewidth{0.501875pt}%
\definecolor{currentstroke}{rgb}{0.000000,0.000000,0.000000}%
\pgfsetstrokecolor{currentstroke}%
\pgfsetdash{}{0pt}%
\pgfsys@defobject{currentmarker}{\pgfqpoint{0.000000in}{-0.020833in}}{\pgfqpoint{0.000000in}{0.000000in}}{%
\pgfpathmoveto{\pgfqpoint{0.000000in}{0.000000in}}%
\pgfpathlineto{\pgfqpoint{0.000000in}{-0.020833in}}%
\pgfusepath{stroke,fill}%
}%
\begin{pgfscope}%
\pgfsys@transformshift{0.987696in}{3.988523in}%
\pgfsys@useobject{currentmarker}{}%
\end{pgfscope}%
\end{pgfscope}%
\begin{pgfscope}%
\pgfsetbuttcap%
\pgfsetroundjoin%
\definecolor{currentfill}{rgb}{0.000000,0.000000,0.000000}%
\pgfsetfillcolor{currentfill}%
\pgfsetlinewidth{0.501875pt}%
\definecolor{currentstroke}{rgb}{0.000000,0.000000,0.000000}%
\pgfsetstrokecolor{currentstroke}%
\pgfsetdash{}{0pt}%
\pgfsys@defobject{currentmarker}{\pgfqpoint{0.000000in}{0.000000in}}{\pgfqpoint{0.000000in}{0.020833in}}{%
\pgfpathmoveto{\pgfqpoint{0.000000in}{0.000000in}}%
\pgfpathlineto{\pgfqpoint{0.000000in}{0.020833in}}%
\pgfusepath{stroke,fill}%
}%
\begin{pgfscope}%
\pgfsys@transformshift{1.275000in}{0.413680in}%
\pgfsys@useobject{currentmarker}{}%
\end{pgfscope}%
\end{pgfscope}%
\begin{pgfscope}%
\pgfsetbuttcap%
\pgfsetroundjoin%
\definecolor{currentfill}{rgb}{0.000000,0.000000,0.000000}%
\pgfsetfillcolor{currentfill}%
\pgfsetlinewidth{0.501875pt}%
\definecolor{currentstroke}{rgb}{0.000000,0.000000,0.000000}%
\pgfsetstrokecolor{currentstroke}%
\pgfsetdash{}{0pt}%
\pgfsys@defobject{currentmarker}{\pgfqpoint{0.000000in}{-0.020833in}}{\pgfqpoint{0.000000in}{0.000000in}}{%
\pgfpathmoveto{\pgfqpoint{0.000000in}{0.000000in}}%
\pgfpathlineto{\pgfqpoint{0.000000in}{-0.020833in}}%
\pgfusepath{stroke,fill}%
}%
\begin{pgfscope}%
\pgfsys@transformshift{1.275000in}{3.988523in}%
\pgfsys@useobject{currentmarker}{}%
\end{pgfscope}%
\end{pgfscope}%
\begin{pgfscope}%
\pgfsetbuttcap%
\pgfsetroundjoin%
\definecolor{currentfill}{rgb}{0.000000,0.000000,0.000000}%
\pgfsetfillcolor{currentfill}%
\pgfsetlinewidth{0.501875pt}%
\definecolor{currentstroke}{rgb}{0.000000,0.000000,0.000000}%
\pgfsetstrokecolor{currentstroke}%
\pgfsetdash{}{0pt}%
\pgfsys@defobject{currentmarker}{\pgfqpoint{0.000000in}{0.000000in}}{\pgfqpoint{0.000000in}{0.020833in}}{%
\pgfpathmoveto{\pgfqpoint{0.000000in}{0.000000in}}%
\pgfpathlineto{\pgfqpoint{0.000000in}{0.020833in}}%
\pgfusepath{stroke,fill}%
}%
\begin{pgfscope}%
\pgfsys@transformshift{1.849606in}{0.413680in}%
\pgfsys@useobject{currentmarker}{}%
\end{pgfscope}%
\end{pgfscope}%
\begin{pgfscope}%
\pgfsetbuttcap%
\pgfsetroundjoin%
\definecolor{currentfill}{rgb}{0.000000,0.000000,0.000000}%
\pgfsetfillcolor{currentfill}%
\pgfsetlinewidth{0.501875pt}%
\definecolor{currentstroke}{rgb}{0.000000,0.000000,0.000000}%
\pgfsetstrokecolor{currentstroke}%
\pgfsetdash{}{0pt}%
\pgfsys@defobject{currentmarker}{\pgfqpoint{0.000000in}{-0.020833in}}{\pgfqpoint{0.000000in}{0.000000in}}{%
\pgfpathmoveto{\pgfqpoint{0.000000in}{0.000000in}}%
\pgfpathlineto{\pgfqpoint{0.000000in}{-0.020833in}}%
\pgfusepath{stroke,fill}%
}%
\begin{pgfscope}%
\pgfsys@transformshift{1.849606in}{3.988523in}%
\pgfsys@useobject{currentmarker}{}%
\end{pgfscope}%
\end{pgfscope}%
\begin{pgfscope}%
\pgfsetbuttcap%
\pgfsetroundjoin%
\definecolor{currentfill}{rgb}{0.000000,0.000000,0.000000}%
\pgfsetfillcolor{currentfill}%
\pgfsetlinewidth{0.501875pt}%
\definecolor{currentstroke}{rgb}{0.000000,0.000000,0.000000}%
\pgfsetstrokecolor{currentstroke}%
\pgfsetdash{}{0pt}%
\pgfsys@defobject{currentmarker}{\pgfqpoint{0.000000in}{0.000000in}}{\pgfqpoint{0.000000in}{0.020833in}}{%
\pgfpathmoveto{\pgfqpoint{0.000000in}{0.000000in}}%
\pgfpathlineto{\pgfqpoint{0.000000in}{0.020833in}}%
\pgfusepath{stroke,fill}%
}%
\begin{pgfscope}%
\pgfsys@transformshift{2.136909in}{0.413680in}%
\pgfsys@useobject{currentmarker}{}%
\end{pgfscope}%
\end{pgfscope}%
\begin{pgfscope}%
\pgfsetbuttcap%
\pgfsetroundjoin%
\definecolor{currentfill}{rgb}{0.000000,0.000000,0.000000}%
\pgfsetfillcolor{currentfill}%
\pgfsetlinewidth{0.501875pt}%
\definecolor{currentstroke}{rgb}{0.000000,0.000000,0.000000}%
\pgfsetstrokecolor{currentstroke}%
\pgfsetdash{}{0pt}%
\pgfsys@defobject{currentmarker}{\pgfqpoint{0.000000in}{-0.020833in}}{\pgfqpoint{0.000000in}{0.000000in}}{%
\pgfpathmoveto{\pgfqpoint{0.000000in}{0.000000in}}%
\pgfpathlineto{\pgfqpoint{0.000000in}{-0.020833in}}%
\pgfusepath{stroke,fill}%
}%
\begin{pgfscope}%
\pgfsys@transformshift{2.136909in}{3.988523in}%
\pgfsys@useobject{currentmarker}{}%
\end{pgfscope}%
\end{pgfscope}%
\begin{pgfscope}%
\pgfsetbuttcap%
\pgfsetroundjoin%
\definecolor{currentfill}{rgb}{0.000000,0.000000,0.000000}%
\pgfsetfillcolor{currentfill}%
\pgfsetlinewidth{0.501875pt}%
\definecolor{currentstroke}{rgb}{0.000000,0.000000,0.000000}%
\pgfsetstrokecolor{currentstroke}%
\pgfsetdash{}{0pt}%
\pgfsys@defobject{currentmarker}{\pgfqpoint{0.000000in}{0.000000in}}{\pgfqpoint{0.000000in}{0.020833in}}{%
\pgfpathmoveto{\pgfqpoint{0.000000in}{0.000000in}}%
\pgfpathlineto{\pgfqpoint{0.000000in}{0.020833in}}%
\pgfusepath{stroke,fill}%
}%
\begin{pgfscope}%
\pgfsys@transformshift{2.424212in}{0.413680in}%
\pgfsys@useobject{currentmarker}{}%
\end{pgfscope}%
\end{pgfscope}%
\begin{pgfscope}%
\pgfsetbuttcap%
\pgfsetroundjoin%
\definecolor{currentfill}{rgb}{0.000000,0.000000,0.000000}%
\pgfsetfillcolor{currentfill}%
\pgfsetlinewidth{0.501875pt}%
\definecolor{currentstroke}{rgb}{0.000000,0.000000,0.000000}%
\pgfsetstrokecolor{currentstroke}%
\pgfsetdash{}{0pt}%
\pgfsys@defobject{currentmarker}{\pgfqpoint{0.000000in}{-0.020833in}}{\pgfqpoint{0.000000in}{0.000000in}}{%
\pgfpathmoveto{\pgfqpoint{0.000000in}{0.000000in}}%
\pgfpathlineto{\pgfqpoint{0.000000in}{-0.020833in}}%
\pgfusepath{stroke,fill}%
}%
\begin{pgfscope}%
\pgfsys@transformshift{2.424212in}{3.988523in}%
\pgfsys@useobject{currentmarker}{}%
\end{pgfscope}%
\end{pgfscope}%
\begin{pgfscope}%
\pgfsetbuttcap%
\pgfsetroundjoin%
\definecolor{currentfill}{rgb}{0.000000,0.000000,0.000000}%
\pgfsetfillcolor{currentfill}%
\pgfsetlinewidth{0.501875pt}%
\definecolor{currentstroke}{rgb}{0.000000,0.000000,0.000000}%
\pgfsetstrokecolor{currentstroke}%
\pgfsetdash{}{0pt}%
\pgfsys@defobject{currentmarker}{\pgfqpoint{0.000000in}{0.000000in}}{\pgfqpoint{0.000000in}{0.020833in}}{%
\pgfpathmoveto{\pgfqpoint{0.000000in}{0.000000in}}%
\pgfpathlineto{\pgfqpoint{0.000000in}{0.020833in}}%
\pgfusepath{stroke,fill}%
}%
\begin{pgfscope}%
\pgfsys@transformshift{2.998819in}{0.413680in}%
\pgfsys@useobject{currentmarker}{}%
\end{pgfscope}%
\end{pgfscope}%
\begin{pgfscope}%
\pgfsetbuttcap%
\pgfsetroundjoin%
\definecolor{currentfill}{rgb}{0.000000,0.000000,0.000000}%
\pgfsetfillcolor{currentfill}%
\pgfsetlinewidth{0.501875pt}%
\definecolor{currentstroke}{rgb}{0.000000,0.000000,0.000000}%
\pgfsetstrokecolor{currentstroke}%
\pgfsetdash{}{0pt}%
\pgfsys@defobject{currentmarker}{\pgfqpoint{0.000000in}{-0.020833in}}{\pgfqpoint{0.000000in}{0.000000in}}{%
\pgfpathmoveto{\pgfqpoint{0.000000in}{0.000000in}}%
\pgfpathlineto{\pgfqpoint{0.000000in}{-0.020833in}}%
\pgfusepath{stroke,fill}%
}%
\begin{pgfscope}%
\pgfsys@transformshift{2.998819in}{3.988523in}%
\pgfsys@useobject{currentmarker}{}%
\end{pgfscope}%
\end{pgfscope}%
\begin{pgfscope}%
\pgfsetbuttcap%
\pgfsetroundjoin%
\definecolor{currentfill}{rgb}{0.000000,0.000000,0.000000}%
\pgfsetfillcolor{currentfill}%
\pgfsetlinewidth{0.501875pt}%
\definecolor{currentstroke}{rgb}{0.000000,0.000000,0.000000}%
\pgfsetstrokecolor{currentstroke}%
\pgfsetdash{}{0pt}%
\pgfsys@defobject{currentmarker}{\pgfqpoint{0.000000in}{0.000000in}}{\pgfqpoint{0.000000in}{0.020833in}}{%
\pgfpathmoveto{\pgfqpoint{0.000000in}{0.000000in}}%
\pgfpathlineto{\pgfqpoint{0.000000in}{0.020833in}}%
\pgfusepath{stroke,fill}%
}%
\begin{pgfscope}%
\pgfsys@transformshift{3.286122in}{0.413680in}%
\pgfsys@useobject{currentmarker}{}%
\end{pgfscope}%
\end{pgfscope}%
\begin{pgfscope}%
\pgfsetbuttcap%
\pgfsetroundjoin%
\definecolor{currentfill}{rgb}{0.000000,0.000000,0.000000}%
\pgfsetfillcolor{currentfill}%
\pgfsetlinewidth{0.501875pt}%
\definecolor{currentstroke}{rgb}{0.000000,0.000000,0.000000}%
\pgfsetstrokecolor{currentstroke}%
\pgfsetdash{}{0pt}%
\pgfsys@defobject{currentmarker}{\pgfqpoint{0.000000in}{-0.020833in}}{\pgfqpoint{0.000000in}{0.000000in}}{%
\pgfpathmoveto{\pgfqpoint{0.000000in}{0.000000in}}%
\pgfpathlineto{\pgfqpoint{0.000000in}{-0.020833in}}%
\pgfusepath{stroke,fill}%
}%
\begin{pgfscope}%
\pgfsys@transformshift{3.286122in}{3.988523in}%
\pgfsys@useobject{currentmarker}{}%
\end{pgfscope}%
\end{pgfscope}%
\begin{pgfscope}%
\pgfsetbuttcap%
\pgfsetroundjoin%
\definecolor{currentfill}{rgb}{0.000000,0.000000,0.000000}%
\pgfsetfillcolor{currentfill}%
\pgfsetlinewidth{0.501875pt}%
\definecolor{currentstroke}{rgb}{0.000000,0.000000,0.000000}%
\pgfsetstrokecolor{currentstroke}%
\pgfsetdash{}{0pt}%
\pgfsys@defobject{currentmarker}{\pgfqpoint{0.000000in}{0.000000in}}{\pgfqpoint{0.000000in}{0.020833in}}{%
\pgfpathmoveto{\pgfqpoint{0.000000in}{0.000000in}}%
\pgfpathlineto{\pgfqpoint{0.000000in}{0.020833in}}%
\pgfusepath{stroke,fill}%
}%
\begin{pgfscope}%
\pgfsys@transformshift{3.573425in}{0.413680in}%
\pgfsys@useobject{currentmarker}{}%
\end{pgfscope}%
\end{pgfscope}%
\begin{pgfscope}%
\pgfsetbuttcap%
\pgfsetroundjoin%
\definecolor{currentfill}{rgb}{0.000000,0.000000,0.000000}%
\pgfsetfillcolor{currentfill}%
\pgfsetlinewidth{0.501875pt}%
\definecolor{currentstroke}{rgb}{0.000000,0.000000,0.000000}%
\pgfsetstrokecolor{currentstroke}%
\pgfsetdash{}{0pt}%
\pgfsys@defobject{currentmarker}{\pgfqpoint{0.000000in}{-0.020833in}}{\pgfqpoint{0.000000in}{0.000000in}}{%
\pgfpathmoveto{\pgfqpoint{0.000000in}{0.000000in}}%
\pgfpathlineto{\pgfqpoint{0.000000in}{-0.020833in}}%
\pgfusepath{stroke,fill}%
}%
\begin{pgfscope}%
\pgfsys@transformshift{3.573425in}{3.988523in}%
\pgfsys@useobject{currentmarker}{}%
\end{pgfscope}%
\end{pgfscope}%
\begin{pgfscope}%
\pgfsetbuttcap%
\pgfsetroundjoin%
\definecolor{currentfill}{rgb}{0.000000,0.000000,0.000000}%
\pgfsetfillcolor{currentfill}%
\pgfsetlinewidth{0.501875pt}%
\definecolor{currentstroke}{rgb}{0.000000,0.000000,0.000000}%
\pgfsetstrokecolor{currentstroke}%
\pgfsetdash{}{0pt}%
\pgfsys@defobject{currentmarker}{\pgfqpoint{0.000000in}{0.000000in}}{\pgfqpoint{0.000000in}{0.020833in}}{%
\pgfpathmoveto{\pgfqpoint{0.000000in}{0.000000in}}%
\pgfpathlineto{\pgfqpoint{0.000000in}{0.020833in}}%
\pgfusepath{stroke,fill}%
}%
\begin{pgfscope}%
\pgfsys@transformshift{4.148032in}{0.413680in}%
\pgfsys@useobject{currentmarker}{}%
\end{pgfscope}%
\end{pgfscope}%
\begin{pgfscope}%
\pgfsetbuttcap%
\pgfsetroundjoin%
\definecolor{currentfill}{rgb}{0.000000,0.000000,0.000000}%
\pgfsetfillcolor{currentfill}%
\pgfsetlinewidth{0.501875pt}%
\definecolor{currentstroke}{rgb}{0.000000,0.000000,0.000000}%
\pgfsetstrokecolor{currentstroke}%
\pgfsetdash{}{0pt}%
\pgfsys@defobject{currentmarker}{\pgfqpoint{0.000000in}{-0.020833in}}{\pgfqpoint{0.000000in}{0.000000in}}{%
\pgfpathmoveto{\pgfqpoint{0.000000in}{0.000000in}}%
\pgfpathlineto{\pgfqpoint{0.000000in}{-0.020833in}}%
\pgfusepath{stroke,fill}%
}%
\begin{pgfscope}%
\pgfsys@transformshift{4.148032in}{3.988523in}%
\pgfsys@useobject{currentmarker}{}%
\end{pgfscope}%
\end{pgfscope}%
\begin{pgfscope}%
\pgfsetbuttcap%
\pgfsetroundjoin%
\definecolor{currentfill}{rgb}{0.000000,0.000000,0.000000}%
\pgfsetfillcolor{currentfill}%
\pgfsetlinewidth{0.501875pt}%
\definecolor{currentstroke}{rgb}{0.000000,0.000000,0.000000}%
\pgfsetstrokecolor{currentstroke}%
\pgfsetdash{}{0pt}%
\pgfsys@defobject{currentmarker}{\pgfqpoint{0.000000in}{0.000000in}}{\pgfqpoint{0.000000in}{0.020833in}}{%
\pgfpathmoveto{\pgfqpoint{0.000000in}{0.000000in}}%
\pgfpathlineto{\pgfqpoint{0.000000in}{0.020833in}}%
\pgfusepath{stroke,fill}%
}%
\begin{pgfscope}%
\pgfsys@transformshift{4.435335in}{0.413680in}%
\pgfsys@useobject{currentmarker}{}%
\end{pgfscope}%
\end{pgfscope}%
\begin{pgfscope}%
\pgfsetbuttcap%
\pgfsetroundjoin%
\definecolor{currentfill}{rgb}{0.000000,0.000000,0.000000}%
\pgfsetfillcolor{currentfill}%
\pgfsetlinewidth{0.501875pt}%
\definecolor{currentstroke}{rgb}{0.000000,0.000000,0.000000}%
\pgfsetstrokecolor{currentstroke}%
\pgfsetdash{}{0pt}%
\pgfsys@defobject{currentmarker}{\pgfqpoint{0.000000in}{-0.020833in}}{\pgfqpoint{0.000000in}{0.000000in}}{%
\pgfpathmoveto{\pgfqpoint{0.000000in}{0.000000in}}%
\pgfpathlineto{\pgfqpoint{0.000000in}{-0.020833in}}%
\pgfusepath{stroke,fill}%
}%
\begin{pgfscope}%
\pgfsys@transformshift{4.435335in}{3.988523in}%
\pgfsys@useobject{currentmarker}{}%
\end{pgfscope}%
\end{pgfscope}%
\begin{pgfscope}%
\pgfsetbuttcap%
\pgfsetroundjoin%
\definecolor{currentfill}{rgb}{0.000000,0.000000,0.000000}%
\pgfsetfillcolor{currentfill}%
\pgfsetlinewidth{0.501875pt}%
\definecolor{currentstroke}{rgb}{0.000000,0.000000,0.000000}%
\pgfsetstrokecolor{currentstroke}%
\pgfsetdash{}{0pt}%
\pgfsys@defobject{currentmarker}{\pgfqpoint{0.000000in}{0.000000in}}{\pgfqpoint{0.000000in}{0.020833in}}{%
\pgfpathmoveto{\pgfqpoint{0.000000in}{0.000000in}}%
\pgfpathlineto{\pgfqpoint{0.000000in}{0.020833in}}%
\pgfusepath{stroke,fill}%
}%
\begin{pgfscope}%
\pgfsys@transformshift{4.722638in}{0.413680in}%
\pgfsys@useobject{currentmarker}{}%
\end{pgfscope}%
\end{pgfscope}%
\begin{pgfscope}%
\pgfsetbuttcap%
\pgfsetroundjoin%
\definecolor{currentfill}{rgb}{0.000000,0.000000,0.000000}%
\pgfsetfillcolor{currentfill}%
\pgfsetlinewidth{0.501875pt}%
\definecolor{currentstroke}{rgb}{0.000000,0.000000,0.000000}%
\pgfsetstrokecolor{currentstroke}%
\pgfsetdash{}{0pt}%
\pgfsys@defobject{currentmarker}{\pgfqpoint{0.000000in}{-0.020833in}}{\pgfqpoint{0.000000in}{0.000000in}}{%
\pgfpathmoveto{\pgfqpoint{0.000000in}{0.000000in}}%
\pgfpathlineto{\pgfqpoint{0.000000in}{-0.020833in}}%
\pgfusepath{stroke,fill}%
}%
\begin{pgfscope}%
\pgfsys@transformshift{4.722638in}{3.988523in}%
\pgfsys@useobject{currentmarker}{}%
\end{pgfscope}%
\end{pgfscope}%
\begin{pgfscope}%
\pgfsetbuttcap%
\pgfsetroundjoin%
\definecolor{currentfill}{rgb}{0.000000,0.000000,0.000000}%
\pgfsetfillcolor{currentfill}%
\pgfsetlinewidth{0.501875pt}%
\definecolor{currentstroke}{rgb}{0.000000,0.000000,0.000000}%
\pgfsetstrokecolor{currentstroke}%
\pgfsetdash{}{0pt}%
\pgfsys@defobject{currentmarker}{\pgfqpoint{0.000000in}{0.000000in}}{\pgfqpoint{0.000000in}{0.020833in}}{%
\pgfpathmoveto{\pgfqpoint{0.000000in}{0.000000in}}%
\pgfpathlineto{\pgfqpoint{0.000000in}{0.020833in}}%
\pgfusepath{stroke,fill}%
}%
\begin{pgfscope}%
\pgfsys@transformshift{5.297244in}{0.413680in}%
\pgfsys@useobject{currentmarker}{}%
\end{pgfscope}%
\end{pgfscope}%
\begin{pgfscope}%
\pgfsetbuttcap%
\pgfsetroundjoin%
\definecolor{currentfill}{rgb}{0.000000,0.000000,0.000000}%
\pgfsetfillcolor{currentfill}%
\pgfsetlinewidth{0.501875pt}%
\definecolor{currentstroke}{rgb}{0.000000,0.000000,0.000000}%
\pgfsetstrokecolor{currentstroke}%
\pgfsetdash{}{0pt}%
\pgfsys@defobject{currentmarker}{\pgfqpoint{0.000000in}{-0.020833in}}{\pgfqpoint{0.000000in}{0.000000in}}{%
\pgfpathmoveto{\pgfqpoint{0.000000in}{0.000000in}}%
\pgfpathlineto{\pgfqpoint{0.000000in}{-0.020833in}}%
\pgfusepath{stroke,fill}%
}%
\begin{pgfscope}%
\pgfsys@transformshift{5.297244in}{3.988523in}%
\pgfsys@useobject{currentmarker}{}%
\end{pgfscope}%
\end{pgfscope}%
\begin{pgfscope}%
\pgfsetbuttcap%
\pgfsetroundjoin%
\definecolor{currentfill}{rgb}{0.000000,0.000000,0.000000}%
\pgfsetfillcolor{currentfill}%
\pgfsetlinewidth{0.501875pt}%
\definecolor{currentstroke}{rgb}{0.000000,0.000000,0.000000}%
\pgfsetstrokecolor{currentstroke}%
\pgfsetdash{}{0pt}%
\pgfsys@defobject{currentmarker}{\pgfqpoint{0.000000in}{0.000000in}}{\pgfqpoint{0.000000in}{0.020833in}}{%
\pgfpathmoveto{\pgfqpoint{0.000000in}{0.000000in}}%
\pgfpathlineto{\pgfqpoint{0.000000in}{0.020833in}}%
\pgfusepath{stroke,fill}%
}%
\begin{pgfscope}%
\pgfsys@transformshift{5.584548in}{0.413680in}%
\pgfsys@useobject{currentmarker}{}%
\end{pgfscope}%
\end{pgfscope}%
\begin{pgfscope}%
\pgfsetbuttcap%
\pgfsetroundjoin%
\definecolor{currentfill}{rgb}{0.000000,0.000000,0.000000}%
\pgfsetfillcolor{currentfill}%
\pgfsetlinewidth{0.501875pt}%
\definecolor{currentstroke}{rgb}{0.000000,0.000000,0.000000}%
\pgfsetstrokecolor{currentstroke}%
\pgfsetdash{}{0pt}%
\pgfsys@defobject{currentmarker}{\pgfqpoint{0.000000in}{-0.020833in}}{\pgfqpoint{0.000000in}{0.000000in}}{%
\pgfpathmoveto{\pgfqpoint{0.000000in}{0.000000in}}%
\pgfpathlineto{\pgfqpoint{0.000000in}{-0.020833in}}%
\pgfusepath{stroke,fill}%
}%
\begin{pgfscope}%
\pgfsys@transformshift{5.584548in}{3.988523in}%
\pgfsys@useobject{currentmarker}{}%
\end{pgfscope}%
\end{pgfscope}%
\begin{pgfscope}%
\pgfsetbuttcap%
\pgfsetroundjoin%
\definecolor{currentfill}{rgb}{0.000000,0.000000,0.000000}%
\pgfsetfillcolor{currentfill}%
\pgfsetlinewidth{0.501875pt}%
\definecolor{currentstroke}{rgb}{0.000000,0.000000,0.000000}%
\pgfsetstrokecolor{currentstroke}%
\pgfsetdash{}{0pt}%
\pgfsys@defobject{currentmarker}{\pgfqpoint{0.000000in}{0.000000in}}{\pgfqpoint{0.000000in}{0.020833in}}{%
\pgfpathmoveto{\pgfqpoint{0.000000in}{0.000000in}}%
\pgfpathlineto{\pgfqpoint{0.000000in}{0.020833in}}%
\pgfusepath{stroke,fill}%
}%
\begin{pgfscope}%
\pgfsys@transformshift{5.871851in}{0.413680in}%
\pgfsys@useobject{currentmarker}{}%
\end{pgfscope}%
\end{pgfscope}%
\begin{pgfscope}%
\pgfsetbuttcap%
\pgfsetroundjoin%
\definecolor{currentfill}{rgb}{0.000000,0.000000,0.000000}%
\pgfsetfillcolor{currentfill}%
\pgfsetlinewidth{0.501875pt}%
\definecolor{currentstroke}{rgb}{0.000000,0.000000,0.000000}%
\pgfsetstrokecolor{currentstroke}%
\pgfsetdash{}{0pt}%
\pgfsys@defobject{currentmarker}{\pgfqpoint{0.000000in}{-0.020833in}}{\pgfqpoint{0.000000in}{0.000000in}}{%
\pgfpathmoveto{\pgfqpoint{0.000000in}{0.000000in}}%
\pgfpathlineto{\pgfqpoint{0.000000in}{-0.020833in}}%
\pgfusepath{stroke,fill}%
}%
\begin{pgfscope}%
\pgfsys@transformshift{5.871851in}{3.988523in}%
\pgfsys@useobject{currentmarker}{}%
\end{pgfscope}%
\end{pgfscope}%
\begin{pgfscope}%
\definecolor{textcolor}{rgb}{0.000000,0.000000,0.000000}%
\pgfsetstrokecolor{textcolor}%
\pgfsetfillcolor{textcolor}%
\pgftext[x=3.286122in,y=0.179757in,,top]{\color{textcolor}{\rmfamily\fontsize{10.000000}{12.000000}\selectfont\catcode`\^=\active\def^{\ifmmode\sp\else\^{}\fi}\catcode`\%=\active\def%{\%}$t$}}%
\end{pgfscope}%
\begin{pgfscope}%
\pgfsetbuttcap%
\pgfsetroundjoin%
\definecolor{currentfill}{rgb}{0.000000,0.000000,0.000000}%
\pgfsetfillcolor{currentfill}%
\pgfsetlinewidth{0.501875pt}%
\definecolor{currentstroke}{rgb}{0.000000,0.000000,0.000000}%
\pgfsetstrokecolor{currentstroke}%
\pgfsetdash{}{0pt}%
\pgfsys@defobject{currentmarker}{\pgfqpoint{0.000000in}{0.000000in}}{\pgfqpoint{0.041667in}{0.000000in}}{%
\pgfpathmoveto{\pgfqpoint{0.000000in}{0.000000in}}%
\pgfpathlineto{\pgfqpoint{0.041667in}{0.000000in}}%
\pgfusepath{stroke,fill}%
}%
\begin{pgfscope}%
\pgfsys@transformshift{0.413090in}{0.948442in}%
\pgfsys@useobject{currentmarker}{}%
\end{pgfscope}%
\end{pgfscope}%
\begin{pgfscope}%
\pgfsetbuttcap%
\pgfsetroundjoin%
\definecolor{currentfill}{rgb}{0.000000,0.000000,0.000000}%
\pgfsetfillcolor{currentfill}%
\pgfsetlinewidth{0.501875pt}%
\definecolor{currentstroke}{rgb}{0.000000,0.000000,0.000000}%
\pgfsetstrokecolor{currentstroke}%
\pgfsetdash{}{0pt}%
\pgfsys@defobject{currentmarker}{\pgfqpoint{-0.041667in}{0.000000in}}{\pgfqpoint{-0.000000in}{0.000000in}}{%
\pgfpathmoveto{\pgfqpoint{-0.000000in}{0.000000in}}%
\pgfpathlineto{\pgfqpoint{-0.041667in}{0.000000in}}%
\pgfusepath{stroke,fill}%
}%
\begin{pgfscope}%
\pgfsys@transformshift{6.159154in}{0.948442in}%
\pgfsys@useobject{currentmarker}{}%
\end{pgfscope}%
\end{pgfscope}%
\begin{pgfscope}%
\definecolor{textcolor}{rgb}{0.000000,0.000000,0.000000}%
\pgfsetstrokecolor{textcolor}%
\pgfsetfillcolor{textcolor}%
\pgftext[x=0.299896in, y=0.899744in, left, base]{\color{textcolor}{\rmfamily\fontsize{10.000000}{12.000000}\selectfont\catcode`\^=\active\def^{\ifmmode\sp\else\^{}\fi}\catcode`\%=\active\def%{\%}$\mathdefault{2}$}}%
\end{pgfscope}%
\begin{pgfscope}%
\pgfsetbuttcap%
\pgfsetroundjoin%
\definecolor{currentfill}{rgb}{0.000000,0.000000,0.000000}%
\pgfsetfillcolor{currentfill}%
\pgfsetlinewidth{0.501875pt}%
\definecolor{currentstroke}{rgb}{0.000000,0.000000,0.000000}%
\pgfsetstrokecolor{currentstroke}%
\pgfsetdash{}{0pt}%
\pgfsys@defobject{currentmarker}{\pgfqpoint{0.000000in}{0.000000in}}{\pgfqpoint{0.041667in}{0.000000in}}{%
\pgfpathmoveto{\pgfqpoint{0.000000in}{0.000000in}}%
\pgfpathlineto{\pgfqpoint{0.041667in}{0.000000in}}%
\pgfusepath{stroke,fill}%
}%
\begin{pgfscope}%
\pgfsys@transformshift{0.413090in}{1.623216in}%
\pgfsys@useobject{currentmarker}{}%
\end{pgfscope}%
\end{pgfscope}%
\begin{pgfscope}%
\pgfsetbuttcap%
\pgfsetroundjoin%
\definecolor{currentfill}{rgb}{0.000000,0.000000,0.000000}%
\pgfsetfillcolor{currentfill}%
\pgfsetlinewidth{0.501875pt}%
\definecolor{currentstroke}{rgb}{0.000000,0.000000,0.000000}%
\pgfsetstrokecolor{currentstroke}%
\pgfsetdash{}{0pt}%
\pgfsys@defobject{currentmarker}{\pgfqpoint{-0.041667in}{0.000000in}}{\pgfqpoint{-0.000000in}{0.000000in}}{%
\pgfpathmoveto{\pgfqpoint{-0.000000in}{0.000000in}}%
\pgfpathlineto{\pgfqpoint{-0.041667in}{0.000000in}}%
\pgfusepath{stroke,fill}%
}%
\begin{pgfscope}%
\pgfsys@transformshift{6.159154in}{1.623216in}%
\pgfsys@useobject{currentmarker}{}%
\end{pgfscope}%
\end{pgfscope}%
\begin{pgfscope}%
\definecolor{textcolor}{rgb}{0.000000,0.000000,0.000000}%
\pgfsetstrokecolor{textcolor}%
\pgfsetfillcolor{textcolor}%
\pgftext[x=0.299896in, y=1.574518in, left, base]{\color{textcolor}{\rmfamily\fontsize{10.000000}{12.000000}\selectfont\catcode`\^=\active\def^{\ifmmode\sp\else\^{}\fi}\catcode`\%=\active\def%{\%}$\mathdefault{4}$}}%
\end{pgfscope}%
\begin{pgfscope}%
\pgfsetbuttcap%
\pgfsetroundjoin%
\definecolor{currentfill}{rgb}{0.000000,0.000000,0.000000}%
\pgfsetfillcolor{currentfill}%
\pgfsetlinewidth{0.501875pt}%
\definecolor{currentstroke}{rgb}{0.000000,0.000000,0.000000}%
\pgfsetstrokecolor{currentstroke}%
\pgfsetdash{}{0pt}%
\pgfsys@defobject{currentmarker}{\pgfqpoint{0.000000in}{0.000000in}}{\pgfqpoint{0.041667in}{0.000000in}}{%
\pgfpathmoveto{\pgfqpoint{0.000000in}{0.000000in}}%
\pgfpathlineto{\pgfqpoint{0.041667in}{0.000000in}}%
\pgfusepath{stroke,fill}%
}%
\begin{pgfscope}%
\pgfsys@transformshift{0.413090in}{2.297990in}%
\pgfsys@useobject{currentmarker}{}%
\end{pgfscope}%
\end{pgfscope}%
\begin{pgfscope}%
\pgfsetbuttcap%
\pgfsetroundjoin%
\definecolor{currentfill}{rgb}{0.000000,0.000000,0.000000}%
\pgfsetfillcolor{currentfill}%
\pgfsetlinewidth{0.501875pt}%
\definecolor{currentstroke}{rgb}{0.000000,0.000000,0.000000}%
\pgfsetstrokecolor{currentstroke}%
\pgfsetdash{}{0pt}%
\pgfsys@defobject{currentmarker}{\pgfqpoint{-0.041667in}{0.000000in}}{\pgfqpoint{-0.000000in}{0.000000in}}{%
\pgfpathmoveto{\pgfqpoint{-0.000000in}{0.000000in}}%
\pgfpathlineto{\pgfqpoint{-0.041667in}{0.000000in}}%
\pgfusepath{stroke,fill}%
}%
\begin{pgfscope}%
\pgfsys@transformshift{6.159154in}{2.297990in}%
\pgfsys@useobject{currentmarker}{}%
\end{pgfscope}%
\end{pgfscope}%
\begin{pgfscope}%
\definecolor{textcolor}{rgb}{0.000000,0.000000,0.000000}%
\pgfsetstrokecolor{textcolor}%
\pgfsetfillcolor{textcolor}%
\pgftext[x=0.299896in, y=2.249292in, left, base]{\color{textcolor}{\rmfamily\fontsize{10.000000}{12.000000}\selectfont\catcode`\^=\active\def^{\ifmmode\sp\else\^{}\fi}\catcode`\%=\active\def%{\%}$\mathdefault{6}$}}%
\end{pgfscope}%
\begin{pgfscope}%
\pgfsetbuttcap%
\pgfsetroundjoin%
\definecolor{currentfill}{rgb}{0.000000,0.000000,0.000000}%
\pgfsetfillcolor{currentfill}%
\pgfsetlinewidth{0.501875pt}%
\definecolor{currentstroke}{rgb}{0.000000,0.000000,0.000000}%
\pgfsetstrokecolor{currentstroke}%
\pgfsetdash{}{0pt}%
\pgfsys@defobject{currentmarker}{\pgfqpoint{0.000000in}{0.000000in}}{\pgfqpoint{0.041667in}{0.000000in}}{%
\pgfpathmoveto{\pgfqpoint{0.000000in}{0.000000in}}%
\pgfpathlineto{\pgfqpoint{0.041667in}{0.000000in}}%
\pgfusepath{stroke,fill}%
}%
\begin{pgfscope}%
\pgfsys@transformshift{0.413090in}{2.972764in}%
\pgfsys@useobject{currentmarker}{}%
\end{pgfscope}%
\end{pgfscope}%
\begin{pgfscope}%
\pgfsetbuttcap%
\pgfsetroundjoin%
\definecolor{currentfill}{rgb}{0.000000,0.000000,0.000000}%
\pgfsetfillcolor{currentfill}%
\pgfsetlinewidth{0.501875pt}%
\definecolor{currentstroke}{rgb}{0.000000,0.000000,0.000000}%
\pgfsetstrokecolor{currentstroke}%
\pgfsetdash{}{0pt}%
\pgfsys@defobject{currentmarker}{\pgfqpoint{-0.041667in}{0.000000in}}{\pgfqpoint{-0.000000in}{0.000000in}}{%
\pgfpathmoveto{\pgfqpoint{-0.000000in}{0.000000in}}%
\pgfpathlineto{\pgfqpoint{-0.041667in}{0.000000in}}%
\pgfusepath{stroke,fill}%
}%
\begin{pgfscope}%
\pgfsys@transformshift{6.159154in}{2.972764in}%
\pgfsys@useobject{currentmarker}{}%
\end{pgfscope}%
\end{pgfscope}%
\begin{pgfscope}%
\definecolor{textcolor}{rgb}{0.000000,0.000000,0.000000}%
\pgfsetstrokecolor{textcolor}%
\pgfsetfillcolor{textcolor}%
\pgftext[x=0.299896in, y=2.924066in, left, base]{\color{textcolor}{\rmfamily\fontsize{10.000000}{12.000000}\selectfont\catcode`\^=\active\def^{\ifmmode\sp\else\^{}\fi}\catcode`\%=\active\def%{\%}$\mathdefault{8}$}}%
\end{pgfscope}%
\begin{pgfscope}%
\pgfsetbuttcap%
\pgfsetroundjoin%
\definecolor{currentfill}{rgb}{0.000000,0.000000,0.000000}%
\pgfsetfillcolor{currentfill}%
\pgfsetlinewidth{0.501875pt}%
\definecolor{currentstroke}{rgb}{0.000000,0.000000,0.000000}%
\pgfsetstrokecolor{currentstroke}%
\pgfsetdash{}{0pt}%
\pgfsys@defobject{currentmarker}{\pgfqpoint{0.000000in}{0.000000in}}{\pgfqpoint{0.041667in}{0.000000in}}{%
\pgfpathmoveto{\pgfqpoint{0.000000in}{0.000000in}}%
\pgfpathlineto{\pgfqpoint{0.041667in}{0.000000in}}%
\pgfusepath{stroke,fill}%
}%
\begin{pgfscope}%
\pgfsys@transformshift{0.413090in}{3.647538in}%
\pgfsys@useobject{currentmarker}{}%
\end{pgfscope}%
\end{pgfscope}%
\begin{pgfscope}%
\pgfsetbuttcap%
\pgfsetroundjoin%
\definecolor{currentfill}{rgb}{0.000000,0.000000,0.000000}%
\pgfsetfillcolor{currentfill}%
\pgfsetlinewidth{0.501875pt}%
\definecolor{currentstroke}{rgb}{0.000000,0.000000,0.000000}%
\pgfsetstrokecolor{currentstroke}%
\pgfsetdash{}{0pt}%
\pgfsys@defobject{currentmarker}{\pgfqpoint{-0.041667in}{0.000000in}}{\pgfqpoint{-0.000000in}{0.000000in}}{%
\pgfpathmoveto{\pgfqpoint{-0.000000in}{0.000000in}}%
\pgfpathlineto{\pgfqpoint{-0.041667in}{0.000000in}}%
\pgfusepath{stroke,fill}%
}%
\begin{pgfscope}%
\pgfsys@transformshift{6.159154in}{3.647538in}%
\pgfsys@useobject{currentmarker}{}%
\end{pgfscope}%
\end{pgfscope}%
\begin{pgfscope}%
\definecolor{textcolor}{rgb}{0.000000,0.000000,0.000000}%
\pgfsetstrokecolor{textcolor}%
\pgfsetfillcolor{textcolor}%
\pgftext[x=0.235312in, y=3.598840in, left, base]{\color{textcolor}{\rmfamily\fontsize{10.000000}{12.000000}\selectfont\catcode`\^=\active\def^{\ifmmode\sp\else\^{}\fi}\catcode`\%=\active\def%{\%}$\mathdefault{10}$}}%
\end{pgfscope}%
\begin{pgfscope}%
\pgfsetbuttcap%
\pgfsetroundjoin%
\definecolor{currentfill}{rgb}{0.000000,0.000000,0.000000}%
\pgfsetfillcolor{currentfill}%
\pgfsetlinewidth{0.501875pt}%
\definecolor{currentstroke}{rgb}{0.000000,0.000000,0.000000}%
\pgfsetstrokecolor{currentstroke}%
\pgfsetdash{}{0pt}%
\pgfsys@defobject{currentmarker}{\pgfqpoint{0.000000in}{0.000000in}}{\pgfqpoint{0.020833in}{0.000000in}}{%
\pgfpathmoveto{\pgfqpoint{0.000000in}{0.000000in}}%
\pgfpathlineto{\pgfqpoint{0.020833in}{0.000000in}}%
\pgfusepath{stroke,fill}%
}%
\begin{pgfscope}%
\pgfsys@transformshift{0.413090in}{0.442362in}%
\pgfsys@useobject{currentmarker}{}%
\end{pgfscope}%
\end{pgfscope}%
\begin{pgfscope}%
\pgfsetbuttcap%
\pgfsetroundjoin%
\definecolor{currentfill}{rgb}{0.000000,0.000000,0.000000}%
\pgfsetfillcolor{currentfill}%
\pgfsetlinewidth{0.501875pt}%
\definecolor{currentstroke}{rgb}{0.000000,0.000000,0.000000}%
\pgfsetstrokecolor{currentstroke}%
\pgfsetdash{}{0pt}%
\pgfsys@defobject{currentmarker}{\pgfqpoint{-0.020833in}{0.000000in}}{\pgfqpoint{-0.000000in}{0.000000in}}{%
\pgfpathmoveto{\pgfqpoint{-0.000000in}{0.000000in}}%
\pgfpathlineto{\pgfqpoint{-0.020833in}{0.000000in}}%
\pgfusepath{stroke,fill}%
}%
\begin{pgfscope}%
\pgfsys@transformshift{6.159154in}{0.442362in}%
\pgfsys@useobject{currentmarker}{}%
\end{pgfscope}%
\end{pgfscope}%
\begin{pgfscope}%
\pgfsetbuttcap%
\pgfsetroundjoin%
\definecolor{currentfill}{rgb}{0.000000,0.000000,0.000000}%
\pgfsetfillcolor{currentfill}%
\pgfsetlinewidth{0.501875pt}%
\definecolor{currentstroke}{rgb}{0.000000,0.000000,0.000000}%
\pgfsetstrokecolor{currentstroke}%
\pgfsetdash{}{0pt}%
\pgfsys@defobject{currentmarker}{\pgfqpoint{0.000000in}{0.000000in}}{\pgfqpoint{0.020833in}{0.000000in}}{%
\pgfpathmoveto{\pgfqpoint{0.000000in}{0.000000in}}%
\pgfpathlineto{\pgfqpoint{0.020833in}{0.000000in}}%
\pgfusepath{stroke,fill}%
}%
\begin{pgfscope}%
\pgfsys@transformshift{0.413090in}{0.611055in}%
\pgfsys@useobject{currentmarker}{}%
\end{pgfscope}%
\end{pgfscope}%
\begin{pgfscope}%
\pgfsetbuttcap%
\pgfsetroundjoin%
\definecolor{currentfill}{rgb}{0.000000,0.000000,0.000000}%
\pgfsetfillcolor{currentfill}%
\pgfsetlinewidth{0.501875pt}%
\definecolor{currentstroke}{rgb}{0.000000,0.000000,0.000000}%
\pgfsetstrokecolor{currentstroke}%
\pgfsetdash{}{0pt}%
\pgfsys@defobject{currentmarker}{\pgfqpoint{-0.020833in}{0.000000in}}{\pgfqpoint{-0.000000in}{0.000000in}}{%
\pgfpathmoveto{\pgfqpoint{-0.000000in}{0.000000in}}%
\pgfpathlineto{\pgfqpoint{-0.020833in}{0.000000in}}%
\pgfusepath{stroke,fill}%
}%
\begin{pgfscope}%
\pgfsys@transformshift{6.159154in}{0.611055in}%
\pgfsys@useobject{currentmarker}{}%
\end{pgfscope}%
\end{pgfscope}%
\begin{pgfscope}%
\pgfsetbuttcap%
\pgfsetroundjoin%
\definecolor{currentfill}{rgb}{0.000000,0.000000,0.000000}%
\pgfsetfillcolor{currentfill}%
\pgfsetlinewidth{0.501875pt}%
\definecolor{currentstroke}{rgb}{0.000000,0.000000,0.000000}%
\pgfsetstrokecolor{currentstroke}%
\pgfsetdash{}{0pt}%
\pgfsys@defobject{currentmarker}{\pgfqpoint{0.000000in}{0.000000in}}{\pgfqpoint{0.020833in}{0.000000in}}{%
\pgfpathmoveto{\pgfqpoint{0.000000in}{0.000000in}}%
\pgfpathlineto{\pgfqpoint{0.020833in}{0.000000in}}%
\pgfusepath{stroke,fill}%
}%
\begin{pgfscope}%
\pgfsys@transformshift{0.413090in}{0.779749in}%
\pgfsys@useobject{currentmarker}{}%
\end{pgfscope}%
\end{pgfscope}%
\begin{pgfscope}%
\pgfsetbuttcap%
\pgfsetroundjoin%
\definecolor{currentfill}{rgb}{0.000000,0.000000,0.000000}%
\pgfsetfillcolor{currentfill}%
\pgfsetlinewidth{0.501875pt}%
\definecolor{currentstroke}{rgb}{0.000000,0.000000,0.000000}%
\pgfsetstrokecolor{currentstroke}%
\pgfsetdash{}{0pt}%
\pgfsys@defobject{currentmarker}{\pgfqpoint{-0.020833in}{0.000000in}}{\pgfqpoint{-0.000000in}{0.000000in}}{%
\pgfpathmoveto{\pgfqpoint{-0.000000in}{0.000000in}}%
\pgfpathlineto{\pgfqpoint{-0.020833in}{0.000000in}}%
\pgfusepath{stroke,fill}%
}%
\begin{pgfscope}%
\pgfsys@transformshift{6.159154in}{0.779749in}%
\pgfsys@useobject{currentmarker}{}%
\end{pgfscope}%
\end{pgfscope}%
\begin{pgfscope}%
\pgfsetbuttcap%
\pgfsetroundjoin%
\definecolor{currentfill}{rgb}{0.000000,0.000000,0.000000}%
\pgfsetfillcolor{currentfill}%
\pgfsetlinewidth{0.501875pt}%
\definecolor{currentstroke}{rgb}{0.000000,0.000000,0.000000}%
\pgfsetstrokecolor{currentstroke}%
\pgfsetdash{}{0pt}%
\pgfsys@defobject{currentmarker}{\pgfqpoint{0.000000in}{0.000000in}}{\pgfqpoint{0.020833in}{0.000000in}}{%
\pgfpathmoveto{\pgfqpoint{0.000000in}{0.000000in}}%
\pgfpathlineto{\pgfqpoint{0.020833in}{0.000000in}}%
\pgfusepath{stroke,fill}%
}%
\begin{pgfscope}%
\pgfsys@transformshift{0.413090in}{1.117136in}%
\pgfsys@useobject{currentmarker}{}%
\end{pgfscope}%
\end{pgfscope}%
\begin{pgfscope}%
\pgfsetbuttcap%
\pgfsetroundjoin%
\definecolor{currentfill}{rgb}{0.000000,0.000000,0.000000}%
\pgfsetfillcolor{currentfill}%
\pgfsetlinewidth{0.501875pt}%
\definecolor{currentstroke}{rgb}{0.000000,0.000000,0.000000}%
\pgfsetstrokecolor{currentstroke}%
\pgfsetdash{}{0pt}%
\pgfsys@defobject{currentmarker}{\pgfqpoint{-0.020833in}{0.000000in}}{\pgfqpoint{-0.000000in}{0.000000in}}{%
\pgfpathmoveto{\pgfqpoint{-0.000000in}{0.000000in}}%
\pgfpathlineto{\pgfqpoint{-0.020833in}{0.000000in}}%
\pgfusepath{stroke,fill}%
}%
\begin{pgfscope}%
\pgfsys@transformshift{6.159154in}{1.117136in}%
\pgfsys@useobject{currentmarker}{}%
\end{pgfscope}%
\end{pgfscope}%
\begin{pgfscope}%
\pgfsetbuttcap%
\pgfsetroundjoin%
\definecolor{currentfill}{rgb}{0.000000,0.000000,0.000000}%
\pgfsetfillcolor{currentfill}%
\pgfsetlinewidth{0.501875pt}%
\definecolor{currentstroke}{rgb}{0.000000,0.000000,0.000000}%
\pgfsetstrokecolor{currentstroke}%
\pgfsetdash{}{0pt}%
\pgfsys@defobject{currentmarker}{\pgfqpoint{0.000000in}{0.000000in}}{\pgfqpoint{0.020833in}{0.000000in}}{%
\pgfpathmoveto{\pgfqpoint{0.000000in}{0.000000in}}%
\pgfpathlineto{\pgfqpoint{0.020833in}{0.000000in}}%
\pgfusepath{stroke,fill}%
}%
\begin{pgfscope}%
\pgfsys@transformshift{0.413090in}{1.285829in}%
\pgfsys@useobject{currentmarker}{}%
\end{pgfscope}%
\end{pgfscope}%
\begin{pgfscope}%
\pgfsetbuttcap%
\pgfsetroundjoin%
\definecolor{currentfill}{rgb}{0.000000,0.000000,0.000000}%
\pgfsetfillcolor{currentfill}%
\pgfsetlinewidth{0.501875pt}%
\definecolor{currentstroke}{rgb}{0.000000,0.000000,0.000000}%
\pgfsetstrokecolor{currentstroke}%
\pgfsetdash{}{0pt}%
\pgfsys@defobject{currentmarker}{\pgfqpoint{-0.020833in}{0.000000in}}{\pgfqpoint{-0.000000in}{0.000000in}}{%
\pgfpathmoveto{\pgfqpoint{-0.000000in}{0.000000in}}%
\pgfpathlineto{\pgfqpoint{-0.020833in}{0.000000in}}%
\pgfusepath{stroke,fill}%
}%
\begin{pgfscope}%
\pgfsys@transformshift{6.159154in}{1.285829in}%
\pgfsys@useobject{currentmarker}{}%
\end{pgfscope}%
\end{pgfscope}%
\begin{pgfscope}%
\pgfsetbuttcap%
\pgfsetroundjoin%
\definecolor{currentfill}{rgb}{0.000000,0.000000,0.000000}%
\pgfsetfillcolor{currentfill}%
\pgfsetlinewidth{0.501875pt}%
\definecolor{currentstroke}{rgb}{0.000000,0.000000,0.000000}%
\pgfsetstrokecolor{currentstroke}%
\pgfsetdash{}{0pt}%
\pgfsys@defobject{currentmarker}{\pgfqpoint{0.000000in}{0.000000in}}{\pgfqpoint{0.020833in}{0.000000in}}{%
\pgfpathmoveto{\pgfqpoint{0.000000in}{0.000000in}}%
\pgfpathlineto{\pgfqpoint{0.020833in}{0.000000in}}%
\pgfusepath{stroke,fill}%
}%
\begin{pgfscope}%
\pgfsys@transformshift{0.413090in}{1.454523in}%
\pgfsys@useobject{currentmarker}{}%
\end{pgfscope}%
\end{pgfscope}%
\begin{pgfscope}%
\pgfsetbuttcap%
\pgfsetroundjoin%
\definecolor{currentfill}{rgb}{0.000000,0.000000,0.000000}%
\pgfsetfillcolor{currentfill}%
\pgfsetlinewidth{0.501875pt}%
\definecolor{currentstroke}{rgb}{0.000000,0.000000,0.000000}%
\pgfsetstrokecolor{currentstroke}%
\pgfsetdash{}{0pt}%
\pgfsys@defobject{currentmarker}{\pgfqpoint{-0.020833in}{0.000000in}}{\pgfqpoint{-0.000000in}{0.000000in}}{%
\pgfpathmoveto{\pgfqpoint{-0.000000in}{0.000000in}}%
\pgfpathlineto{\pgfqpoint{-0.020833in}{0.000000in}}%
\pgfusepath{stroke,fill}%
}%
\begin{pgfscope}%
\pgfsys@transformshift{6.159154in}{1.454523in}%
\pgfsys@useobject{currentmarker}{}%
\end{pgfscope}%
\end{pgfscope}%
\begin{pgfscope}%
\pgfsetbuttcap%
\pgfsetroundjoin%
\definecolor{currentfill}{rgb}{0.000000,0.000000,0.000000}%
\pgfsetfillcolor{currentfill}%
\pgfsetlinewidth{0.501875pt}%
\definecolor{currentstroke}{rgb}{0.000000,0.000000,0.000000}%
\pgfsetstrokecolor{currentstroke}%
\pgfsetdash{}{0pt}%
\pgfsys@defobject{currentmarker}{\pgfqpoint{0.000000in}{0.000000in}}{\pgfqpoint{0.020833in}{0.000000in}}{%
\pgfpathmoveto{\pgfqpoint{0.000000in}{0.000000in}}%
\pgfpathlineto{\pgfqpoint{0.020833in}{0.000000in}}%
\pgfusepath{stroke,fill}%
}%
\begin{pgfscope}%
\pgfsys@transformshift{0.413090in}{1.791910in}%
\pgfsys@useobject{currentmarker}{}%
\end{pgfscope}%
\end{pgfscope}%
\begin{pgfscope}%
\pgfsetbuttcap%
\pgfsetroundjoin%
\definecolor{currentfill}{rgb}{0.000000,0.000000,0.000000}%
\pgfsetfillcolor{currentfill}%
\pgfsetlinewidth{0.501875pt}%
\definecolor{currentstroke}{rgb}{0.000000,0.000000,0.000000}%
\pgfsetstrokecolor{currentstroke}%
\pgfsetdash{}{0pt}%
\pgfsys@defobject{currentmarker}{\pgfqpoint{-0.020833in}{0.000000in}}{\pgfqpoint{-0.000000in}{0.000000in}}{%
\pgfpathmoveto{\pgfqpoint{-0.000000in}{0.000000in}}%
\pgfpathlineto{\pgfqpoint{-0.020833in}{0.000000in}}%
\pgfusepath{stroke,fill}%
}%
\begin{pgfscope}%
\pgfsys@transformshift{6.159154in}{1.791910in}%
\pgfsys@useobject{currentmarker}{}%
\end{pgfscope}%
\end{pgfscope}%
\begin{pgfscope}%
\pgfsetbuttcap%
\pgfsetroundjoin%
\definecolor{currentfill}{rgb}{0.000000,0.000000,0.000000}%
\pgfsetfillcolor{currentfill}%
\pgfsetlinewidth{0.501875pt}%
\definecolor{currentstroke}{rgb}{0.000000,0.000000,0.000000}%
\pgfsetstrokecolor{currentstroke}%
\pgfsetdash{}{0pt}%
\pgfsys@defobject{currentmarker}{\pgfqpoint{0.000000in}{0.000000in}}{\pgfqpoint{0.020833in}{0.000000in}}{%
\pgfpathmoveto{\pgfqpoint{0.000000in}{0.000000in}}%
\pgfpathlineto{\pgfqpoint{0.020833in}{0.000000in}}%
\pgfusepath{stroke,fill}%
}%
\begin{pgfscope}%
\pgfsys@transformshift{0.413090in}{1.960603in}%
\pgfsys@useobject{currentmarker}{}%
\end{pgfscope}%
\end{pgfscope}%
\begin{pgfscope}%
\pgfsetbuttcap%
\pgfsetroundjoin%
\definecolor{currentfill}{rgb}{0.000000,0.000000,0.000000}%
\pgfsetfillcolor{currentfill}%
\pgfsetlinewidth{0.501875pt}%
\definecolor{currentstroke}{rgb}{0.000000,0.000000,0.000000}%
\pgfsetstrokecolor{currentstroke}%
\pgfsetdash{}{0pt}%
\pgfsys@defobject{currentmarker}{\pgfqpoint{-0.020833in}{0.000000in}}{\pgfqpoint{-0.000000in}{0.000000in}}{%
\pgfpathmoveto{\pgfqpoint{-0.000000in}{0.000000in}}%
\pgfpathlineto{\pgfqpoint{-0.020833in}{0.000000in}}%
\pgfusepath{stroke,fill}%
}%
\begin{pgfscope}%
\pgfsys@transformshift{6.159154in}{1.960603in}%
\pgfsys@useobject{currentmarker}{}%
\end{pgfscope}%
\end{pgfscope}%
\begin{pgfscope}%
\pgfsetbuttcap%
\pgfsetroundjoin%
\definecolor{currentfill}{rgb}{0.000000,0.000000,0.000000}%
\pgfsetfillcolor{currentfill}%
\pgfsetlinewidth{0.501875pt}%
\definecolor{currentstroke}{rgb}{0.000000,0.000000,0.000000}%
\pgfsetstrokecolor{currentstroke}%
\pgfsetdash{}{0pt}%
\pgfsys@defobject{currentmarker}{\pgfqpoint{0.000000in}{0.000000in}}{\pgfqpoint{0.020833in}{0.000000in}}{%
\pgfpathmoveto{\pgfqpoint{0.000000in}{0.000000in}}%
\pgfpathlineto{\pgfqpoint{0.020833in}{0.000000in}}%
\pgfusepath{stroke,fill}%
}%
\begin{pgfscope}%
\pgfsys@transformshift{0.413090in}{2.129297in}%
\pgfsys@useobject{currentmarker}{}%
\end{pgfscope}%
\end{pgfscope}%
\begin{pgfscope}%
\pgfsetbuttcap%
\pgfsetroundjoin%
\definecolor{currentfill}{rgb}{0.000000,0.000000,0.000000}%
\pgfsetfillcolor{currentfill}%
\pgfsetlinewidth{0.501875pt}%
\definecolor{currentstroke}{rgb}{0.000000,0.000000,0.000000}%
\pgfsetstrokecolor{currentstroke}%
\pgfsetdash{}{0pt}%
\pgfsys@defobject{currentmarker}{\pgfqpoint{-0.020833in}{0.000000in}}{\pgfqpoint{-0.000000in}{0.000000in}}{%
\pgfpathmoveto{\pgfqpoint{-0.000000in}{0.000000in}}%
\pgfpathlineto{\pgfqpoint{-0.020833in}{0.000000in}}%
\pgfusepath{stroke,fill}%
}%
\begin{pgfscope}%
\pgfsys@transformshift{6.159154in}{2.129297in}%
\pgfsys@useobject{currentmarker}{}%
\end{pgfscope}%
\end{pgfscope}%
\begin{pgfscope}%
\pgfsetbuttcap%
\pgfsetroundjoin%
\definecolor{currentfill}{rgb}{0.000000,0.000000,0.000000}%
\pgfsetfillcolor{currentfill}%
\pgfsetlinewidth{0.501875pt}%
\definecolor{currentstroke}{rgb}{0.000000,0.000000,0.000000}%
\pgfsetstrokecolor{currentstroke}%
\pgfsetdash{}{0pt}%
\pgfsys@defobject{currentmarker}{\pgfqpoint{0.000000in}{0.000000in}}{\pgfqpoint{0.020833in}{0.000000in}}{%
\pgfpathmoveto{\pgfqpoint{0.000000in}{0.000000in}}%
\pgfpathlineto{\pgfqpoint{0.020833in}{0.000000in}}%
\pgfusepath{stroke,fill}%
}%
\begin{pgfscope}%
\pgfsys@transformshift{0.413090in}{2.466683in}%
\pgfsys@useobject{currentmarker}{}%
\end{pgfscope}%
\end{pgfscope}%
\begin{pgfscope}%
\pgfsetbuttcap%
\pgfsetroundjoin%
\definecolor{currentfill}{rgb}{0.000000,0.000000,0.000000}%
\pgfsetfillcolor{currentfill}%
\pgfsetlinewidth{0.501875pt}%
\definecolor{currentstroke}{rgb}{0.000000,0.000000,0.000000}%
\pgfsetstrokecolor{currentstroke}%
\pgfsetdash{}{0pt}%
\pgfsys@defobject{currentmarker}{\pgfqpoint{-0.020833in}{0.000000in}}{\pgfqpoint{-0.000000in}{0.000000in}}{%
\pgfpathmoveto{\pgfqpoint{-0.000000in}{0.000000in}}%
\pgfpathlineto{\pgfqpoint{-0.020833in}{0.000000in}}%
\pgfusepath{stroke,fill}%
}%
\begin{pgfscope}%
\pgfsys@transformshift{6.159154in}{2.466683in}%
\pgfsys@useobject{currentmarker}{}%
\end{pgfscope}%
\end{pgfscope}%
\begin{pgfscope}%
\pgfsetbuttcap%
\pgfsetroundjoin%
\definecolor{currentfill}{rgb}{0.000000,0.000000,0.000000}%
\pgfsetfillcolor{currentfill}%
\pgfsetlinewidth{0.501875pt}%
\definecolor{currentstroke}{rgb}{0.000000,0.000000,0.000000}%
\pgfsetstrokecolor{currentstroke}%
\pgfsetdash{}{0pt}%
\pgfsys@defobject{currentmarker}{\pgfqpoint{0.000000in}{0.000000in}}{\pgfqpoint{0.020833in}{0.000000in}}{%
\pgfpathmoveto{\pgfqpoint{0.000000in}{0.000000in}}%
\pgfpathlineto{\pgfqpoint{0.020833in}{0.000000in}}%
\pgfusepath{stroke,fill}%
}%
\begin{pgfscope}%
\pgfsys@transformshift{0.413090in}{2.635377in}%
\pgfsys@useobject{currentmarker}{}%
\end{pgfscope}%
\end{pgfscope}%
\begin{pgfscope}%
\pgfsetbuttcap%
\pgfsetroundjoin%
\definecolor{currentfill}{rgb}{0.000000,0.000000,0.000000}%
\pgfsetfillcolor{currentfill}%
\pgfsetlinewidth{0.501875pt}%
\definecolor{currentstroke}{rgb}{0.000000,0.000000,0.000000}%
\pgfsetstrokecolor{currentstroke}%
\pgfsetdash{}{0pt}%
\pgfsys@defobject{currentmarker}{\pgfqpoint{-0.020833in}{0.000000in}}{\pgfqpoint{-0.000000in}{0.000000in}}{%
\pgfpathmoveto{\pgfqpoint{-0.000000in}{0.000000in}}%
\pgfpathlineto{\pgfqpoint{-0.020833in}{0.000000in}}%
\pgfusepath{stroke,fill}%
}%
\begin{pgfscope}%
\pgfsys@transformshift{6.159154in}{2.635377in}%
\pgfsys@useobject{currentmarker}{}%
\end{pgfscope}%
\end{pgfscope}%
\begin{pgfscope}%
\pgfsetbuttcap%
\pgfsetroundjoin%
\definecolor{currentfill}{rgb}{0.000000,0.000000,0.000000}%
\pgfsetfillcolor{currentfill}%
\pgfsetlinewidth{0.501875pt}%
\definecolor{currentstroke}{rgb}{0.000000,0.000000,0.000000}%
\pgfsetstrokecolor{currentstroke}%
\pgfsetdash{}{0pt}%
\pgfsys@defobject{currentmarker}{\pgfqpoint{0.000000in}{0.000000in}}{\pgfqpoint{0.020833in}{0.000000in}}{%
\pgfpathmoveto{\pgfqpoint{0.000000in}{0.000000in}}%
\pgfpathlineto{\pgfqpoint{0.020833in}{0.000000in}}%
\pgfusepath{stroke,fill}%
}%
\begin{pgfscope}%
\pgfsys@transformshift{0.413090in}{2.804070in}%
\pgfsys@useobject{currentmarker}{}%
\end{pgfscope}%
\end{pgfscope}%
\begin{pgfscope}%
\pgfsetbuttcap%
\pgfsetroundjoin%
\definecolor{currentfill}{rgb}{0.000000,0.000000,0.000000}%
\pgfsetfillcolor{currentfill}%
\pgfsetlinewidth{0.501875pt}%
\definecolor{currentstroke}{rgb}{0.000000,0.000000,0.000000}%
\pgfsetstrokecolor{currentstroke}%
\pgfsetdash{}{0pt}%
\pgfsys@defobject{currentmarker}{\pgfqpoint{-0.020833in}{0.000000in}}{\pgfqpoint{-0.000000in}{0.000000in}}{%
\pgfpathmoveto{\pgfqpoint{-0.000000in}{0.000000in}}%
\pgfpathlineto{\pgfqpoint{-0.020833in}{0.000000in}}%
\pgfusepath{stroke,fill}%
}%
\begin{pgfscope}%
\pgfsys@transformshift{6.159154in}{2.804070in}%
\pgfsys@useobject{currentmarker}{}%
\end{pgfscope}%
\end{pgfscope}%
\begin{pgfscope}%
\pgfsetbuttcap%
\pgfsetroundjoin%
\definecolor{currentfill}{rgb}{0.000000,0.000000,0.000000}%
\pgfsetfillcolor{currentfill}%
\pgfsetlinewidth{0.501875pt}%
\definecolor{currentstroke}{rgb}{0.000000,0.000000,0.000000}%
\pgfsetstrokecolor{currentstroke}%
\pgfsetdash{}{0pt}%
\pgfsys@defobject{currentmarker}{\pgfqpoint{0.000000in}{0.000000in}}{\pgfqpoint{0.020833in}{0.000000in}}{%
\pgfpathmoveto{\pgfqpoint{0.000000in}{0.000000in}}%
\pgfpathlineto{\pgfqpoint{0.020833in}{0.000000in}}%
\pgfusepath{stroke,fill}%
}%
\begin{pgfscope}%
\pgfsys@transformshift{0.413090in}{3.141457in}%
\pgfsys@useobject{currentmarker}{}%
\end{pgfscope}%
\end{pgfscope}%
\begin{pgfscope}%
\pgfsetbuttcap%
\pgfsetroundjoin%
\definecolor{currentfill}{rgb}{0.000000,0.000000,0.000000}%
\pgfsetfillcolor{currentfill}%
\pgfsetlinewidth{0.501875pt}%
\definecolor{currentstroke}{rgb}{0.000000,0.000000,0.000000}%
\pgfsetstrokecolor{currentstroke}%
\pgfsetdash{}{0pt}%
\pgfsys@defobject{currentmarker}{\pgfqpoint{-0.020833in}{0.000000in}}{\pgfqpoint{-0.000000in}{0.000000in}}{%
\pgfpathmoveto{\pgfqpoint{-0.000000in}{0.000000in}}%
\pgfpathlineto{\pgfqpoint{-0.020833in}{0.000000in}}%
\pgfusepath{stroke,fill}%
}%
\begin{pgfscope}%
\pgfsys@transformshift{6.159154in}{3.141457in}%
\pgfsys@useobject{currentmarker}{}%
\end{pgfscope}%
\end{pgfscope}%
\begin{pgfscope}%
\pgfsetbuttcap%
\pgfsetroundjoin%
\definecolor{currentfill}{rgb}{0.000000,0.000000,0.000000}%
\pgfsetfillcolor{currentfill}%
\pgfsetlinewidth{0.501875pt}%
\definecolor{currentstroke}{rgb}{0.000000,0.000000,0.000000}%
\pgfsetstrokecolor{currentstroke}%
\pgfsetdash{}{0pt}%
\pgfsys@defobject{currentmarker}{\pgfqpoint{0.000000in}{0.000000in}}{\pgfqpoint{0.020833in}{0.000000in}}{%
\pgfpathmoveto{\pgfqpoint{0.000000in}{0.000000in}}%
\pgfpathlineto{\pgfqpoint{0.020833in}{0.000000in}}%
\pgfusepath{stroke,fill}%
}%
\begin{pgfscope}%
\pgfsys@transformshift{0.413090in}{3.310151in}%
\pgfsys@useobject{currentmarker}{}%
\end{pgfscope}%
\end{pgfscope}%
\begin{pgfscope}%
\pgfsetbuttcap%
\pgfsetroundjoin%
\definecolor{currentfill}{rgb}{0.000000,0.000000,0.000000}%
\pgfsetfillcolor{currentfill}%
\pgfsetlinewidth{0.501875pt}%
\definecolor{currentstroke}{rgb}{0.000000,0.000000,0.000000}%
\pgfsetstrokecolor{currentstroke}%
\pgfsetdash{}{0pt}%
\pgfsys@defobject{currentmarker}{\pgfqpoint{-0.020833in}{0.000000in}}{\pgfqpoint{-0.000000in}{0.000000in}}{%
\pgfpathmoveto{\pgfqpoint{-0.000000in}{0.000000in}}%
\pgfpathlineto{\pgfqpoint{-0.020833in}{0.000000in}}%
\pgfusepath{stroke,fill}%
}%
\begin{pgfscope}%
\pgfsys@transformshift{6.159154in}{3.310151in}%
\pgfsys@useobject{currentmarker}{}%
\end{pgfscope}%
\end{pgfscope}%
\begin{pgfscope}%
\pgfsetbuttcap%
\pgfsetroundjoin%
\definecolor{currentfill}{rgb}{0.000000,0.000000,0.000000}%
\pgfsetfillcolor{currentfill}%
\pgfsetlinewidth{0.501875pt}%
\definecolor{currentstroke}{rgb}{0.000000,0.000000,0.000000}%
\pgfsetstrokecolor{currentstroke}%
\pgfsetdash{}{0pt}%
\pgfsys@defobject{currentmarker}{\pgfqpoint{0.000000in}{0.000000in}}{\pgfqpoint{0.020833in}{0.000000in}}{%
\pgfpathmoveto{\pgfqpoint{0.000000in}{0.000000in}}%
\pgfpathlineto{\pgfqpoint{0.020833in}{0.000000in}}%
\pgfusepath{stroke,fill}%
}%
\begin{pgfscope}%
\pgfsys@transformshift{0.413090in}{3.478844in}%
\pgfsys@useobject{currentmarker}{}%
\end{pgfscope}%
\end{pgfscope}%
\begin{pgfscope}%
\pgfsetbuttcap%
\pgfsetroundjoin%
\definecolor{currentfill}{rgb}{0.000000,0.000000,0.000000}%
\pgfsetfillcolor{currentfill}%
\pgfsetlinewidth{0.501875pt}%
\definecolor{currentstroke}{rgb}{0.000000,0.000000,0.000000}%
\pgfsetstrokecolor{currentstroke}%
\pgfsetdash{}{0pt}%
\pgfsys@defobject{currentmarker}{\pgfqpoint{-0.020833in}{0.000000in}}{\pgfqpoint{-0.000000in}{0.000000in}}{%
\pgfpathmoveto{\pgfqpoint{-0.000000in}{0.000000in}}%
\pgfpathlineto{\pgfqpoint{-0.020833in}{0.000000in}}%
\pgfusepath{stroke,fill}%
}%
\begin{pgfscope}%
\pgfsys@transformshift{6.159154in}{3.478844in}%
\pgfsys@useobject{currentmarker}{}%
\end{pgfscope}%
\end{pgfscope}%
\begin{pgfscope}%
\pgfsetbuttcap%
\pgfsetroundjoin%
\definecolor{currentfill}{rgb}{0.000000,0.000000,0.000000}%
\pgfsetfillcolor{currentfill}%
\pgfsetlinewidth{0.501875pt}%
\definecolor{currentstroke}{rgb}{0.000000,0.000000,0.000000}%
\pgfsetstrokecolor{currentstroke}%
\pgfsetdash{}{0pt}%
\pgfsys@defobject{currentmarker}{\pgfqpoint{0.000000in}{0.000000in}}{\pgfqpoint{0.020833in}{0.000000in}}{%
\pgfpathmoveto{\pgfqpoint{0.000000in}{0.000000in}}%
\pgfpathlineto{\pgfqpoint{0.020833in}{0.000000in}}%
\pgfusepath{stroke,fill}%
}%
\begin{pgfscope}%
\pgfsys@transformshift{0.413090in}{3.816231in}%
\pgfsys@useobject{currentmarker}{}%
\end{pgfscope}%
\end{pgfscope}%
\begin{pgfscope}%
\pgfsetbuttcap%
\pgfsetroundjoin%
\definecolor{currentfill}{rgb}{0.000000,0.000000,0.000000}%
\pgfsetfillcolor{currentfill}%
\pgfsetlinewidth{0.501875pt}%
\definecolor{currentstroke}{rgb}{0.000000,0.000000,0.000000}%
\pgfsetstrokecolor{currentstroke}%
\pgfsetdash{}{0pt}%
\pgfsys@defobject{currentmarker}{\pgfqpoint{-0.020833in}{0.000000in}}{\pgfqpoint{-0.000000in}{0.000000in}}{%
\pgfpathmoveto{\pgfqpoint{-0.000000in}{0.000000in}}%
\pgfpathlineto{\pgfqpoint{-0.020833in}{0.000000in}}%
\pgfusepath{stroke,fill}%
}%
\begin{pgfscope}%
\pgfsys@transformshift{6.159154in}{3.816231in}%
\pgfsys@useobject{currentmarker}{}%
\end{pgfscope}%
\end{pgfscope}%
\begin{pgfscope}%
\pgfsetbuttcap%
\pgfsetroundjoin%
\definecolor{currentfill}{rgb}{0.000000,0.000000,0.000000}%
\pgfsetfillcolor{currentfill}%
\pgfsetlinewidth{0.501875pt}%
\definecolor{currentstroke}{rgb}{0.000000,0.000000,0.000000}%
\pgfsetstrokecolor{currentstroke}%
\pgfsetdash{}{0pt}%
\pgfsys@defobject{currentmarker}{\pgfqpoint{0.000000in}{0.000000in}}{\pgfqpoint{0.020833in}{0.000000in}}{%
\pgfpathmoveto{\pgfqpoint{0.000000in}{0.000000in}}%
\pgfpathlineto{\pgfqpoint{0.020833in}{0.000000in}}%
\pgfusepath{stroke,fill}%
}%
\begin{pgfscope}%
\pgfsys@transformshift{0.413090in}{3.984925in}%
\pgfsys@useobject{currentmarker}{}%
\end{pgfscope}%
\end{pgfscope}%
\begin{pgfscope}%
\pgfsetbuttcap%
\pgfsetroundjoin%
\definecolor{currentfill}{rgb}{0.000000,0.000000,0.000000}%
\pgfsetfillcolor{currentfill}%
\pgfsetlinewidth{0.501875pt}%
\definecolor{currentstroke}{rgb}{0.000000,0.000000,0.000000}%
\pgfsetstrokecolor{currentstroke}%
\pgfsetdash{}{0pt}%
\pgfsys@defobject{currentmarker}{\pgfqpoint{-0.020833in}{0.000000in}}{\pgfqpoint{-0.000000in}{0.000000in}}{%
\pgfpathmoveto{\pgfqpoint{-0.000000in}{0.000000in}}%
\pgfpathlineto{\pgfqpoint{-0.020833in}{0.000000in}}%
\pgfusepath{stroke,fill}%
}%
\begin{pgfscope}%
\pgfsys@transformshift{6.159154in}{3.984925in}%
\pgfsys@useobject{currentmarker}{}%
\end{pgfscope}%
\end{pgfscope}%
\begin{pgfscope}%
\definecolor{textcolor}{rgb}{0.000000,0.000000,0.000000}%
\pgfsetstrokecolor{textcolor}%
\pgfsetfillcolor{textcolor}%
\pgftext[x=0.179757in,y=2.201101in,,bottom,rotate=90.000000]{\color{textcolor}{\rmfamily\fontsize{10.000000}{12.000000}\selectfont\catcode`\^=\active\def^{\ifmmode\sp\else\^{}\fi}\catcode`\%=\active\def%{\%}$X_t$}}%
\end{pgfscope}%
\begin{pgfscope}%
\pgfpathrectangle{\pgfqpoint{0.413090in}{0.413680in}}{\pgfqpoint{5.746064in}{3.574842in}}%
\pgfusepath{clip}%
\pgfsetrectcap%
\pgfsetroundjoin%
\pgfsetlinewidth{1.003750pt}%
\definecolor{currentstroke}{rgb}{0.000000,0.000000,0.000000}%
\pgfsetstrokecolor{currentstroke}%
\pgfsetdash{}{0pt}%
\pgfpathmoveto{\pgfqpoint{0.413090in}{0.611055in}}%
\pgfpathlineto{\pgfqpoint{0.435536in}{0.614112in}}%
\pgfpathlineto{\pgfqpoint{0.457981in}{0.626982in}}%
\pgfpathlineto{\pgfqpoint{0.480427in}{0.607612in}}%
\pgfpathlineto{\pgfqpoint{0.502872in}{0.624456in}}%
\pgfpathlineto{\pgfqpoint{0.525318in}{0.617412in}}%
\pgfpathlineto{\pgfqpoint{0.547763in}{0.586830in}}%
\pgfpathlineto{\pgfqpoint{0.570209in}{0.576173in}}%
\pgfpathlineto{\pgfqpoint{0.592655in}{0.589550in}}%
\pgfpathlineto{\pgfqpoint{0.615100in}{0.598073in}}%
\pgfpathlineto{\pgfqpoint{0.637546in}{0.577397in}}%
\pgfpathlineto{\pgfqpoint{0.659991in}{0.591226in}}%
\pgfpathlineto{\pgfqpoint{0.682437in}{0.591341in}}%
\pgfpathlineto{\pgfqpoint{0.704882in}{0.601811in}}%
\pgfpathlineto{\pgfqpoint{0.727328in}{0.610665in}}%
\pgfpathlineto{\pgfqpoint{0.749773in}{0.589002in}}%
\pgfpathlineto{\pgfqpoint{0.772219in}{0.597890in}}%
\pgfpathlineto{\pgfqpoint{0.794665in}{0.587657in}}%
\pgfpathlineto{\pgfqpoint{0.817110in}{0.621564in}}%
\pgfpathlineto{\pgfqpoint{0.839556in}{0.639792in}}%
\pgfpathlineto{\pgfqpoint{0.862001in}{0.643007in}}%
\pgfpathlineto{\pgfqpoint{0.884447in}{0.623040in}}%
\pgfpathlineto{\pgfqpoint{0.906892in}{0.626290in}}%
\pgfpathlineto{\pgfqpoint{0.929338in}{0.631925in}}%
\pgfpathlineto{\pgfqpoint{0.951784in}{0.607805in}}%
\pgfpathlineto{\pgfqpoint{0.974229in}{0.659874in}}%
\pgfpathlineto{\pgfqpoint{0.996675in}{0.671836in}}%
\pgfpathlineto{\pgfqpoint{1.019120in}{0.718859in}}%
\pgfpathlineto{\pgfqpoint{1.041566in}{0.691389in}}%
\pgfpathlineto{\pgfqpoint{1.064011in}{0.739072in}}%
\pgfpathlineto{\pgfqpoint{1.086457in}{0.703634in}}%
\pgfpathlineto{\pgfqpoint{1.108902in}{0.688911in}}%
\pgfpathlineto{\pgfqpoint{1.131348in}{0.705755in}}%
\pgfpathlineto{\pgfqpoint{1.153794in}{0.709600in}}%
\pgfpathlineto{\pgfqpoint{1.176239in}{0.652073in}}%
\pgfpathlineto{\pgfqpoint{1.198685in}{0.628884in}}%
\pgfpathlineto{\pgfqpoint{1.221130in}{0.650180in}}%
\pgfpathlineto{\pgfqpoint{1.243576in}{0.705183in}}%
\pgfpathlineto{\pgfqpoint{1.266021in}{0.733260in}}%
\pgfpathlineto{\pgfqpoint{1.288467in}{0.728049in}}%
\pgfpathlineto{\pgfqpoint{1.310913in}{0.754047in}}%
\pgfpathlineto{\pgfqpoint{1.333358in}{0.723533in}}%
\pgfpathlineto{\pgfqpoint{1.355804in}{0.740831in}}%
\pgfpathlineto{\pgfqpoint{1.378249in}{0.705623in}}%
\pgfpathlineto{\pgfqpoint{1.400695in}{0.707907in}}%
\pgfpathlineto{\pgfqpoint{1.423140in}{0.706905in}}%
\pgfpathlineto{\pgfqpoint{1.445586in}{0.731854in}}%
\pgfpathlineto{\pgfqpoint{1.468031in}{0.725902in}}%
\pgfpathlineto{\pgfqpoint{1.490477in}{0.764050in}}%
\pgfpathlineto{\pgfqpoint{1.512923in}{0.760903in}}%
\pgfpathlineto{\pgfqpoint{1.557814in}{0.819282in}}%
\pgfpathlineto{\pgfqpoint{1.580259in}{0.762841in}}%
\pgfpathlineto{\pgfqpoint{1.625150in}{0.770885in}}%
\pgfpathlineto{\pgfqpoint{1.647596in}{0.750048in}}%
\pgfpathlineto{\pgfqpoint{1.670042in}{0.758250in}}%
\pgfpathlineto{\pgfqpoint{1.692487in}{0.718713in}}%
\pgfpathlineto{\pgfqpoint{1.737378in}{0.754981in}}%
\pgfpathlineto{\pgfqpoint{1.759824in}{0.766986in}}%
\pgfpathlineto{\pgfqpoint{1.782269in}{0.852956in}}%
\pgfpathlineto{\pgfqpoint{1.804715in}{0.869335in}}%
\pgfpathlineto{\pgfqpoint{1.827160in}{0.875367in}}%
\pgfpathlineto{\pgfqpoint{1.849606in}{0.851828in}}%
\pgfpathlineto{\pgfqpoint{1.872052in}{0.866127in}}%
\pgfpathlineto{\pgfqpoint{1.894497in}{0.834841in}}%
\pgfpathlineto{\pgfqpoint{1.916943in}{0.884839in}}%
\pgfpathlineto{\pgfqpoint{1.939388in}{0.864044in}}%
\pgfpathlineto{\pgfqpoint{1.961834in}{0.873424in}}%
\pgfpathlineto{\pgfqpoint{1.984279in}{0.956851in}}%
\pgfpathlineto{\pgfqpoint{2.006725in}{1.012731in}}%
\pgfpathlineto{\pgfqpoint{2.029171in}{1.028435in}}%
\pgfpathlineto{\pgfqpoint{2.051616in}{1.101254in}}%
\pgfpathlineto{\pgfqpoint{2.074062in}{1.103284in}}%
\pgfpathlineto{\pgfqpoint{2.096507in}{1.150683in}}%
\pgfpathlineto{\pgfqpoint{2.118953in}{1.156890in}}%
\pgfpathlineto{\pgfqpoint{2.141398in}{1.113201in}}%
\pgfpathlineto{\pgfqpoint{2.163844in}{1.141152in}}%
\pgfpathlineto{\pgfqpoint{2.186289in}{1.198741in}}%
\pgfpathlineto{\pgfqpoint{2.208735in}{1.184265in}}%
\pgfpathlineto{\pgfqpoint{2.231181in}{1.128915in}}%
\pgfpathlineto{\pgfqpoint{2.253626in}{1.106005in}}%
\pgfpathlineto{\pgfqpoint{2.276072in}{0.995188in}}%
\pgfpathlineto{\pgfqpoint{2.298517in}{0.972381in}}%
\pgfpathlineto{\pgfqpoint{2.320963in}{0.930806in}}%
\pgfpathlineto{\pgfqpoint{2.343408in}{1.023728in}}%
\pgfpathlineto{\pgfqpoint{2.365854in}{1.001964in}}%
\pgfpathlineto{\pgfqpoint{2.388300in}{0.938899in}}%
\pgfpathlineto{\pgfqpoint{2.410745in}{0.980237in}}%
\pgfpathlineto{\pgfqpoint{2.433191in}{0.976649in}}%
\pgfpathlineto{\pgfqpoint{2.455636in}{0.982931in}}%
\pgfpathlineto{\pgfqpoint{2.500527in}{1.009365in}}%
\pgfpathlineto{\pgfqpoint{2.522973in}{0.996571in}}%
\pgfpathlineto{\pgfqpoint{2.545418in}{1.045528in}}%
\pgfpathlineto{\pgfqpoint{2.567864in}{1.075874in}}%
\pgfpathlineto{\pgfqpoint{2.590310in}{1.055966in}}%
\pgfpathlineto{\pgfqpoint{2.612755in}{1.052175in}}%
\pgfpathlineto{\pgfqpoint{2.635201in}{1.011868in}}%
\pgfpathlineto{\pgfqpoint{2.657646in}{1.029356in}}%
\pgfpathlineto{\pgfqpoint{2.680092in}{1.069691in}}%
\pgfpathlineto{\pgfqpoint{2.724983in}{1.145363in}}%
\pgfpathlineto{\pgfqpoint{2.747429in}{1.143577in}}%
\pgfpathlineto{\pgfqpoint{2.769874in}{1.239440in}}%
\pgfpathlineto{\pgfqpoint{2.792320in}{1.232618in}}%
\pgfpathlineto{\pgfqpoint{2.814765in}{1.240067in}}%
\pgfpathlineto{\pgfqpoint{2.837211in}{1.225535in}}%
\pgfpathlineto{\pgfqpoint{2.859656in}{1.236004in}}%
\pgfpathlineto{\pgfqpoint{2.882102in}{1.227216in}}%
\pgfpathlineto{\pgfqpoint{2.904547in}{1.362718in}}%
\pgfpathlineto{\pgfqpoint{2.926993in}{1.398665in}}%
\pgfpathlineto{\pgfqpoint{2.949439in}{1.292294in}}%
\pgfpathlineto{\pgfqpoint{2.971884in}{1.270868in}}%
\pgfpathlineto{\pgfqpoint{2.994330in}{1.191759in}}%
\pgfpathlineto{\pgfqpoint{3.016775in}{1.249362in}}%
\pgfpathlineto{\pgfqpoint{3.039221in}{1.343140in}}%
\pgfpathlineto{\pgfqpoint{3.061666in}{1.319523in}}%
\pgfpathlineto{\pgfqpoint{3.084112in}{1.366792in}}%
\pgfpathlineto{\pgfqpoint{3.106558in}{1.358868in}}%
\pgfpathlineto{\pgfqpoint{3.129003in}{1.393443in}}%
\pgfpathlineto{\pgfqpoint{3.151449in}{1.528826in}}%
\pgfpathlineto{\pgfqpoint{3.173894in}{1.508927in}}%
\pgfpathlineto{\pgfqpoint{3.196340in}{1.424021in}}%
\pgfpathlineto{\pgfqpoint{3.218785in}{1.483910in}}%
\pgfpathlineto{\pgfqpoint{3.241231in}{1.530966in}}%
\pgfpathlineto{\pgfqpoint{3.263676in}{1.592751in}}%
\pgfpathlineto{\pgfqpoint{3.286122in}{1.688413in}}%
\pgfpathlineto{\pgfqpoint{3.308568in}{1.637416in}}%
\pgfpathlineto{\pgfqpoint{3.331013in}{1.486818in}}%
\pgfpathlineto{\pgfqpoint{3.353459in}{1.575974in}}%
\pgfpathlineto{\pgfqpoint{3.375904in}{1.550594in}}%
\pgfpathlineto{\pgfqpoint{3.398350in}{1.464738in}}%
\pgfpathlineto{\pgfqpoint{3.420795in}{1.387374in}}%
\pgfpathlineto{\pgfqpoint{3.443241in}{1.475775in}}%
\pgfpathlineto{\pgfqpoint{3.465687in}{1.496583in}}%
\pgfpathlineto{\pgfqpoint{3.488132in}{1.527279in}}%
\pgfpathlineto{\pgfqpoint{3.510578in}{1.603479in}}%
\pgfpathlineto{\pgfqpoint{3.533023in}{1.555884in}}%
\pgfpathlineto{\pgfqpoint{3.555469in}{1.475840in}}%
\pgfpathlineto{\pgfqpoint{3.577914in}{1.420129in}}%
\pgfpathlineto{\pgfqpoint{3.600360in}{1.349940in}}%
\pgfpathlineto{\pgfqpoint{3.622805in}{1.336746in}}%
\pgfpathlineto{\pgfqpoint{3.645251in}{1.285075in}}%
\pgfpathlineto{\pgfqpoint{3.667697in}{1.293478in}}%
\pgfpathlineto{\pgfqpoint{3.690142in}{1.269101in}}%
\pgfpathlineto{\pgfqpoint{3.712588in}{1.269958in}}%
\pgfpathlineto{\pgfqpoint{3.735033in}{1.317350in}}%
\pgfpathlineto{\pgfqpoint{3.757479in}{1.419572in}}%
\pgfpathlineto{\pgfqpoint{3.779924in}{1.544027in}}%
\pgfpathlineto{\pgfqpoint{3.802370in}{1.560703in}}%
\pgfpathlineto{\pgfqpoint{3.824816in}{1.585564in}}%
\pgfpathlineto{\pgfqpoint{3.847261in}{1.552004in}}%
\pgfpathlineto{\pgfqpoint{3.869707in}{1.642958in}}%
\pgfpathlineto{\pgfqpoint{3.892152in}{1.789385in}}%
\pgfpathlineto{\pgfqpoint{3.914598in}{1.742765in}}%
\pgfpathlineto{\pgfqpoint{3.937043in}{1.823682in}}%
\pgfpathlineto{\pgfqpoint{3.959489in}{1.801005in}}%
\pgfpathlineto{\pgfqpoint{3.981934in}{1.793638in}}%
\pgfpathlineto{\pgfqpoint{4.004380in}{1.820281in}}%
\pgfpathlineto{\pgfqpoint{4.026826in}{1.790848in}}%
\pgfpathlineto{\pgfqpoint{4.049271in}{1.887573in}}%
\pgfpathlineto{\pgfqpoint{4.071717in}{1.752164in}}%
\pgfpathlineto{\pgfqpoint{4.094162in}{1.697600in}}%
\pgfpathlineto{\pgfqpoint{4.116608in}{1.695319in}}%
\pgfpathlineto{\pgfqpoint{4.139053in}{1.613275in}}%
\pgfpathlineto{\pgfqpoint{4.161499in}{1.651864in}}%
\pgfpathlineto{\pgfqpoint{4.183945in}{1.534729in}}%
\pgfpathlineto{\pgfqpoint{4.206390in}{1.546522in}}%
\pgfpathlineto{\pgfqpoint{4.228836in}{1.443786in}}%
\pgfpathlineto{\pgfqpoint{4.251281in}{1.440026in}}%
\pgfpathlineto{\pgfqpoint{4.273727in}{1.484733in}}%
\pgfpathlineto{\pgfqpoint{4.296172in}{1.671723in}}%
\pgfpathlineto{\pgfqpoint{4.318618in}{1.777916in}}%
\pgfpathlineto{\pgfqpoint{4.341063in}{1.947692in}}%
\pgfpathlineto{\pgfqpoint{4.363509in}{1.964698in}}%
\pgfpathlineto{\pgfqpoint{4.385955in}{2.215110in}}%
\pgfpathlineto{\pgfqpoint{4.430846in}{2.178338in}}%
\pgfpathlineto{\pgfqpoint{4.453291in}{2.042002in}}%
\pgfpathlineto{\pgfqpoint{4.475737in}{2.077274in}}%
\pgfpathlineto{\pgfqpoint{4.498182in}{2.211359in}}%
\pgfpathlineto{\pgfqpoint{4.520628in}{2.147041in}}%
\pgfpathlineto{\pgfqpoint{4.543074in}{2.114904in}}%
\pgfpathlineto{\pgfqpoint{4.565519in}{2.072837in}}%
\pgfpathlineto{\pgfqpoint{4.587965in}{2.046306in}}%
\pgfpathlineto{\pgfqpoint{4.610410in}{1.992129in}}%
\pgfpathlineto{\pgfqpoint{4.632856in}{1.862360in}}%
\pgfpathlineto{\pgfqpoint{4.655301in}{2.015750in}}%
\pgfpathlineto{\pgfqpoint{4.677747in}{1.822817in}}%
\pgfpathlineto{\pgfqpoint{4.700192in}{1.744197in}}%
\pgfpathlineto{\pgfqpoint{4.722638in}{1.784961in}}%
\pgfpathlineto{\pgfqpoint{4.745084in}{1.893055in}}%
\pgfpathlineto{\pgfqpoint{4.767529in}{1.881472in}}%
\pgfpathlineto{\pgfqpoint{4.789975in}{1.862542in}}%
\pgfpathlineto{\pgfqpoint{4.812420in}{1.876829in}}%
\pgfpathlineto{\pgfqpoint{4.834866in}{1.860128in}}%
\pgfpathlineto{\pgfqpoint{4.857311in}{1.959334in}}%
\pgfpathlineto{\pgfqpoint{4.879757in}{1.918832in}}%
\pgfpathlineto{\pgfqpoint{4.902203in}{1.829656in}}%
\pgfpathlineto{\pgfqpoint{4.924648in}{1.793240in}}%
\pgfpathlineto{\pgfqpoint{4.947094in}{1.811394in}}%
\pgfpathlineto{\pgfqpoint{4.969539in}{1.943946in}}%
\pgfpathlineto{\pgfqpoint{4.991985in}{2.130444in}}%
\pgfpathlineto{\pgfqpoint{5.014430in}{2.095309in}}%
\pgfpathlineto{\pgfqpoint{5.036876in}{1.945777in}}%
\pgfpathlineto{\pgfqpoint{5.059321in}{1.816244in}}%
\pgfpathlineto{\pgfqpoint{5.081767in}{2.045733in}}%
\pgfpathlineto{\pgfqpoint{5.104213in}{2.011825in}}%
\pgfpathlineto{\pgfqpoint{5.126658in}{2.072799in}}%
\pgfpathlineto{\pgfqpoint{5.149104in}{2.313758in}}%
\pgfpathlineto{\pgfqpoint{5.171549in}{2.240346in}}%
\pgfpathlineto{\pgfqpoint{5.193995in}{2.227992in}}%
\pgfpathlineto{\pgfqpoint{5.216440in}{2.169422in}}%
\pgfpathlineto{\pgfqpoint{5.238886in}{2.054323in}}%
\pgfpathlineto{\pgfqpoint{5.261332in}{2.156591in}}%
\pgfpathlineto{\pgfqpoint{5.283777in}{2.298103in}}%
\pgfpathlineto{\pgfqpoint{5.306223in}{2.084962in}}%
\pgfpathlineto{\pgfqpoint{5.328668in}{2.210895in}}%
\pgfpathlineto{\pgfqpoint{5.351114in}{2.193753in}}%
\pgfpathlineto{\pgfqpoint{5.373559in}{2.227329in}}%
\pgfpathlineto{\pgfqpoint{5.396005in}{2.263264in}}%
\pgfpathlineto{\pgfqpoint{5.418451in}{2.118028in}}%
\pgfpathlineto{\pgfqpoint{5.440896in}{2.093169in}}%
\pgfpathlineto{\pgfqpoint{5.463342in}{1.993952in}}%
\pgfpathlineto{\pgfqpoint{5.485787in}{2.224102in}}%
\pgfpathlineto{\pgfqpoint{5.508233in}{2.522761in}}%
\pgfpathlineto{\pgfqpoint{5.530678in}{2.547891in}}%
\pgfpathlineto{\pgfqpoint{5.553124in}{2.398755in}}%
\pgfpathlineto{\pgfqpoint{5.575569in}{2.571408in}}%
\pgfpathlineto{\pgfqpoint{5.598015in}{2.557626in}}%
\pgfpathlineto{\pgfqpoint{5.620461in}{2.835851in}}%
\pgfpathlineto{\pgfqpoint{5.665352in}{3.528628in}}%
\pgfpathlineto{\pgfqpoint{5.687797in}{3.467865in}}%
\pgfpathlineto{\pgfqpoint{5.710243in}{3.292335in}}%
\pgfpathlineto{\pgfqpoint{5.732688in}{3.186816in}}%
\pgfpathlineto{\pgfqpoint{5.755134in}{3.186343in}}%
\pgfpathlineto{\pgfqpoint{5.777580in}{3.039090in}}%
\pgfpathlineto{\pgfqpoint{5.800025in}{3.270747in}}%
\pgfpathlineto{\pgfqpoint{5.822471in}{3.414819in}}%
\pgfpathlineto{\pgfqpoint{5.844916in}{3.426578in}}%
\pgfpathlineto{\pgfqpoint{5.867362in}{3.366727in}}%
\pgfpathlineto{\pgfqpoint{5.889807in}{3.513761in}}%
\pgfpathlineto{\pgfqpoint{5.912253in}{3.826030in}}%
\pgfpathlineto{\pgfqpoint{5.934698in}{3.440906in}}%
\pgfpathlineto{\pgfqpoint{5.957144in}{3.506428in}}%
\pgfpathlineto{\pgfqpoint{5.979590in}{3.464931in}}%
\pgfpathlineto{\pgfqpoint{6.002035in}{3.337792in}}%
\pgfpathlineto{\pgfqpoint{6.024481in}{3.290757in}}%
\pgfpathlineto{\pgfqpoint{6.046926in}{3.334110in}}%
\pgfpathlineto{\pgfqpoint{6.069372in}{3.733660in}}%
\pgfpathlineto{\pgfqpoint{6.091817in}{3.766427in}}%
\pgfpathlineto{\pgfqpoint{6.114263in}{3.693840in}}%
\pgfpathlineto{\pgfqpoint{6.136709in}{3.459938in}}%
\pgfpathlineto{\pgfqpoint{6.159154in}{3.637515in}}%
\pgfpathlineto{\pgfqpoint{6.159154in}{3.637515in}}%
\pgfusepath{stroke}%
\end{pgfscope}%
\begin{pgfscope}%
\pgfpathrectangle{\pgfqpoint{0.413090in}{0.413680in}}{\pgfqpoint{5.746064in}{3.574842in}}%
\pgfusepath{clip}%
\pgfsetbuttcap%
\pgfsetroundjoin%
\pgfsetlinewidth{1.003750pt}%
\definecolor{currentstroke}{rgb}{0.690196,0.188235,0.376471}%
\pgfsetstrokecolor{currentstroke}%
\pgfsetdash{{3.700000pt}{1.600000pt}}{0.000000pt}%
\pgfpathmoveto{\pgfqpoint{0.413090in}{0.611055in}}%
\pgfpathlineto{\pgfqpoint{0.435536in}{0.614100in}}%
\pgfpathlineto{\pgfqpoint{0.457981in}{0.626898in}}%
\pgfpathlineto{\pgfqpoint{0.480427in}{0.607666in}}%
\pgfpathlineto{\pgfqpoint{0.502872in}{0.624415in}}%
\pgfpathlineto{\pgfqpoint{0.525318in}{0.617421in}}%
\pgfpathlineto{\pgfqpoint{0.547763in}{0.587077in}}%
\pgfpathlineto{\pgfqpoint{0.570209in}{0.576482in}}%
\pgfpathlineto{\pgfqpoint{0.592655in}{0.589797in}}%
\pgfpathlineto{\pgfqpoint{0.615100in}{0.598282in}}%
\pgfpathlineto{\pgfqpoint{0.637546in}{0.577739in}}%
\pgfpathlineto{\pgfqpoint{0.659991in}{0.591505in}}%
\pgfpathlineto{\pgfqpoint{0.682437in}{0.591624in}}%
\pgfpathlineto{\pgfqpoint{0.704882in}{0.602047in}}%
\pgfpathlineto{\pgfqpoint{0.727328in}{0.610861in}}%
\pgfpathlineto{\pgfqpoint{0.749773in}{0.589339in}}%
\pgfpathlineto{\pgfqpoint{0.772219in}{0.598189in}}%
\pgfpathlineto{\pgfqpoint{0.794665in}{0.588015in}}%
\pgfpathlineto{\pgfqpoint{0.817110in}{0.621719in}}%
\pgfpathlineto{\pgfqpoint{0.839556in}{0.639849in}}%
\pgfpathlineto{\pgfqpoint{0.862001in}{0.643052in}}%
\pgfpathlineto{\pgfqpoint{0.884447in}{0.623220in}}%
\pgfpathlineto{\pgfqpoint{0.906892in}{0.626459in}}%
\pgfpathlineto{\pgfqpoint{0.929338in}{0.632069in}}%
\pgfpathlineto{\pgfqpoint{0.951784in}{0.608112in}}%
\pgfpathlineto{\pgfqpoint{0.974229in}{0.659775in}}%
\pgfpathlineto{\pgfqpoint{0.996675in}{0.671669in}}%
\pgfpathlineto{\pgfqpoint{1.019120in}{0.718325in}}%
\pgfpathlineto{\pgfqpoint{1.041566in}{0.691080in}}%
\pgfpathlineto{\pgfqpoint{1.064011in}{0.738381in}}%
\pgfpathlineto{\pgfqpoint{1.086457in}{0.703256in}}%
\pgfpathlineto{\pgfqpoint{1.108902in}{0.688644in}}%
\pgfpathlineto{\pgfqpoint{1.131348in}{0.705383in}}%
\pgfpathlineto{\pgfqpoint{1.153794in}{0.709210in}}%
\pgfpathlineto{\pgfqpoint{1.176239in}{0.652302in}}%
\pgfpathlineto{\pgfqpoint{1.198685in}{0.629262in}}%
\pgfpathlineto{\pgfqpoint{1.221130in}{0.650453in}}%
\pgfpathlineto{\pgfqpoint{1.243576in}{0.705036in}}%
\pgfpathlineto{\pgfqpoint{1.266021in}{0.732934in}}%
\pgfpathlineto{\pgfqpoint{1.288467in}{0.727765in}}%
\pgfpathlineto{\pgfqpoint{1.310913in}{0.753594in}}%
\pgfpathlineto{\pgfqpoint{1.333358in}{0.723324in}}%
\pgfpathlineto{\pgfqpoint{1.355804in}{0.740517in}}%
\pgfpathlineto{\pgfqpoint{1.378249in}{0.705592in}}%
\pgfpathlineto{\pgfqpoint{1.400695in}{0.707870in}}%
\pgfpathlineto{\pgfqpoint{1.423140in}{0.706882in}}%
\pgfpathlineto{\pgfqpoint{1.445586in}{0.731682in}}%
\pgfpathlineto{\pgfqpoint{1.468031in}{0.725775in}}%
\pgfpathlineto{\pgfqpoint{1.490477in}{0.763664in}}%
\pgfpathlineto{\pgfqpoint{1.512923in}{0.760547in}}%
\pgfpathlineto{\pgfqpoint{1.557814in}{0.818528in}}%
\pgfpathlineto{\pgfqpoint{1.580259in}{0.762638in}}%
\pgfpathlineto{\pgfqpoint{1.625150in}{0.770649in}}%
\pgfpathlineto{\pgfqpoint{1.647596in}{0.749960in}}%
\pgfpathlineto{\pgfqpoint{1.670042in}{0.758122in}}%
\pgfpathlineto{\pgfqpoint{1.692487in}{0.718894in}}%
\pgfpathlineto{\pgfqpoint{1.737378in}{0.754970in}}%
\pgfpathlineto{\pgfqpoint{1.759824in}{0.766913in}}%
\pgfpathlineto{\pgfqpoint{1.782269in}{0.852060in}}%
\pgfpathlineto{\pgfqpoint{1.804715in}{0.868326in}}%
\pgfpathlineto{\pgfqpoint{1.827160in}{0.874323in}}%
\pgfpathlineto{\pgfqpoint{1.849606in}{0.850982in}}%
\pgfpathlineto{\pgfqpoint{1.872052in}{0.865185in}}%
\pgfpathlineto{\pgfqpoint{1.894497in}{0.834163in}}%
\pgfpathlineto{\pgfqpoint{1.916943in}{0.883771in}}%
\pgfpathlineto{\pgfqpoint{1.939388in}{0.863150in}}%
\pgfpathlineto{\pgfqpoint{1.961834in}{0.872472in}}%
\pgfpathlineto{\pgfqpoint{1.984279in}{0.955091in}}%
\pgfpathlineto{\pgfqpoint{2.006725in}{1.010469in}}%
\pgfpathlineto{\pgfqpoint{2.029171in}{1.026046in}}%
\pgfpathlineto{\pgfqpoint{2.051616in}{1.098141in}}%
\pgfpathlineto{\pgfqpoint{2.074062in}{1.100166in}}%
\pgfpathlineto{\pgfqpoint{2.096507in}{1.147108in}}%
\pgfpathlineto{\pgfqpoint{2.118953in}{1.153268in}}%
\pgfpathlineto{\pgfqpoint{2.141398in}{1.110054in}}%
\pgfpathlineto{\pgfqpoint{2.163844in}{1.137749in}}%
\pgfpathlineto{\pgfqpoint{2.186289in}{1.194763in}}%
\pgfpathlineto{\pgfqpoint{2.208735in}{1.180451in}}%
\pgfpathlineto{\pgfqpoint{2.231181in}{1.125717in}}%
\pgfpathlineto{\pgfqpoint{2.253626in}{1.103046in}}%
\pgfpathlineto{\pgfqpoint{2.276072in}{0.993721in}}%
\pgfpathlineto{\pgfqpoint{2.298517in}{0.971112in}}%
\pgfpathlineto{\pgfqpoint{2.320963in}{0.929901in}}%
\pgfpathlineto{\pgfqpoint{2.343408in}{1.021937in}}%
\pgfpathlineto{\pgfqpoint{2.365854in}{1.000370in}}%
\pgfpathlineto{\pgfqpoint{2.388300in}{0.937929in}}%
\pgfpathlineto{\pgfqpoint{2.410745in}{0.978959in}}%
\pgfpathlineto{\pgfqpoint{2.433191in}{0.975411in}}%
\pgfpathlineto{\pgfqpoint{2.455636in}{0.981657in}}%
\pgfpathlineto{\pgfqpoint{2.500527in}{1.007912in}}%
\pgfpathlineto{\pgfqpoint{2.522973in}{0.995232in}}%
\pgfpathlineto{\pgfqpoint{2.545418in}{1.043799in}}%
\pgfpathlineto{\pgfqpoint{2.567864in}{1.073909in}}%
\pgfpathlineto{\pgfqpoint{2.590310in}{1.054184in}}%
\pgfpathlineto{\pgfqpoint{2.612755in}{1.050437in}}%
\pgfpathlineto{\pgfqpoint{2.635201in}{1.010495in}}%
\pgfpathlineto{\pgfqpoint{2.657646in}{1.027860in}}%
\pgfpathlineto{\pgfqpoint{2.680092in}{1.067880in}}%
\pgfpathlineto{\pgfqpoint{2.724983in}{1.142941in}}%
\pgfpathlineto{\pgfqpoint{2.747429in}{1.141185in}}%
\pgfpathlineto{\pgfqpoint{2.769874in}{1.236098in}}%
\pgfpathlineto{\pgfqpoint{2.792320in}{1.229356in}}%
\pgfpathlineto{\pgfqpoint{2.814765in}{1.236753in}}%
\pgfpathlineto{\pgfqpoint{2.837211in}{1.222373in}}%
\pgfpathlineto{\pgfqpoint{2.859656in}{1.232763in}}%
\pgfpathlineto{\pgfqpoint{2.882102in}{1.224073in}}%
\pgfpathlineto{\pgfqpoint{2.904547in}{1.358020in}}%
\pgfpathlineto{\pgfqpoint{2.926993in}{1.393618in}}%
\pgfpathlineto{\pgfqpoint{2.949439in}{1.288573in}}%
\pgfpathlineto{\pgfqpoint{2.971884in}{1.267371in}}%
\pgfpathlineto{\pgfqpoint{2.994330in}{1.189129in}}%
\pgfpathlineto{\pgfqpoint{3.016775in}{1.246222in}}%
\pgfpathlineto{\pgfqpoint{3.039221in}{1.339064in}}%
\pgfpathlineto{\pgfqpoint{3.061666in}{1.315696in}}%
\pgfpathlineto{\pgfqpoint{3.084112in}{1.362525in}}%
\pgfpathlineto{\pgfqpoint{3.106558in}{1.354698in}}%
\pgfpathlineto{\pgfqpoint{3.129003in}{1.388953in}}%
\pgfpathlineto{\pgfqpoint{3.151449in}{1.522783in}}%
\pgfpathlineto{\pgfqpoint{3.173894in}{1.503119in}}%
\pgfpathlineto{\pgfqpoint{3.196340in}{1.419219in}}%
\pgfpathlineto{\pgfqpoint{3.218785in}{1.478512in}}%
\pgfpathlineto{\pgfqpoint{3.241231in}{1.525098in}}%
\pgfpathlineto{\pgfqpoint{3.263676in}{1.586241in}}%
\pgfpathlineto{\pgfqpoint{3.286122in}{1.680839in}}%
\pgfpathlineto{\pgfqpoint{3.308568in}{1.630453in}}%
\pgfpathlineto{\pgfqpoint{3.331013in}{1.481930in}}%
\pgfpathlineto{\pgfqpoint{3.353459in}{1.570172in}}%
\pgfpathlineto{\pgfqpoint{3.375904in}{1.545078in}}%
\pgfpathlineto{\pgfqpoint{3.398350in}{1.460204in}}%
\pgfpathlineto{\pgfqpoint{3.420795in}{1.383681in}}%
\pgfpathlineto{\pgfqpoint{3.443241in}{1.471227in}}%
\pgfpathlineto{\pgfqpoint{3.465687in}{1.491855in}}%
\pgfpathlineto{\pgfqpoint{3.488132in}{1.522273in}}%
\pgfpathlineto{\pgfqpoint{3.510578in}{1.597715in}}%
\pgfpathlineto{\pgfqpoint{3.533023in}{1.550643in}}%
\pgfpathlineto{\pgfqpoint{3.555469in}{1.471486in}}%
\pgfpathlineto{\pgfqpoint{3.577914in}{1.416351in}}%
\pgfpathlineto{\pgfqpoint{3.600360in}{1.346883in}}%
\pgfpathlineto{\pgfqpoint{3.622805in}{1.333824in}}%
\pgfpathlineto{\pgfqpoint{3.645251in}{1.282644in}}%
\pgfpathlineto{\pgfqpoint{3.667697in}{1.290995in}}%
\pgfpathlineto{\pgfqpoint{3.690142in}{1.266841in}}%
\pgfpathlineto{\pgfqpoint{3.712588in}{1.267709in}}%
\pgfpathlineto{\pgfqpoint{3.735033in}{1.314723in}}%
\pgfpathlineto{\pgfqpoint{3.757479in}{1.415991in}}%
\pgfpathlineto{\pgfqpoint{3.779924in}{1.539172in}}%
\pgfpathlineto{\pgfqpoint{3.802370in}{1.555709in}}%
\pgfpathlineto{\pgfqpoint{3.824816in}{1.580351in}}%
\pgfpathlineto{\pgfqpoint{3.847261in}{1.547148in}}%
\pgfpathlineto{\pgfqpoint{3.869707in}{1.637196in}}%
\pgfpathlineto{\pgfqpoint{3.892152in}{1.781974in}}%
\pgfpathlineto{\pgfqpoint{3.914598in}{1.735891in}}%
\pgfpathlineto{\pgfqpoint{3.937043in}{1.815957in}}%
\pgfpathlineto{\pgfqpoint{3.959489in}{1.793553in}}%
\pgfpathlineto{\pgfqpoint{3.981934in}{1.786291in}}%
\pgfpathlineto{\pgfqpoint{4.004380in}{1.812677in}}%
\pgfpathlineto{\pgfqpoint{4.026826in}{1.783590in}}%
\pgfpathlineto{\pgfqpoint{4.049271in}{1.879271in}}%
\pgfpathlineto{\pgfqpoint{4.071717in}{1.745586in}}%
\pgfpathlineto{\pgfqpoint{4.094162in}{1.691628in}}%
\pgfpathlineto{\pgfqpoint{4.116608in}{1.689394in}}%
\pgfpathlineto{\pgfqpoint{4.139053in}{1.608258in}}%
\pgfpathlineto{\pgfqpoint{4.161499in}{1.646499in}}%
\pgfpathlineto{\pgfqpoint{4.183945in}{1.530714in}}%
\pgfpathlineto{\pgfqpoint{4.206390in}{1.542422in}}%
\pgfpathlineto{\pgfqpoint{4.228836in}{1.440788in}}%
\pgfpathlineto{\pgfqpoint{4.251281in}{1.437080in}}%
\pgfpathlineto{\pgfqpoint{4.273727in}{1.481427in}}%
\pgfpathlineto{\pgfqpoint{4.296172in}{1.666293in}}%
\pgfpathlineto{\pgfqpoint{4.318618in}{1.771408in}}%
\pgfpathlineto{\pgfqpoint{4.341063in}{1.939223in}}%
\pgfpathlineto{\pgfqpoint{4.363509in}{1.956073in}}%
\pgfpathlineto{\pgfqpoint{4.385955in}{2.203110in}}%
\pgfpathlineto{\pgfqpoint{4.430846in}{2.166845in}}%
\pgfpathlineto{\pgfqpoint{4.453291in}{2.032316in}}%
\pgfpathlineto{\pgfqpoint{4.475737in}{2.067219in}}%
\pgfpathlineto{\pgfqpoint{4.498182in}{2.199722in}}%
\pgfpathlineto{\pgfqpoint{4.520628in}{2.136216in}}%
\pgfpathlineto{\pgfqpoint{4.543074in}{2.104487in}}%
\pgfpathlineto{\pgfqpoint{4.565519in}{2.062941in}}%
\pgfpathlineto{\pgfqpoint{4.587965in}{2.036743in}}%
\pgfpathlineto{\pgfqpoint{4.610410in}{1.983218in}}%
\pgfpathlineto{\pgfqpoint{4.632856in}{1.855073in}}%
\pgfpathlineto{\pgfqpoint{4.655301in}{2.006722in}}%
\pgfpathlineto{\pgfqpoint{4.677747in}{1.816464in}}%
\pgfpathlineto{\pgfqpoint{4.700192in}{1.738700in}}%
\pgfpathlineto{\pgfqpoint{4.722638in}{1.779097in}}%
\pgfpathlineto{\pgfqpoint{4.745084in}{1.886103in}}%
\pgfpathlineto{\pgfqpoint{4.767529in}{1.874666in}}%
\pgfpathlineto{\pgfqpoint{4.789975in}{1.855955in}}%
\pgfpathlineto{\pgfqpoint{4.812420in}{1.870127in}}%
\pgfpathlineto{\pgfqpoint{4.834866in}{1.853621in}}%
\pgfpathlineto{\pgfqpoint{4.857311in}{1.951826in}}%
\pgfpathlineto{\pgfqpoint{4.879757in}{1.911777in}}%
\pgfpathlineto{\pgfqpoint{4.902203in}{1.823593in}}%
\pgfpathlineto{\pgfqpoint{4.924648in}{1.787564in}}%
\pgfpathlineto{\pgfqpoint{4.947094in}{1.805570in}}%
\pgfpathlineto{\pgfqpoint{4.969539in}{1.936757in}}%
\pgfpathlineto{\pgfqpoint{4.991985in}{2.121112in}}%
\pgfpathlineto{\pgfqpoint{5.014430in}{2.086394in}}%
\pgfpathlineto{\pgfqpoint{5.036876in}{1.938720in}}%
\pgfpathlineto{\pgfqpoint{5.059321in}{1.810691in}}%
\pgfpathlineto{\pgfqpoint{5.081767in}{2.037424in}}%
\pgfpathlineto{\pgfqpoint{5.104213in}{2.003908in}}%
\pgfpathlineto{\pgfqpoint{5.126658in}{2.064272in}}%
\pgfpathlineto{\pgfqpoint{5.149104in}{2.302189in}}%
\pgfpathlineto{\pgfqpoint{5.171549in}{2.229680in}}%
\pgfpathlineto{\pgfqpoint{5.193995in}{2.217499in}}%
\pgfpathlineto{\pgfqpoint{5.216440in}{2.159632in}}%
\pgfpathlineto{\pgfqpoint{5.238886in}{2.045928in}}%
\pgfpathlineto{\pgfqpoint{5.261332in}{2.147112in}}%
\pgfpathlineto{\pgfqpoint{5.283777in}{2.287028in}}%
\pgfpathlineto{\pgfqpoint{5.306223in}{2.076850in}}%
\pgfpathlineto{\pgfqpoint{5.328668in}{2.201447in}}%
\pgfpathlineto{\pgfqpoint{5.351114in}{2.184521in}}%
\pgfpathlineto{\pgfqpoint{5.373559in}{2.217774in}}%
\pgfpathlineto{\pgfqpoint{5.396005in}{2.253356in}}%
\pgfpathlineto{\pgfqpoint{5.418451in}{2.109898in}}%
\pgfpathlineto{\pgfqpoint{5.440896in}{2.085326in}}%
\pgfpathlineto{\pgfqpoint{5.463342in}{1.987216in}}%
\pgfpathlineto{\pgfqpoint{5.485787in}{2.214650in}}%
\pgfpathlineto{\pgfqpoint{5.508233in}{2.509305in}}%
\pgfpathlineto{\pgfqpoint{5.530678in}{2.534177in}}%
\pgfpathlineto{\pgfqpoint{5.553124in}{2.386994in}}%
\pgfpathlineto{\pgfqpoint{5.575569in}{2.557588in}}%
\pgfpathlineto{\pgfqpoint{5.598015in}{2.544009in}}%
\pgfpathlineto{\pgfqpoint{5.620461in}{2.818494in}}%
\pgfpathlineto{\pgfqpoint{5.665352in}{3.500943in}}%
\pgfpathlineto{\pgfqpoint{5.687797in}{3.441112in}}%
\pgfpathlineto{\pgfqpoint{5.710243in}{3.268250in}}%
\pgfpathlineto{\pgfqpoint{5.732688in}{3.164270in}}%
\pgfpathlineto{\pgfqpoint{5.755134in}{3.163854in}}%
\pgfpathlineto{\pgfqpoint{5.777580in}{3.018732in}}%
\pgfpathlineto{\pgfqpoint{5.800025in}{3.247194in}}%
\pgfpathlineto{\pgfqpoint{5.822471in}{3.389314in}}%
\pgfpathlineto{\pgfqpoint{5.844916in}{3.400963in}}%
\pgfpathlineto{\pgfqpoint{5.867362in}{3.342007in}}%
\pgfpathlineto{\pgfqpoint{5.889807in}{3.487032in}}%
\pgfpathlineto{\pgfqpoint{5.912253in}{3.794620in}}%
\pgfpathlineto{\pgfqpoint{5.934698in}{3.416191in}}%
\pgfpathlineto{\pgfqpoint{5.957144in}{3.480875in}}%
\pgfpathlineto{\pgfqpoint{5.979590in}{3.440004in}}%
\pgfpathlineto{\pgfqpoint{6.002035in}{3.314702in}}%
\pgfpathlineto{\pgfqpoint{6.024481in}{3.268352in}}%
\pgfpathlineto{\pgfqpoint{6.046926in}{3.311182in}}%
\pgfpathlineto{\pgfqpoint{6.069372in}{3.704643in}}%
\pgfpathlineto{\pgfqpoint{6.091817in}{3.737003in}}%
\pgfpathlineto{\pgfqpoint{6.114263in}{3.665519in}}%
\pgfpathlineto{\pgfqpoint{6.136709in}{3.435217in}}%
\pgfpathlineto{\pgfqpoint{6.159154in}{3.610381in}}%
\pgfpathlineto{\pgfqpoint{6.159154in}{3.610381in}}%
\pgfusepath{stroke}%
\end{pgfscope}%
\begin{pgfscope}%
\pgfpathrectangle{\pgfqpoint{0.413090in}{0.413680in}}{\pgfqpoint{5.746064in}{3.574842in}}%
\pgfusepath{clip}%
\pgfsetbuttcap%
\pgfsetroundjoin%
\pgfsetlinewidth{1.003750pt}%
\definecolor{currentstroke}{rgb}{0.690196,0.690196,0.690196}%
\pgfsetstrokecolor{currentstroke}%
\pgfsetdash{{3.700000pt}{1.600000pt}}{0.000000pt}%
\pgfpathmoveto{\pgfqpoint{0.413090in}{0.611055in}}%
\pgfpathlineto{\pgfqpoint{0.502872in}{0.624237in}}%
\pgfpathlineto{\pgfqpoint{0.592655in}{0.590376in}}%
\pgfpathlineto{\pgfqpoint{0.682437in}{0.592217in}}%
\pgfpathlineto{\pgfqpoint{0.772219in}{0.598719in}}%
\pgfpathlineto{\pgfqpoint{0.862001in}{0.642923in}}%
\pgfpathlineto{\pgfqpoint{0.951784in}{0.608758in}}%
\pgfpathlineto{\pgfqpoint{1.041566in}{0.690261in}}%
\pgfpathlineto{\pgfqpoint{1.131348in}{0.704370in}}%
\pgfpathlineto{\pgfqpoint{1.221130in}{0.650663in}}%
\pgfpathlineto{\pgfqpoint{1.310913in}{0.751647in}}%
\pgfpathlineto{\pgfqpoint{1.400695in}{0.707080in}}%
\pgfpathlineto{\pgfqpoint{1.490477in}{0.761866in}}%
\pgfpathlineto{\pgfqpoint{1.580259in}{0.760825in}}%
\pgfpathlineto{\pgfqpoint{1.670042in}{0.756495in}}%
\pgfpathlineto{\pgfqpoint{1.759824in}{0.765129in}}%
\pgfpathlineto{\pgfqpoint{1.849606in}{0.847617in}}%
\pgfpathlineto{\pgfqpoint{1.939388in}{0.859619in}}%
\pgfpathlineto{\pgfqpoint{2.029171in}{1.018025in}}%
\pgfpathlineto{\pgfqpoint{2.118953in}{1.141670in}}%
\pgfpathlineto{\pgfqpoint{2.208735in}{1.168186in}}%
\pgfpathlineto{\pgfqpoint{2.298517in}{0.968438in}}%
\pgfpathlineto{\pgfqpoint{2.388300in}{0.936149in}}%
\pgfpathlineto{\pgfqpoint{2.478082in}{0.991428in}}%
\pgfpathlineto{\pgfqpoint{2.567864in}{1.069559in}}%
\pgfpathlineto{\pgfqpoint{2.657646in}{1.024768in}}%
\pgfpathlineto{\pgfqpoint{2.747429in}{1.135663in}}%
\pgfpathlineto{\pgfqpoint{2.837211in}{1.215185in}}%
\pgfpathlineto{\pgfqpoint{2.926993in}{1.382240in}}%
\pgfpathlineto{\pgfqpoint{3.016775in}{1.239636in}}%
\pgfpathlineto{\pgfqpoint{3.106558in}{1.345708in}}%
\pgfpathlineto{\pgfqpoint{3.196340in}{1.408865in}}%
\pgfpathlineto{\pgfqpoint{3.286122in}{1.662816in}}%
\pgfpathlineto{\pgfqpoint{3.375904in}{1.531451in}}%
\pgfpathlineto{\pgfqpoint{3.465687in}{1.479887in}}%
\pgfpathlineto{\pgfqpoint{3.555469in}{1.460242in}}%
\pgfpathlineto{\pgfqpoint{3.645251in}{1.278061in}}%
\pgfpathlineto{\pgfqpoint{3.735033in}{1.309644in}}%
\pgfpathlineto{\pgfqpoint{3.824816in}{1.568419in}}%
\pgfpathlineto{\pgfqpoint{3.914598in}{1.720165in}}%
\pgfpathlineto{\pgfqpoint{4.004380in}{1.795109in}}%
\pgfpathlineto{\pgfqpoint{4.094162in}{1.677930in}}%
\pgfpathlineto{\pgfqpoint{4.183945in}{1.522100in}}%
\pgfpathlineto{\pgfqpoint{4.273727in}{1.474144in}}%
\pgfpathlineto{\pgfqpoint{4.363509in}{1.932741in}}%
\pgfpathlineto{\pgfqpoint{4.453291in}{2.007276in}}%
\pgfpathlineto{\pgfqpoint{4.543074in}{2.077562in}}%
\pgfpathlineto{\pgfqpoint{4.632856in}{1.837803in}}%
\pgfpathlineto{\pgfqpoint{4.722638in}{1.763893in}}%
\pgfpathlineto{\pgfqpoint{4.812420in}{1.852777in}}%
\pgfpathlineto{\pgfqpoint{4.902203in}{1.807858in}}%
\pgfpathlineto{\pgfqpoint{4.991985in}{2.097023in}}%
\pgfpathlineto{\pgfqpoint{5.081767in}{2.016363in}}%
\pgfpathlineto{\pgfqpoint{5.171549in}{2.203636in}}%
\pgfpathlineto{\pgfqpoint{5.261332in}{2.123969in}}%
\pgfpathlineto{\pgfqpoint{5.351114in}{2.160145in}}%
\pgfpathlineto{\pgfqpoint{5.440896in}{2.064267in}}%
\pgfpathlineto{\pgfqpoint{5.530678in}{2.499390in}}%
\pgfpathlineto{\pgfqpoint{5.620461in}{2.774954in}}%
\pgfpathlineto{\pgfqpoint{5.710243in}{3.209748in}}%
\pgfpathlineto{\pgfqpoint{5.800025in}{3.190052in}}%
\pgfpathlineto{\pgfqpoint{5.889807in}{3.421759in}}%
\pgfpathlineto{\pgfqpoint{5.979590in}{3.376302in}}%
\pgfpathlineto{\pgfqpoint{6.069372in}{3.632263in}}%
\pgfpathlineto{\pgfqpoint{6.159154in}{3.542088in}}%
\pgfusepath{stroke}%
\end{pgfscope}%
\begin{pgfscope}%
\pgfsetrectcap%
\pgfsetmiterjoin%
\pgfsetlinewidth{0.501875pt}%
\definecolor{currentstroke}{rgb}{0.000000,0.000000,0.000000}%
\pgfsetstrokecolor{currentstroke}%
\pgfsetdash{}{0pt}%
\pgfpathmoveto{\pgfqpoint{0.413090in}{0.413680in}}%
\pgfpathlineto{\pgfqpoint{0.413090in}{3.988523in}}%
\pgfusepath{stroke}%
\end{pgfscope}%
\begin{pgfscope}%
\pgfsetrectcap%
\pgfsetmiterjoin%
\pgfsetlinewidth{0.501875pt}%
\definecolor{currentstroke}{rgb}{0.000000,0.000000,0.000000}%
\pgfsetstrokecolor{currentstroke}%
\pgfsetdash{}{0pt}%
\pgfpathmoveto{\pgfqpoint{6.159154in}{0.413680in}}%
\pgfpathlineto{\pgfqpoint{6.159154in}{3.988523in}}%
\pgfusepath{stroke}%
\end{pgfscope}%
\begin{pgfscope}%
\pgfsetrectcap%
\pgfsetmiterjoin%
\pgfsetlinewidth{0.501875pt}%
\definecolor{currentstroke}{rgb}{0.000000,0.000000,0.000000}%
\pgfsetstrokecolor{currentstroke}%
\pgfsetdash{}{0pt}%
\pgfpathmoveto{\pgfqpoint{0.413090in}{0.413680in}}%
\pgfpathlineto{\pgfqpoint{6.159154in}{0.413680in}}%
\pgfusepath{stroke}%
\end{pgfscope}%
\begin{pgfscope}%
\pgfsetrectcap%
\pgfsetmiterjoin%
\pgfsetlinewidth{0.501875pt}%
\definecolor{currentstroke}{rgb}{0.000000,0.000000,0.000000}%
\pgfsetstrokecolor{currentstroke}%
\pgfsetdash{}{0pt}%
\pgfpathmoveto{\pgfqpoint{0.413090in}{3.988523in}}%
\pgfpathlineto{\pgfqpoint{6.159154in}{3.988523in}}%
\pgfusepath{stroke}%
\end{pgfscope}%
\begin{pgfscope}%
\pgfsetbuttcap%
\pgfsetmiterjoin%
\definecolor{currentfill}{rgb}{1.000000,1.000000,1.000000}%
\pgfsetfillcolor{currentfill}%
\pgfsetfillopacity{0.800000}%
\pgfsetlinewidth{1.003750pt}%
\definecolor{currentstroke}{rgb}{0.800000,0.800000,0.800000}%
\pgfsetstrokecolor{currentstroke}%
\pgfsetstrokeopacity{0.800000}%
\pgfsetdash{}{0pt}%
\pgfpathmoveto{\pgfqpoint{0.510312in}{3.279808in}}%
\pgfpathlineto{\pgfqpoint{2.427181in}{3.279808in}}%
\pgfpathquadraticcurveto{\pgfqpoint{2.454959in}{3.279808in}}{\pgfqpoint{2.454959in}{3.307586in}}%
\pgfpathlineto{\pgfqpoint{2.454959in}{3.891300in}}%
\pgfpathquadraticcurveto{\pgfqpoint{2.454959in}{3.919078in}}{\pgfqpoint{2.427181in}{3.919078in}}%
\pgfpathlineto{\pgfqpoint{0.510312in}{3.919078in}}%
\pgfpathquadraticcurveto{\pgfqpoint{0.482534in}{3.919078in}}{\pgfqpoint{0.482534in}{3.891300in}}%
\pgfpathlineto{\pgfqpoint{0.482534in}{3.307586in}}%
\pgfpathquadraticcurveto{\pgfqpoint{0.482534in}{3.279808in}}{\pgfqpoint{0.510312in}{3.279808in}}%
\pgfpathlineto{\pgfqpoint{0.510312in}{3.279808in}}%
\pgfpathclose%
\pgfusepath{stroke,fill}%
\end{pgfscope}%
\begin{pgfscope}%
\pgfsetrectcap%
\pgfsetroundjoin%
\pgfsetlinewidth{1.003750pt}%
\definecolor{currentstroke}{rgb}{0.000000,0.000000,0.000000}%
\pgfsetstrokecolor{currentstroke}%
\pgfsetdash{}{0pt}%
\pgfpathmoveto{\pgfqpoint{0.538090in}{3.814738in}}%
\pgfpathlineto{\pgfqpoint{0.676979in}{3.814738in}}%
\pgfpathlineto{\pgfqpoint{0.815868in}{3.814738in}}%
\pgfusepath{stroke}%
\end{pgfscope}%
\begin{pgfscope}%
\definecolor{textcolor}{rgb}{0.000000,0.000000,0.000000}%
\pgfsetstrokecolor{textcolor}%
\pgfsetfillcolor{textcolor}%
\pgftext[x=0.926979in,y=3.766127in,left,base]{\color{textcolor}{\rmfamily\fontsize{10.000000}{12.000000}\selectfont\catcode`\^=\active\def^{\ifmmode\sp\else\^{}\fi}\catcode`\%=\active\def%{\%}Exact ($Y_t$)}}%
\end{pgfscope}%
\begin{pgfscope}%
\pgfsetbuttcap%
\pgfsetroundjoin%
\pgfsetlinewidth{1.003750pt}%
\definecolor{currentstroke}{rgb}{0.690196,0.188235,0.376471}%
\pgfsetstrokecolor{currentstroke}%
\pgfsetdash{{3.700000pt}{1.600000pt}}{0.000000pt}%
\pgfpathmoveto{\pgfqpoint{0.538090in}{3.615537in}}%
\pgfpathlineto{\pgfqpoint{0.676979in}{3.615537in}}%
\pgfpathlineto{\pgfqpoint{0.815868in}{3.615537in}}%
\pgfusepath{stroke}%
\end{pgfscope}%
\begin{pgfscope}%
\definecolor{textcolor}{rgb}{0.000000,0.000000,0.000000}%
\pgfsetstrokecolor{textcolor}%
\pgfsetfillcolor{textcolor}%
\pgftext[x=0.926979in,y=3.566926in,left,base]{\color{textcolor}{\rmfamily\fontsize{10.000000}{12.000000}\selectfont\catcode`\^=\active\def^{\ifmmode\sp\else\^{}\fi}\catcode`\%=\active\def%{\%}Milstein ($X_t$): Fine Grid}}%
\end{pgfscope}%
\begin{pgfscope}%
\pgfsetbuttcap%
\pgfsetroundjoin%
\pgfsetlinewidth{1.003750pt}%
\definecolor{currentstroke}{rgb}{0.690196,0.690196,0.690196}%
\pgfsetstrokecolor{currentstroke}%
\pgfsetdash{{3.700000pt}{1.600000pt}}{0.000000pt}%
\pgfpathmoveto{\pgfqpoint{0.538090in}{3.416335in}}%
\pgfpathlineto{\pgfqpoint{0.676979in}{3.416335in}}%
\pgfpathlineto{\pgfqpoint{0.815868in}{3.416335in}}%
\pgfusepath{stroke}%
\end{pgfscope}%
\begin{pgfscope}%
\definecolor{textcolor}{rgb}{0.000000,0.000000,0.000000}%
\pgfsetstrokecolor{textcolor}%
\pgfsetfillcolor{textcolor}%
\pgftext[x=0.926979in,y=3.367724in,left,base]{\color{textcolor}{\rmfamily\fontsize{10.000000}{12.000000}\selectfont\catcode`\^=\active\def^{\ifmmode\sp\else\^{}\fi}\catcode`\%=\active\def%{\%}Milstein ($X_t$): Coarse Grid}}%
\end{pgfscope}%
\end{pgfpicture}%
\makeatother%
\endgroup%

    \caption{Comparison of Milstein approximations against the exact analytical solution for a GBM. The simulation uses a fine grid step size of $\Delta t = 2^{-8}$ and a coarse grid step size of $\Delta t_{\text{coarse}} = 4\Delta t$. Parameters: $\mu=1.0$, $\sigma=0.5$, $X_0=1.0$.}
    \label{fig:milsteinCoarseFineComparison}
\end{figure}